% Limitation des dysfonctionnements des structures responsables du contrôle de la motricité — SVT Tle D — Cours signé OnBuch+
% Programme officiel MINESEC. Compiler avec : tectonic N10-motricite-volontaire-aires-cerebrales.tex
\newcommand{\DOCMATIERE}{SVT}
\newcommand{\DOCNIVEAU}{Tle D}
\newcommand{\DOCMODULE}{Module I : Le Monde Vivant — Activités cérébrales et motricité volontaire}
\newcommand{\DOCLECON}{N10}
\newcommand{\DOCTITRE}{Limitation des dysfonctionnements des structures responsables du contrôle de la motricité}
% OnBuch+ — préambule « cours signés » ENRICHI (compilé avec tectonic / XeLaTeX)
% Système visuel « Pop » d'OnBuch+ : fond crème (jamais blanc), bordures encre
% épaisses, ombres « dures » décalées, titres Archivo Black en capitales.
% Version MATHÉMATIQUES : graphes (pgfplots/tikz), boîtes pédagogiques
% (définition, propriété, méthode, attention, le savais-tu, exercice résolu,
% exercice, TP, bilan), page de garde et signature OnBuch+ ancrée en bas.
%
% Macros attendues dans main.tex AVANT \input{preamble} :
%   \DOCMATIERE  \DOCNIVEAU  \DOCMODULE  \DOCLECON (numéro)  \DOCTITRE
\documentclass[11pt]{article}

\usepackage{fontspec}
\usepackage[french]{babel}
\usepackage[a4paper,margin=2cm,top=3.4cm,bottom=2.8cm,headheight=1.4cm,footskip=1.3cm]{geometry}
\usepackage{amsmath,amssymb,amsfonts}
\usepackage{mathtools} % \xmapsto, \coloneqq... souvent utilisés par les modèles
\usepackage{cancel} % simplifications barrées (bilans de forces)
% \abs{z} (module, valeur absolue) et \norm{u} (norme), utilisés par les modèles
\providecommand{\abs}{}\let\abs\relax\DeclarePairedDelimiter{\abs}{\lvert}{\rvert}
\providecommand{\norm}{}\let\norm\relax\DeclarePairedDelimiter{\norm}{\lVert}{\rVert}
\usepackage[table]{xcolor} % [table] : fournit aussi \rowcolors (alternance de lignes)
\usepackage{pagecolor}
\usepackage{tikz}
\usetikzlibrary{arrows.meta,positioning,calc,shapes.geometric,shapes.misc,shapes.arrows,decorations.pathmorphing,decorations.pathreplacing,decorations.markings,patterns,shadows,mindmap,trees,fit,backgrounds,matrix,intersections,angles,quotes}
\usepackage{pgfplots}
\usepackage[version=4]{mhchem} % équations nucléaires \ce{^{235}_{92}U}
\usepackage[european,siunitx]{circuitikz} % schémas électriques
\pgfplotsset{compat=1.18}
\usepgfplotslibrary{fillbetween,groupplots} % aires entre courbes (calcul intégral)
\usepackage{enumitem}
\usepackage{fancyhdr}
% [most] charge toutes les bibliothèques tcolorbox (listings, minted, raster,
% documentation...) et gonfle inutilement la mémoire TeX sur un document long
% (~300 boîtes) : on ne charge que ce qui est réellement utilisé.
\usepackage[skins,breakable]{tcolorbox}
\usepackage{graphicx}
\usepackage{colortbl}
\usepackage{booktabs}
\usepackage{tabularx}
\usepackage{diagbox} % case d'en-tête diagonale des tableaux à double entrée
% Colonnes extensibles centrée / gauche / droite (C, L, R), souvent utilisées par les modèles
\newcolumntype{C}{>{\centering\arraybackslash}X}
\newcolumntype{L}{>{\raggedright\arraybackslash}X}
\newcolumntype{R}{>{\raggedleft\arraybackslash}X}
\usepackage{array}
\usepackage{multicol}
\usepackage{siunitx}
% parse-numbers=false : accepte aussi \SI{2,5 \times 10^{-3}}{...}
\sisetup{locale=FR,output-decimal-marker={,},group-minimum-digits=4,parse-numbers=false}
\DeclareSIUnit\ohm{\ensuremath{\symup{\Omega}}} % U+2126 absent de la police
\DeclareSIUnit\division{div} % réglages d'oscilloscope (ms/div, V/div)
\DeclareSIUnit\tr{tr} % tours (vitesse de rotation, ex. \SI{1500}{\tr\per\minute})
\DeclareSIUnit\molar{M} % concentration molaire (ex. \SI{3}{\molar}, \SI{1}{\micro\molar})
\DeclareSIUnit\an{an} % année (le modèle confond parfois avec \year, un compteur TeX) — ex. \SI{5}{\kilo\meter\per\an}
\DeclareSIUnit\FCFA{FCFA} % franc CFA (ex. \SI{500}{\FCFA\per\kilo\gram})
\DeclareSIUnit\degreeBrix{\textdegree Brix} % teneur en sucre d'un sirop/jus (réfractométrie)
\DeclareSIUnit\month{mois} % le modèle utilise parfois \month, un compteur TeX, comme unité de durée
\usepackage{qrcode}
% \degree : commande fréquemment utilisée par les modèles pour un angle ou une
% température en dehors de \SI{}{} (non fournie par les paquets chargés).
\providecommand{\degree}{^{\circ}}
% \qed (fin de démonstration), utilisé par les modèles sans amsthm
\providecommand{\qed}{\hfill$\square$}
% \barre : double barre des valeurs interdites dans un tableau de variations (tkz-tab)
\providecommand{\barre}{\ensuremath{\|}}
% environnement proof (amsthm non chargé) utilisé par les modèles
\ifdefined\proof\else\newenvironment{proof}[1][Démonstration]{\par\noindent\textit{#1.}\ }{\qed\par}\fi
\usepackage{lastpage}
\usepackage{hyperref}
\hypersetup{hidelinks}

% ============================================================
% Palette « Pop » OnBuch+ — mêmes hex que la classe P (plus_ui.dart)
% ============================================================
\definecolor{poporange}{HTML}{F59321}
\definecolor{popdark}{HTML}{E07A0C}
\definecolor{popcream}{HTML}{FFF6E8}
\definecolor{popink}{HTML}{1C1714}
\definecolor{poppurple}{HTML}{7A5AE0}
\definecolor{popgreen}{HTML}{2E9E6A}
\definecolor{poppink}{HTML}{E0457B}
\definecolor{popblue}{HTML}{2F7DE1}
\definecolor{popgold}{HTML}{C98A12}
\definecolor{popmuted}{HTML}{8A7F76}
\definecolor{popline}{HTML}{EADFD2}
\definecolor{popgray}{HTML}{8A7F76}
\colorlet{popgrey}{popgray}
% Alias tolérants (le modèle nomme parfois les couleurs différemment)
\colorlet{popwhite}{white}
\colorlet{popblack}{popink}
\colorlet{popred}{poppink}
\colorlet{popyellow}{popgold}
% Teintes claires (fonds de tableaux/encadrés) — le modèle nomme parfois
% les couleurs avec un suffixe L (ex. popblueL) sans les avoir définies.
\colorlet{poporangeL}{poporange!20}
\colorlet{popdarkL}{popdark!20}
\colorlet{poppurpleL}{poppurple!20}
\colorlet{popgreenL}{popgreen!20}
\colorlet{poppinkL}{poppink!20}
\colorlet{popblueL}{popblue!20}
\colorlet{popgoldL}{popgold!20}
\colorlet{popredL}{popred!20}
\colorlet{popgrayL}{popgray!20}
% Teintes claires pour fonds de tableaux / zones
\colorlet{popblueL}{popblue!10!white}
\colorlet{poporangeL}{poporange!14!white}
\colorlet{popgreenL}{popgreen!12!white}
\colorlet{poppurpleL}{poppurple!12!white}
\colorlet{poppinkL}{poppink!12!white}
\colorlet{popgoldL}{popgold!14!white}

\pagecolor{popcream}

% ============================================================
% Typographie
% ============================================================
\defaultfontfeatures{Ligatures=TeX}
\setmainfont{PlusJakartaSans-Medium.ttf}[Path=../../../fonts/,BoldFont=PlusJakartaSans-Bold.ttf]
% Maths en sans-serif (Fira Math) avec chiffres et lettres droites de Plus
% Jakarta, pour que les formules se fondent dans le texte.
\usepackage[math-style=ISO,bold-style=ISO]{unicode-math}
\setmathfont{FiraMath-Regular.otf}[Path=../../../fonts/]
\setmathfont{PlusJakartaSans-Medium.ttf}[Path=../../../fonts/,range={up/num,up/{Latin,latin},\mathup}]
\usepackage{icomma} % virgule décimale sans espace en mode math
% Caractères mathématiques Unicode absents des polices de texte (Plus Jakarta,
% Archivo Black) : rendus par leur équivalent mathématique, en texte comme en
% formule — sinon ils s'affichent en carré vide (ex. « Arithmétique dans ℤ »).
\usepackage{newunicodechar}
\newunicodechar{ℝ}{\ensuremath{\mathbb{R}}}
\newunicodechar{ℤ}{\ensuremath{\mathbb{Z}}}
\newunicodechar{ℂ}{\ensuremath{\mathbb{C}}}
\newunicodechar{ℕ}{\ensuremath{\mathbb{N}}}
\newunicodechar{ℚ}{\ensuremath{\mathbb{Q}}}
\newunicodechar{𝔻}{\ensuremath{\mathbb{D}}}
\newunicodechar{α}{\ensuremath{\alpha}}
\newunicodechar{β}{\ensuremath{\beta}}
\newunicodechar{γ}{\ensuremath{\gamma}}
\newunicodechar{δ}{\ensuremath{\delta}}
\newunicodechar{ε}{\ensuremath{\varepsilon}}
\newunicodechar{θ}{\ensuremath{\theta}}
\newunicodechar{λ}{\ensuremath{\lambda}}
\newunicodechar{μ}{\ensuremath{\mu}}
\newunicodechar{σ}{\ensuremath{\sigma}}
\newunicodechar{τ}{\ensuremath{\tau}}
\newunicodechar{φ}{\ensuremath{\varphi}}
\newunicodechar{ψ}{\ensuremath{\psi}}
\newunicodechar{ω}{\ensuremath{\omega}}
\newunicodechar{Σ}{\ensuremath{\Sigma}}
% majuscules produites par \MakeUppercase dans les titres (α-aminés -> Α)
\newunicodechar{Α}{\ensuremath{\alpha}}
\newunicodechar{Β}{\ensuremath{\beta}}
\newunicodechar{Γ}{\ensuremath{\Gamma}}
\newunicodechar{Δ}{\ensuremath{\Delta}}
\newunicodechar{Ε}{\ensuremath{\varepsilon}}
\newunicodechar{Θ}{\ensuremath{\Theta}}
\newunicodechar{Λ}{\ensuremath{\Lambda}}
\newunicodechar{Μ}{\ensuremath{\mu}}
\newunicodechar{Τ}{\ensuremath{\tau}}
\newunicodechar{Φ}{\ensuremath{\Phi}}
\newunicodechar{Ψ}{\ensuremath{\Psi}}
\newunicodechar{Ω}{\ensuremath{\Omega}}
\newunicodechar{↦}{\ensuremath{\mapsto}}
\newunicodechar{∈}{\ensuremath{\in}}
\newunicodechar{∉}{\ensuremath{\notin}}
\newunicodechar{∧}{\ensuremath{\wedge}}
\newunicodechar{⇔}{\ensuremath{\Leftrightarrow}}
\newunicodechar{⟺}{\ensuremath{\Longleftrightarrow}}
\newunicodechar{⇒}{\ensuremath{\Rightarrow}}
\newunicodechar{⟹}{\ensuremath{\Longrightarrow}}
\newunicodechar{∘}{\ensuremath{\circ}}
\newunicodechar{∪}{\ensuremath{\cup}}
\newunicodechar{∩}{\ensuremath{\cap}}
\newunicodechar{≡}{\ensuremath{\equiv}}
\newunicodechar{⊂}{\ensuremath{\subset}}
\newunicodechar{⊥}{\ensuremath{\perp}}
\newunicodechar{∃}{\ensuremath{\exists}}
\newunicodechar{∀}{\ensuremath{\forall}}
\newunicodechar{∛}{\ensuremath{\sqrt[3]{\,}}}
\newunicodechar{∜}{\ensuremath{\sqrt[4]{\,}}}
\newunicodechar{⃗}{\ensuremath{{}^{\to}}}
\newunicodechar{ⁿ}{\ifmmode^{n}\else\textsuperscript{\ensuremath{n}}\fi}
\newunicodechar{ˣ}{\ifmmode^{x}\else\textsuperscript{\ensuremath{x}}\fi}
\newunicodechar{ᵇ}{\ifmmode^{b}\else\textsuperscript{\ensuremath{b}}\fi}
\newunicodechar{ʳ}{\ifmmode^{r}\else\textsuperscript{\ensuremath{r}}\fi}
\newunicodechar{ᵘ}{\ifmmode^{u}\else\textsuperscript{\ensuremath{u}}\fi}
\newunicodechar{ʸ}{\ifmmode^{y}\else\textsuperscript{\ensuremath{y}}\fi}
\newunicodechar{ᵃ}{\ifmmode^{a}\else\textsuperscript{\ensuremath{a}}\fi}
\newunicodechar{ᵉ}{\ifmmode^{e}\else\textsuperscript{\ensuremath{e}}\fi}
\newunicodechar{ᵏ}{\ifmmode^{k}\else\textsuperscript{\ensuremath{k}}\fi}
\newunicodechar{ᵖ}{\ifmmode^{p}\else\textsuperscript{\ensuremath{p}}\fi}
\newunicodechar{ᴺ}{\ifmmode^{N}\else\textsuperscript{\ensuremath{N}}\fi}
\newunicodechar{ᶜ}{\ifmmode^{c}\else\textsuperscript{\ensuremath{c}}\fi}
\newunicodechar{ʰ}{\ifmmode^{h}\else\textsuperscript{\ensuremath{h}}\fi}
\newunicodechar{ᵠ}{\ifmmode^{q}\else\textsuperscript{\ensuremath{q}}\fi}
\newunicodechar{ₙ}{\ifmmode_{n}\else\textsubscript{\ensuremath{n}}\fi}
\newunicodechar{ₐ}{\ifmmode_{a}\else\textsubscript{\ensuremath{a}}\fi}
\newunicodechar{ᵢ}{\ifmmode_{i}\else\textsubscript{\ensuremath{i}}\fi}
\newunicodechar{ᵣ}{\ifmmode_{r}\else\textsubscript{\ensuremath{r}}\fi}
\newunicodechar{ₖ}{\ifmmode_{k}\else\textsubscript{\ensuremath{k}}\fi}
\newunicodechar{ᵦ}{\ifmmode_{\beta}\else\textsubscript{\ensuremath{\beta}}\fi}
\newunicodechar{ₓ}{\ifmmode_{x}\else\textsubscript{\ensuremath{x}}\fi}
\newfontfamily\popdisplay{ArchivoBlack-Regular.ttf}[Path=../../../fonts/]
\newfontfamily\poplogo{SpaceGrotesk-Bold.ttf}[Path=../../../fonts/]
\newfontfamily\popsemi{PlusJakartaSans-SemiBold.ttf}[Path=../../../fonts/]
\newfontfamily\popxbold{PlusJakartaSans-ExtraBold.ttf}[Path=../../../fonts/]
\newfontfamily\popxboldls{PlusJakartaSans-ExtraBold.ttf}[Path=../../../fonts/,LetterSpace=3.0]

\setlist[enumerate]{leftmargin=1.7em,itemsep=5pt,topsep=4pt}
\setlist[itemize]{leftmargin=1.7em,itemsep=4pt,topsep=4pt,label={\color{poporange}\textbullet}}
\setlength{\parindent}{0pt}
\setlength{\parskip}{5pt}
\renewcommand{\arraystretch}{1.25}

% pgfplots : style Pop par défaut
\pgfplotsset{
  every axis/.append style={axis line style={popink,line width=1pt},tick style={popink},
    grid style={popline,line width=0.5pt},label style={font=\small},tick label style={font=\footnotesize}},
  popcurve/.style={poporange,line width=1.8pt,no markers,smooth},
}
% Même style disponible dans un \draw TikZ simple (hors pgfplots)
\tikzset{popcurve/.style={poporange,line width=1.8pt}}
\tikzset{popgrid/.style={popline,very thin}}
\colorlet{popcurve}{poporange} % le nom du style est parfois utilisé comme couleur

% ============================================================
% Pill / squiggle
% ============================================================
\newcommand{\poppill}[3]{%
  \tikz[baseline=-0.55ex]\node[draw=popink,line width=1.2pt,rounded corners=2.6pt,%
    fill=#2,inner xsep=6.5pt,inner ysep=3pt]{%
    \popxboldls\fontsize{7}{7}\selectfont\color{#3}\MakeUppercase{#1}};%
}
\newcommand{\popsquiggle}[1]{%
  \tikz[baseline=-0.4ex]{
    \draw[#1,line width=2.6pt,line cap=round]
      plot [smooth] coordinates {(0,0.09) (0.34,0.01) (0.68,0.21) (1.02,0.01) (1.36,0.10)};
  }%
}
% Mot-clé surligné (marqueur orange sous le texte)
\newcommand{\cle}[1]{{\popxbold\color{popdark}#1}}

% ============================================================
% En-tête / pied
% ============================================================
\pagestyle{fancy}
\fancyhf{}
\renewcommand{\headrulewidth}{0pt}
\renewcommand{\footrulewidth}{0pt}
\lhead{\raisebox{-0.15ex}{\poplogo\fontsize{13}{13}\selectfont\color{popink}OnBuch\color{poporange}+}}
\rhead{\poppill{\DOCMATIERE\ \textbullet\ \DOCNIVEAU\ \textbullet\ Leçon \DOCLECON}{white}{popink}}
\lfoot{\popxbold\fontsize{7.6}{7.6}\selectfont\color{popmuted}\MakeUppercase{Cours signé OnBuch+}\ \textbullet\ {\popsemi\DOCTITRE}}
\rfoot{\tikz[baseline=-0.6ex]\node[rounded corners=8.5pt,fill=popink,minimum height=17pt,minimum width=17pt,inner xsep=5pt,inner sep=0]{\popxbold\fontsize{8}{8}\selectfont\color{popcream}\thepage\,/\,\pageref*{LastPage}};}

% ============================================================
% Titres
% ============================================================
\newcommand{\chaptitre}[2]{%
  \begin{tcolorbox}[enhanced,colback=poporange,colframe=popink,boxrule=2.6pt,arc=11pt,
      drop shadow=popink,left=15pt,right=15pt,top=12pt,bottom=11pt,before skip=3pt,after skip=18pt]
    \poppill{#2}{popink}{popcream}\par\vspace{9pt}
    {\raggedright\hyphenpenalty=10000\exhyphenpenalty=10000\popdisplay\fontsize{22}{24}\selectfont\color{popcream}\MakeUppercase{#1}\par}
  \end{tcolorbox}
}
% Section numérotée : pastille orange avec le numéro + titre Archivo
\newcounter{popsec}
\newcommand{\coursec}[1]{%
  \stepcounter{popsec}\setcounter{popsub}{0}%
  \par\vspace{16pt}\noindent
  \begin{minipage}{\linewidth}
  \tikz[baseline=-0.7ex]\node[draw=popink,line width=1.6pt,fill=poporange,rounded corners=4pt,minimum size=22pt,inner sep=0]{\popdisplay\fontsize{12}{12}\selectfont\color{popcream}\thepopsec};%
  \hspace{8pt}{\popdisplay\fontsize{15}{17}\selectfont\color{popink}\MakeUppercase{#1}}\par
  \vspace{4pt}\noindent{\color{poporange}\rule{36pt}{3.6pt}}
  \end{minipage}\par\nopagebreak\vspace{8pt}
}
\newcounter{popsub}
\newcommand{\courssub}[1]{%
  \stepcounter{popsub}%
  \par\vspace{9pt}\noindent{\popxbold\fontsize{11.5}{13}\selectfont\color{popink}{\color{poporange}\thepopsec.\thepopsub}\hspace{6pt}#1}\par\nopagebreak\vspace{4pt}
}

% ============================================================
% Boîtes pédagogiques — même recette : fond blanc, bordure épaisse colorée,
% ombre dure encre, badge plein. Titre optionnel [#1].
% ============================================================
\tcbset{popbase/.style={breakable,enhanced,colback=white,boxrule=2.3pt,arc=9pt,
  shadow={3pt}{-3pt}{0pt}{popink},left=14pt,right=14pt,top=11pt,bottom=11pt,before skip=11pt,after skip=15pt,
  fontupper=\popsemi\fontsize{10.6}{14.5}\selectfont\color{popink},
  fontlower=\popsemi\fontsize{10.6}{14.5}\selectfont\color{popink}}}
% Coche dessinée (le glyphe ✓ n'existe pas dans les polices du document)
\newcommand{\popcheck}[1]{\tikz[baseline=-0.5ex]\draw[#1,line width=1.6pt,line cap=round,line join=round](-0.13,0.0)--(-0.04,-0.09)--(0.13,0.1);}
\providecommand{\coche}{\popcheck{popgreen}}
\providecommand{\resultat}[1]{\par\smallskip\noindent\fcolorbox{popgreen}{popgreenL}{\parbox{\dimexpr\linewidth-2\fboxsep-2\fboxrule\relax}{\textbf{Résultat.} #1}}\par\smallskip}
\newcommand{\popboxhead}[3]{\poppill{#1}{#2}{white}\ifx\relax#3\relax\else\hspace{6pt}{\popxbold\fontsize{10.6}{12}\selectfont\color{popink}#3}\fi\par\vspace{7pt}}

\newtcolorbox{aretenir}[1][]{popbase,colframe=popgreen,before upper={\popboxhead{À retenir}{popgreen}{#1}}}
\newtcolorbox{exemplebox}[1][]{popbase,colframe=poporange,before upper={\popboxhead{Exemple}{poporange}{#1}}}
\newtcolorbox{definition}[1][]{popbase,colframe=popblue,before upper={\popboxhead{Définition}{popblue}{#1}}}
\newtcolorbox{propriete}[1][]{popbase,colframe=poppurple,before upper={\popboxhead{Propriété}{poppurple}{#1}}}
\newtcolorbox{methode}[1][]{popbase,colframe=popgold,before upper={\popboxhead{Méthode}{popgold}{#1}}}
\newtcolorbox{attention}[1][]{popbase,colframe=poppink,before upper={\popboxhead{Attention}{poppink}{#1}}}
\newtcolorbox{experience}[1][]{popbase,colframe=popblue,colback=popblueL,before upper={\popboxhead{Expérience}{popblue}{#1}}}
% « Le savais-tu ? » : carte violette pleine (contexte, culture, Cameroun)
\newtcolorbox{savaistu}[1][]{popbase,colback=poppurple,colframe=popink,
  fontupper=\popsemi\fontsize{10.4}{14.2}\selectfont\color{white},
  before upper={\poppill{Le savais-tu\,?}{white}{poppurple}\ifx\relax#1\relax\else\hspace{6pt}{\popxbold\color{white}#1}\fi\par\vspace{7pt}}}
% Exercice résolu : énoncé en haut, \tcblower puis la solution sur fond crème
\newcounter{popexo}\newcounter{popexr}
\newtcolorbox{exoresolu}[1][]{popbase,colframe=poporange,colbacklower=popcream,
  segmentation style={popink,line width=1.2pt,dashed},
  before upper={\stepcounter{popexr}\popboxhead{Exercice résolu \thepopexr}{poporange}{#1}},
  before lower={\poppill{Solution}{popink}{popcream}\par\vspace{6pt}}}
% Exercice (non résolu en place) — la correction est dans la partie Corrigés
\newtcolorbox{exercice}[1][]{popbase,colframe=popink,
  before upper={\stepcounter{popexo}\popboxhead{Exercice \thepopexo}{popink}{#1}}}
% Corrigé
\newtcolorbox{corrige}[1][]{popbase,colframe=popgreen,colback=popgreenL,
  before upper={\popboxhead{Corrigé}{popgreen}{#1}}}
% Objectifs / prérequis / bilan
\newtcolorbox{objectifs}{popbase,colframe=popink,colback=poporangeL,
  before upper={\popboxhead{Objectifs de la leçon}{popink}{}}}
\newtcolorbox{prerequis}{popbase,colframe=popink,colback=popblueL,
  before upper={\popboxhead{Prérequis}{popblue}{}}}
\newtcolorbox{bilan}{popbase,colframe=popink,colback=white,boxrule=2.8pt,
  before upper={\popboxhead{Fiche bilan}{poporange}{L'essentiel à connaître pour l'examen}}}
% Figure encadrée avec légende
\newtcolorbox{popfigure}[1][]{enhanced,breakable=false,colback=white,colframe=popline,boxrule=1.4pt,arc=8pt,
  left=10pt,right=10pt,top=10pt,bottom=8pt,before skip=10pt,after skip=12pt,halign=center,
  lower separated=false,halign lower=center,
  fontlower=\popsemi\fontsize{9}{11.5}\selectfont\color{popmuted},
  #1}
\newcommand{\legende}[1]{\par\vspace{6pt}{\popsemi\fontsize{9}{11.5}\selectfont\color{popmuted}#1}}

% ============================================================
% Signature OnBuch+
% ============================================================
\newcommand{\poplogoinline}[1]{{\poplogo\fontsize{#1}{#1}\selectfont\color{popink}OnBuch\color{poporange}+}}
\newcommand{\signatureonbuch}{%
  \begin{tcolorbox}[enhanced,colback=popink,colframe=popink,boxrule=2.2pt,arc=10pt,
      drop shadow=poporange,left=14pt,right=14pt,top=10pt,bottom=10pt,before skip=0pt,after skip=0pt]
    \tikz[baseline=-0.7ex]\node[circle,fill=popgreen,draw=popcream,line width=1.3pt,minimum size=22pt,inner sep=0]{\popcheck{white}};%
    \hspace{9pt}\begin{minipage}[c]{0.62\linewidth}
      {\popxbold\fontsize{12}{13}\selectfont\color{popcream}Signé\ \textbullet\ {\poplogo OnBuch\color{poporange}+}}\par\vspace{2pt}
      {\raggedright\popsemi\fontsize{8}{10.5}\selectfont\color{popline}\MakeUppercase{Équipe pédagogique OnBuch+}\\\MakeUppercase{Conforme au programme officiel MINESEC}\par}
    \end{minipage}\hfill
    {\poplogo\fontsize{22}{22}\selectfont\color{popcream}OnBuch\color{poporange}+}
  \end{tcolorbox}%
}

% ============================================================
% Page de garde
% #1 titre · #2 matière · #3 niveau · #4 module · #5 n° leçon · #6 durée ·
% #7 description / objectifs (texte ou itemize)
% ============================================================
\newcommand{\pagegarde}[7]{%
  \thispagestyle{empty}
  \newgeometry{a4paper,margin=1.7cm,top=1.6cm,bottom=1.4cm}%
  \begin{tikzpicture}[remember picture,overlay]
    \fill[poporange!18] ([xshift=-3.2cm,yshift=-3.6cm]current page.north east) circle (4.6cm);
    \fill[poppurple!14] ([xshift=2.4cm,yshift=7.6cm]current page.south west) circle (3.8cm);
  \end{tikzpicture}%
  {\poplogo\fontsize{30}{30}\selectfont\color{popink}OnBuch\color{poporange}+}\hfill
  \raisebox{0.6ex}{\poppill{Cours signé \textbullet\ Premium}{popink}{popcream}}\par
  \vspace{1pt}\popsquiggle{poppurple}\par
  \vspace{8pt}
  \begin{tcolorbox}[enhanced,colback=poporange,colframe=popink,boxrule=2.8pt,arc=13pt,
      drop shadow=popink,left=18pt,right=18pt,top=12pt,bottom=13pt,after skip=10pt]
    \poppill{#2 \textbullet\ #3}{popink}{popcream}\hspace{5pt}\poppill{#4}{white}{popink}\par\vspace{8pt}
    {\popdisplay\fontsize{14}{14}\selectfont\color{popink}LEÇON #5}\par\vspace{4pt}
    {\raggedright\hyphenpenalty=10000\exhyphenpenalty=10000\popdisplay\fontsize{25}{27}\selectfont\color{popcream}\MakeUppercase{#1}\par}
    \vspace{8pt}
    {\popxbold\fontsize{9.5}{9.5}\selectfont\color{popink}\MakeUppercase{Durée conseillée : #6}}
  \end{tcolorbox}
  \begin{tcolorbox}[enhanced,colback=white,colframe=popink,boxrule=2.2pt,arc=10pt,
      drop shadow=popink,left=16pt,right=16pt,top=10pt,bottom=10pt,after skip=8pt]
    \poppill{Dans cette leçon}{poppurple}{white}\par\vspace{5pt}
    \popsemi\fontsize{9.3}{12.6}\selectfont\color{popink}#7
  \end{tcolorbox}
  \noindent\begin{minipage}[t]{0.487\linewidth}
    \begin{tcolorbox}[enhanced,colback=popgreen,colframe=popink,boxrule=2.2pt,arc=10pt,
        drop shadow=popink,left=11pt,right=11pt,top=10pt,bottom=10pt,height=3.1cm,valign=center]
      \tcbox[colback=white,colframe=popink,boxrule=1.3pt,arc=5pt,boxsep=0pt,nobeforeafter,
        left=3pt,right=3pt,top=3pt,bottom=3pt,valign=center]{\qrcode[height=1.9cm]{https://play.google.com/store/apps/details?id=com.luvvix.onbuch&hl=fr}}%
      \hspace{8pt}\begin{minipage}[c]{0.5\linewidth}
        {\popdisplay\fontsize{11}{12.5}\selectfont\color{white}\MakeUppercase{Télécharge OnBuch}}\par\vspace{4pt}
        {\raggedright\popsemi\fontsize{8.4}{11}\selectfont\color{white}Gratuit \textbullet\ Google Play\\Tuteur IA, annales, concours, résultats\par}
      \end{minipage}
    \end{tcolorbox}
  \end{minipage}\hfill
  \begin{minipage}[t]{0.487\linewidth}
    \begin{tcolorbox}[enhanced,colback=poppurple,colframe=popink,boxrule=2.2pt,arc=10pt,
        drop shadow=popink,left=11pt,right=11pt,top=10pt,bottom=10pt,height=3.1cm,valign=center]
      \tikz[baseline=-0.7ex]\node[circle,fill=white,draw=popink,line width=1.3pt,minimum size=1.6cm,inner sep=0]{\popdisplay\fontsize{24}{24}\selectfont\color{poppurple}+};%
      \hspace{8pt}\begin{minipage}[c]{0.6\linewidth}
        {\popdisplay\fontsize{11}{12.5}\selectfont\color{white}\MakeUppercase{Passe à OnBuch+}}\par\vspace{4pt}
        {\raggedright\popsemi\fontsize{8.4}{11}\selectfont\color{white}Tous les cours signés, Tuteur IA illimité, fiches PDF — onglet OnBuch+ de l'app\par}
      \end{minipage}
    \end{tcolorbox}
  \end{minipage}
  \vspace{8pt}
  \popcontenu
  \vfill
  \signatureonbuch
  \restoregeometry
  \newpage
}

% Rangée « Ce cours contient » (page de garde)
\newcommand{\poptile}[3]{% #1 couleur #2 gros texte #3 légende
  \begin{tcolorbox}[enhanced,colback=#1,colframe=popink,boxrule=2pt,arc=9pt,shadow={2.5pt}{-2.5pt}{0pt}{popink},
      left=6pt,right=6pt,top=9pt,bottom=9pt,halign=center,valign=center,height=1.75cm,width=\linewidth,nobeforeafter]
    {\popdisplay\fontsize{12.5}{13}\selectfont\color{white}\MakeUppercase{#2}}\par\vspace{3pt}
    {\centering\popsemi\fontsize{7.8}{9.5}\selectfont\color{white}#3\par}
  \end{tcolorbox}}
\newcommand{\popcontenu}{%
  \par\noindent{\popxboldls\fontsize{7.5}{7.5}\selectfont\color{popmuted}CE COURS SIGNÉ CONTIENT}\par\vspace{7pt}
  \noindent\begin{minipage}[t]{0.235\linewidth}\poptile{popblue}{Cours}{complet, illustré, pas à pas}\end{minipage}\hfill
  \begin{minipage}[t]{0.235\linewidth}\poptile{poporange}{Méthodes}{et exercices résolus}\end{minipage}\hfill
  \begin{minipage}[t]{0.235\linewidth}\poptile{poppink}{Exercices}{type examen + corrigés}\end{minipage}\hfill
  \begin{minipage}[t]{0.235\linewidth}\poptile{popgreen}{Fiche bilan}{l'essentiel à retenir}\end{minipage}\par}

% Sommaire
\newtcolorbox{sommaire}{enhanced,colback=white,colframe=popink,boxrule=2.2pt,arc=10pt,
  drop shadow=popink,left=18pt,right=18pt,top=13pt,bottom=13pt,before skip=6pt,after skip=20pt,
  before upper={\poppill{Sommaire}{poppurple}{white}\par\vspace{9pt}\popsemi\fontsize{10.6}{14.5}\selectfont\color{popink}}}

% Signature de fin de cours — poussée tout en bas de la dernière page
\newcommand{\findecours}{%
  \par\vfill\nopagebreak
  \begin{center}\popsquiggle{poporange}\end{center}\vspace{4pt}
  \signatureonbuch
}
\AtBeginDocument{\def\llbracket{\mathopen{[\mkern-3.5mu[}}\def\rrbracket{\mathclose{]\mkern-3.5mu]}}} % intervalles d'entiers (glyphe ⟦ absent de la police)
% Glyphes absents de Fira Math : redéfinis à partir de glyphes disponibles
\AtBeginDocument{%
  \def\setminus{\mathbin{\backslash}}\let\smallsetminus\setminus
  \def\vdots{\mathord{\vbox{\baselineskip3.2pt\lineskiplimit0pt\kern4pt\hbox{.}\hbox{.}\hbox{.}}}}%
  \def\ddots{\mathinner{\mkern1mu\raise6.5pt\hbox{.}\mkern2mu\raise3.6pt\hbox{.}\mkern2mu\raise0.7pt\hbox{.}\mkern1mu}}%
  \def\triangle{\mathord{\scalebox{0.9}{$\symup{\Delta}$}}}%
  \def\nabla{\mathord{\raisebox{\depth}{\scalebox{1}[-1]{$\symup{\Delta}$}}}}%
  \def\checkmark{\popcheck{popdark}}%
  \def\star{\mathbin{\ast}}\def\bigstar{\mathord{\ast}}%
  \def\diamond{\mathbin{\scalebox{0.6}{$\Diamond$}}}%
  \def\bigcup{\mathop{\mathchoice{\raisebox{-0.3em}{\scalebox{1.7}{$\cup$}}}{\scalebox{1.25}{$\cup$}}{\cup}{\cup}}\displaylimits}%
  \def\bigcap{\mathop{\mathchoice{\raisebox{-0.3em}{\scalebox{1.7}{$\cap$}}}{\scalebox{1.25}{$\cap$}}{\cap}{\cap}}\displaylimits}%
  \def\lesseqqgtr{\mathrel{\lessgtr}}\def\bigoplus{\mathop{\oplus}\displaylimits}%
}
\providecommand{\ssi}{si et seulement si} % abréviation fréquente
\AtBeginDocument{\providecommand{\wideparen}{\overparen}} % arc de cercle

%% FIGFIX-BEGIN v4 — correctifs globaux des figures (docs/onbuch-generation/tools/figfix)
\usepackage{adjustbox}\usepackage{etoolbox}
% toute figure TikZ/pgfplots plus large que la ligne est réduite à la largeur disponible
\BeforeBeginEnvironment{tikzpicture}{\begin{adjustbox}{max width=\linewidth}}
\AfterEndEnvironment{tikzpicture}{\end{adjustbox}}
% légendes pgfplots lisibles par-dessus les courbes
\pgfplotsset{every axis/.append style={legend style={fill=white,fill opacity=0.92,text opacity=1,draw=black!15,inner sep=3pt}}}
% étiquettes d'arêtes (syntaxe edge["..."]) avec fond blanc
\tikzset{every edge quotes/.append style={fill=white,inner sep=1.2pt}}
%% FIGFIX-END
\begin{document}

\pagegarde{Limitation des dysfonctionnements des structures responsables du contrôle de la motricité}{SVT}{Tle D}{Module I : Le Monde Vivant — Activités cérébrales et motricité volontaire}{N10}{4 h}{Cette leçon étudie l'organisation fonctionnelle de l'encéphale dans le contrôle de la motricité volontaire : localisation des aires motrices et sensitives, trajet de l'influx nerveux lors d'un mouvement volontaire, puis conséquences des affections de ces structures (Parkinson, Huntington, Alzheimer) et moyens de limitation des dysfonctionnements. Les mécanismes sont établis à partir de dissections, d'expériences de stimulations et de lésions corticales et de cas cliniques.
\vspace{4pt}
\begin{multicols}{2}\small
\begin{itemize}
\item À la fin de cette leçon, je sais définir un mouvement volontaire et le distinguer d'un mouvement réflexe.
\item À la fin de cette leçon, je sais disséquer l'encéphale d'un Mammifère, identifier ses principales structures et réaliser un schéma annoté.
\item À la fin de cette leçon, je sais localiser les aires motrices et sensitives du cortex cérébral et décrire leur organisation somatotopique.
\item À la fin de cette leçon, je sais interpréter des expériences de stimulation et de lésion corticales pour établir la fonction d'une aire cérébrale.
\item À la fin de cette leçon, je sais expliquer à l'aide d'un schéma fonctionnel le trajet de l'influx nerveux lors d'un mouvement volontaire.
\item À la fin de cette leçon, je sais décrire les symptômes et la cause des maladies de Parkinson, de Huntington et d'Alzheimer.
\end{itemize}\end{multicols}}

\begin{sommaire}
\begin{multicols}{2}
\begin{enumerate}
\item Le mouvement volontaire et l'encéphale : observations et dissection
\item Localisation des aires motrices et sensitives : preuves expérimentales
\item Trajet de l'influx nerveux lors d'un mouvement volontaire
\item Dysfonctionnements : maladie de Parkinson
\item Dysfonctionnements : maladies de Huntington et d'Alzheimer
\item Synthèse : relier structure, fonction et dysfonctionnement ; prévention
\item Activité d'intégration
\item Méthodes et exercices résolus
\item Exercices
\item Corrigés des exercices
\item Fiche bilan
\end{enumerate}
\end{multicols}
\end{sommaire}

\begin{objectifs}
\begin{itemize}
\item À la fin de cette leçon, je sais définir un mouvement volontaire et le distinguer d'un mouvement réflexe.
\item À la fin de cette leçon, je sais disséquer l'encéphale d'un Mammifère, identifier ses principales structures et réaliser un schéma annoté.
\item À la fin de cette leçon, je sais localiser les aires motrices et sensitives du cortex cérébral et décrire leur organisation somatotopique.
\item À la fin de cette leçon, je sais interpréter des expériences de stimulation et de lésion corticales pour établir la fonction d'une aire cérébrale.
\item À la fin de cette leçon, je sais expliquer à l'aide d'un schéma fonctionnel le trajet de l'influx nerveux lors d'un mouvement volontaire.
\item À la fin de cette leçon, je sais décrire les symptômes et la cause des maladies de Parkinson, de Huntington et d'Alzheimer.
\item À la fin de cette leçon, je sais relier un symptôme moteur ou cognitif à la structure encéphalique atteinte à partir d'un cas clinique.
\item À la fin de cette leçon, je sais citer les moyens de prévention et de limitation des dysfonctionnements du contrôle de la motricité.
\end{itemize}
\end{objectifs}

\begin{prerequis}
\begin{itemize}
\item Structure du système nerveux : neurone, influx nerveux, synapse (Première)
\item Arc réflexe et message nerveux moteur (Première)
\item Organisation générale du système nerveux central : encéphale et moelle épinière (Première)
\item Notion de neurotransmetteur et de transmission synaptique (début de Terminale)
\item Lecture et interprétation de résultats expérimentaux et de courbes (méthodologie SVT)
\end{itemize}
\end{prerequis}

\newpage

\coursec{Le mouvement volontaire et l'encéphale : observations et dissection}

\courssub{Situation d'accroche et problématique}

À Yaoundé, un ancien joueur de football de 62 ans consulte à l'hôpital central : sa main droite tremble au repos, ses gestes deviennent lents et sa démarche se fige parfois, comme si ses pieds étaient collés au sol. Le médecin évoque la \cle{maladie de Parkinson}. Quelques mois plus tard, à Bonabéri (Douala), un jeune moto-taxi victime d'un accident de la route ne peut plus lever le bras gauche... alors que son bras n'est \emph{pas blessé} : la radiographie du membre est normale, mais l'imagerie cérébrale révèle une lésion d'une zone précise du cerveau, du côté \emph{droit}.

Ces deux situations cliniques, fréquentes dans nos hôpitaux, posent des questions fondamentales : comment le cerveau commande-t-il nos mouvements volontaires ? Pourquoi une lésion cérébrale ou la dégénérescence de certaines structures nerveuses entraîne-t-elle des troubles moteurs si précis et si localisés ? Peut-on prévenir ou limiter ces dysfonctionnements ?

Pour répondre à ces questions, il faut d'abord définir ce qu'est un mouvement volontaire, puis identifier l'organe qui le commande : l'\cle{encéphale}. C'est l'objet de cette première partie, qui s'appuie sur une démarche d'investigation : observation, dissection, hypothèse.

\courssub{Qu'est-ce qu'un mouvement volontaire ?}

\textbf{Intuition et observations de la vie courante.}

Observe autour de toi : un élève attrape un ballon lancé vers lui, un autre écrit sur son cahier, une vendeuse au marché Mokolo frappe à la porte d'une boutique. Ces mouvements ont un point commun : ils sont \emph{décidés}. La personne choisit le moment, la direction, la force et la durée du geste. Elle peut l'interrompre à volonté. On dit que ces mouvements sont \cle{volontaires}.

Compare maintenant avec l'expérience classique du \cle{réflexe rotulien} : assis, jambe croisée, on frappe doucement le tendon situé sous la rotule avec le bord de la main : la jambe se détend brusquement, \emph{sans que la personne l'ait décidé}. Mieux : même en essayant de l'empêcher, on n'y parvient qu'imparfaitement. Ce mouvement est \emph{involontaire} : c'est un \cle{réflexe}.

\begin{definition}[Mouvement volontaire]
Un \cle{mouvement volontaire} est un mouvement \textbf{déclenché et contrôlé consciemment} par le sujet : il est décidé, orienté vers un but, et peut être modulé ou interrompu à tout moment. Son centre de commande se situe dans le \textbf{cortex cérébral} (cerveau), au niveau des \cle{aires motrices}. Il met en jeu des muscles squelettiques (muscles striés volontaires).
\end{definition}

\begin{definition}[Mouvement réflexe]
Un \cle{réflexe} est une réponse motrice \textbf{involontaire, rapide et stéréotypée} à une stimulation. Son centre nerveux se situe le plus souvent dans la \textbf{moelle épinière} (réflexes spinaux, comme le réflexe rotulien) ou dans le tronc cérébral, et non dans le cortex cérébral.
\end{definition}

\begin{experience}[Mise en évidence : le réflexe persiste sans commande cérébrale]
\textbf{Protocole.} Chez une grenouille dont on a détruit le cerveau (animal dit \og spinal \fg, protocole historique de laboratoire), on pince la peau d'une patte ou on la met en contact avec une solution acide diluée.

\textbf{Observations.} La patte se retire vivement. Le mouvement se produit alors que le cerveau est détruit.

\textbf{Interprétation.} La commande de ce mouvement ne transite \emph{pas} par le cerveau : elle est assurée par la \textbf{moelle épinière}, qui reçoit le message sensitif et émet directement l'ordre moteur (c'est l'\cle{arc réflexe}). En revanche, chez l'animal intact, attraper une proie ou sauter vers un endroit précis exige le cerveau : ce sont des mouvements volontaires.

\textbf{Conclusion.} Il existe deux niveaux de commande motrice : la moelle épinière suffit pour les réflexes ; le cerveau est indispensable pour les mouvements volontaires.
\end{experience}

\begin{attention}[Ne pas confondre mouvement volontaire et réflexe]
Le piège classique du Baccalauréat : affirmer que \og tout mouvement est commandé par le cerveau \fg. C'est faux ! Le \textbf{réflexe} ne transite \emph{pas obligatoirement} par le cerveau pour sa commande : son centre est médullaire (moelle épinière). Le cerveau est informé \emph{après coup} (c'est pourquoi tu ressens la douleur après avoir retiré ta main d'une flamme), mais il n'a pas commandé le retrait. À l'inverse, le mouvement volontaire naît \emph{dans le cortex cérébral}. Autre erreur : croire qu'un réflexe est \og inconscient parce que lent \fg : c'est l'inverse, le réflexe est plus rapide car son circuit est plus court.
\end{attention}

\begin{center}
\begin{tabularx}{\linewidth}{|l|X|X|}
\hline
\rowcolor{popblueL} \textbf{Critère} & \textbf{Mouvement volontaire} & \textbf{Mouvement réflexe} \\
\hline
Déclenchement & Décidé consciemment & Automatique, provoqué par un stimulus \\
\hline
Centre de commande & Cortex cérébral (aires motrices) & Moelle épinière (ou tronc cérébral) \\
\hline
Vitesse de réponse & Plus lente (circuit long) & Très rapide (circuit court) \\
\hline
Caractère & Modulable, interruptible à volonté & Stéréotypé, invariable \\
\hline
Exemples & Attraper un ballon, écrire, frapper à une porte & Réflexe rotulien, retrait de la main d'un objet brûlant \\
\hline
\end{tabularx}
\end{center}

\courssub{Travail pratique : dissection de l'encéphale d'un Mammifère}

Puisque le cerveau commande la motricité volontaire, il faut connaître son organisation. L'observation directe d'un encéphale de Mammifère (mouton, chèvre ou porc, organes disponibles auprès d'un abattoir) est un savoir-faire exigible du programme.

\begin{methode}[Protocole de dissection de l'encéphale d'un Mammifère]
\begin{enumerate}
\item \textbf{Sécurité d'abord.} Porter des gants, une blouse et des lunettes. L'organe doit être frais et provenir d'un animal sain (abattoir contrôlé). Manipuler les instruments tranchants (scalpel, ciseaux) en dirigeant toujours la lame vers l'extérieur. Se laver les mains après la séance.
\item \textbf{Matériel.} Encéphale de mouton (ou de chèvre/porc), bac à dissection, scalpel, ciseaux fins, pinces, aiguilles à disséquer, gants, eau physiologique (\SI{9}{\g\per\litre} de \ce{NaCl}) pour humidifier l'organe.
\item \textbf{Dégagement de l'organe.} Si l'encéphale est encore dans la boîte crânienne, l'enseignant l'extrait en sectionnant les nerfs crâniens et en coupant la moelle épinière au niveau du trou occipital, afin de préserver le \textbf{bulbe rachidien}.
\item \textbf{Observation de la face dorsale.} Écarter délicatement les méninges (enveloppes protectrices) si elles sont présentes. Observer : les \textbf{deux hémisphères cérébraux} séparés par une fente longitudinale profonde ; leur surface plissée en \textbf{circonvolutions} séparées par des sillons (scissures) ; en arrière et en dessous, le \textbf{cervelet}, très finement plissé.
\item \textbf{Observation de la face ventrale.} Retourner l'organe. Observer : le \textbf{tronc cérébral} (pont, renflement transversal, puis bulbe rachidien) se continuant avec la \textbf{moelle épinière} ; les nerfs crâniens (dont les nerfs optiques qui se croisent en chiasma).
\item \textbf{Coupe sagittale médiane.} À l'aide d'un scalpel long et tranchant, sectionner l'encéphale en deux moitiés symétriques, d'un seul geste ferme, dans le plan médian. Observer sur la coupe : la substance grise périphérique (cortex) et la substance blanche centrale ; le corps calleux (pont de substance blanche reliant les deux hémisphères) ; le cervelet et son aspect ramifié (\og arbre de vie \fg) ; le tronc cérébral.
\item \textbf{Schéma annoté.} Réaliser un schéma propre de la vue sagittale et de la vue dorsale (voir règles ci-dessous).
\end{enumerate}
\end{methode}

\begin{methode}[Règles de réalisation d'un schéma annoté en SVT]
\begin{enumerate}
\item Dessiner au \textbf{crayon à papier}, à \textbf{traits nets et continus}, sans ombres ni hachures inutiles.
\item Donner un \textbf{titre} précis au schéma (ex. : \og Encéphale de mouton en coupe sagittale médiane \fg), souligné.
\item Indiquer les légendes par des \textbf{traits horizontaux ou obliques} qui ne se croisent jamais, aboutissant exactement à la structure nommée ; les noms sont alignés verticalement dans la marge.
\item Respecter les \textbf{proportions} réelles et la position relative des structures.
\item Ne jamais reproduire de couleurs fantaisistes : un schéma d'observation se fait en noir et blanc.
\end{enumerate}
\end{methode}

\courssub{Structures observées : description de l'encéphale}

L'observation de l'encéphale de mouton, transposable à l'Homme, permet d'identifier quatre grandes structures :

\begin{itemize}
\item Le \cle{cerveau}, partie la plus volumineuse, formé de \textbf{deux hémisphères cérébraux} (droit et gauche) séparés par la fente longitudinale cérébrale. Sa surface n'est pas lisse : elle est plissée en \textbf{circonvolutions} (replis) séparées par des \textbf{sillons}. La couche superficielle, grise, est le \cle{cortex cérébral} (substance grise, riche en corps cellulaires des neurones) ; l'intérieur est formé de substance blanche (faisceaux d'axones myélinisés).
\item Le \cle{cervelet}, situé en arrière et sous le cerveau, présente une surface \textbf{très finement plissée} (lamelles parallèles). Il joue un rôle essentiel dans la coordination et l'équilibre des mouvements.
\item Le \cle{tronc cérébral}, qui relie le cerveau à la moelle épinière : il comprend le \textbf{pont} (renflement visible sur la face ventrale) et le \textbf{bulbe rachidien} (élargissement juste au-dessus de la moelle). Le bulbe contient des centres vitaux (contrôle du rythme cardiaque et de la respiration).
\item La \cle{moelle épinière}, cordon nerveux prolongeant le bulbe rachidien dans le canal vertébral : elle est le centre des réflexes et la voie de passage des messages entre l'encéphale et le reste du corps.
\end{itemize}

\begin{popfigure}
\begin{center}
\begin{tikzpicture}[scale=1.0, every node/.style={font=\small}]
% Cerveau (grand contour)
\draw[thick, popblue, fill=popblueL]
  (0.2,0.6) .. controls (0.2,2.2) and (1.2,3.4) .. (3.0,3.4)
  .. controls (5.0,3.4) and (6.2,2.4) .. (6.2,1.5)
  .. controls (6.2,1.0) and (5.9,0.8) .. (5.5,0.8)
  .. controls (4.6,0.8) and (4.2,1.0) .. (3.6,0.9)
  .. controls (2.4,0.7) and (1.2,0.4) .. (0.2,0.6) -- cycle;
% Circonvolutions (sillons)
\draw[popblue!70] (1.0,2.9) .. controls (1.6,2.5) and (1.9,2.9) .. (2.5,2.6);
\draw[popblue!70] (2.8,3.1) .. controls (3.4,2.7) and (3.8,3.0) .. (4.4,2.7);
\draw[popblue!70] (4.7,2.9) .. controls (5.2,2.5) and (5.5,2.4) .. (5.9,2.1);
\draw[popblue!70] (1.2,1.9) .. controls (1.9,1.6) and (2.4,1.9) .. (3.0,1.6);
\draw[popblue!70] (3.4,1.9) .. controls (4.0,1.6) and (4.6,1.8) .. (5.2,1.5);
% Cervelet
\draw[thick, popgreen, fill=popgreenL]
  (5.3,0.75) .. controls (5.6,1.3) and (6.6,1.35) .. (7.0,0.9)
  .. controls (7.3,0.5) and (7.1,-0.1) .. (6.5,-0.3)
  .. controls (5.8,-0.55) and (5.2,-0.2) .. (5.1,0.3)
  .. controls (5.05,0.5) and (5.15,0.65) .. (5.3,0.75) -- cycle;
\draw[popgreen!80] (5.5,0.6) .. controls (5.9,0.45) and (6.4,0.5) .. (6.9,0.7);
\draw[popgreen!80] (5.5,0.25) .. controls (6.0,0.1) and (6.5,0.15) .. (7.0,0.35);
\draw[popgreen!80] (5.7,-0.1) .. controls (6.2,-0.2) and (6.6,-0.15) .. (6.9,0.0);
% Tronc cérébral + moelle
\draw[thick, poporange, fill=poporangeL]
  (4.6,0.85) .. controls (5.0,0.5) and (5.2,0.1) .. (5.15,-0.5)
  .. controls (5.1,-1.0) and (5.0,-1.5) .. (4.95,-2.1)
  -- (4.45,-2.1) .. controls (4.5,-1.4) and (4.55,-0.9) .. (4.5,-0.4)
  .. controls (4.45,0.1) and (4.3,0.5) .. (4.0,0.75) -- cycle;
% Pont (renflement)
\draw[thick, poporange] (4.55,-0.35) .. controls (4.9,-0.15) and (5.15,-0.2) .. (5.2,-0.5);
% Légendes
\draw[popdark, ->] (3.0,3.35) -- (3.0,4.1) node[above]{\textbf{1. Cerveau} (hémisphères cérébraux, circonvolutions)};
\draw[popdark, ->] (6.4,0.4) -- (8.2,0.9) node[right, text width=4.6cm]{\textbf{2. Cervelet} (surface très plissée)};
\draw[popdark, ->] (5.0,-0.4) -- (8.2,-0.6) node[right, text width=4.6cm]{\textbf{3. Tronc cérébral} (pont et bulbe rachidien)};
\draw[popdark, ->] (4.7,-1.9) -- (8.2,-2.1) node[right, text width=4.6cm]{\textbf{4. Moelle épinière}};
\draw[popdark, ->] (1.0,0.62) -- (-0.6,-0.6) node[below left, text width=4.2cm]{\textbf{5. Cortex cérébral} (substance grise périphérique)};
\end{tikzpicture}
\end{center}
\legende{Schéma annoté de l'encéphale de Mammifère en coupe sagittale médiane. 1 : cerveau (deux hémisphères à surface plissée en circonvolutions) ; 2 : cervelet ; 3 : tronc cérébral (pont puis bulbe rachidien) ; 4 : moelle épinière ; 5 : cortex cérébral, siège des aires motrices et sensitives.}
\end{popfigure}

\courssub{De l'observation à l'hypothèse : la démarche d'investigation}

\textbf{Observation.} La dissection montre que le cerveau n'est pas lisse : ses circonvolutions multiplient la surface du cortex. Le cervelet, lui aussi, est extrêmement plissé. De plus, une comparaison entre espèces est instructive : le cerveau du rat est presque lisse, celui du mouton est plissé, celui de l'Homme l'est encore davantage. Or, parmi ces espèces, l'Homme est celle dont les comportements moteurs volontaires sont les plus élaborés (écriture, parole, outils, sport).

\textbf{Question.} Pourquoi le cortex cérébral est-il plissé ? Ce plissement a-t-il un lien avec les fonctions élaborées ?

\textbf{Hypothèse.} La complexité de la surface cérébrale (circonvolutions) serait liée à des fonctions élaborées comme la motricité volontaire : les plis permettraient de loger un \emph{plus grand nombre de neurones corticaux} dans un volume crânien limité.

\textbf{Vérification.} On peut estimer la surface du cortex : chez l'Homme, les deux tiers environ du cortex sont cachés au fond des sillons. Le plissement augmente donc considérablement la surface corticale sans augmenter le volume du crâne. Cette hypothèse sera confirmée dans la suite de la leçon par les expériences de stimulation corticale : c'est bien dans le cortex que se trouvent les aires qui commandent les mouvements volontaires.

\begin{savaistu}[Un cortex grand comme une serviette]
Chez l'Homme, le cortex cérébral, s'il était déplié, couvrirait une surface d'environ \SI{0,25}{m^2} (soit environ \SI{2500}{cm^2}, la taille d'une grande serviette ou d'une feuille de journal ouverte), pour une épaisseur de seulement 2 à \SI{4}{mm}. Il contient de l'ordre de 15 à 20 milliards de neurones. Les circonvolutions sont donc une solution \og d'emballage \fg : loger une immense surface de traitement de l'information dans un crâne de volume limité. C'est cette richesse corticale qui rend possibles les gestes volontaires fins, comme ceux du footballeur contrôlant un ballon ou du griot jouant du mvet.
\end{savaistu}

\courssub{Deux circuits, deux niveaux de contrôle}

L'ensemble de ces observations permet de construire un schéma fonctionnel comparant le circuit du réflexe (commande médullaire) et celui du mouvement volontaire (commande cérébrale). Ce schéma est la clé pour comprendre pourquoi une lésion cérébrale abolit un geste volontaire alors que le réflexe du même membre peut persister.

\begin{popfigure}
\begin{center}
\begin{tikzpicture}[font=\small, >=Stealth,
  box/.style={draw, thick, rounded corners=2pt, align=center, minimum height=0.9cm, text width=3.1cm},
  lab/.style={font=\small\bfseries}]
% --- Arc réflexe (gauche) ---
\node[lab, popgreen] at (1.9,3.6) {A. Arc réflexe (involontaire)};
\node[box, fill=popgreenL, draw=popgreen] (rec) at (0,2.4) {Récepteur (peau, muscle)};
\node[box, fill=popgreenL, draw=popgreen] (moelle) at (3.8,2.4) {Moelle épinière (centre du réflexe)};
\node[box, fill=popgreenL, draw=popgreen] (eff) at (0,0.6) {Muscle effecteur};
\draw[->, thick, popgreen] (rec) -- node[above, text width=2.6cm, align=center]{message sensitif (nerf sensitif)} (moelle);
\draw[->, thick, popgreen] (moelle) -- node[right, text width=2.8cm, align=center]{message moteur (nerf moteur)} (eff);
% --- Circuit volontaire (droite) ---
\node[lab, popblue] at (10.4,3.6) {B. Mouvement volontaire};
\node[box, fill=popblueL, draw=popblue] (cortex) at (10.4,2.4) {Cortex cérébral (aire motrice) : décision};
\node[box, fill=poporangeL, draw=poporange] (moelle2) at (10.4,0.6) {Moelle épinière (relais)};
\node[box, fill=popblueL, draw=popblue] (eff2) at (10.4,-1.2) {Muscle effecteur};
\draw[->, thick, popblue] (cortex) -- node[right, text width=3.4cm, align=center]{influx moteur volontaire (voies pyramidales)} (moelle2);
\draw[->, thick, popblue] (moelle2) -- node[right]{nerf moteur} (eff2);
% Information ascendante vers le cerveau dans le réflexe
\draw[->, dashed, popmuted] (moelle.north) .. controls (5.5,4.4) and (8.0,4.4) .. node[above, text width=4.2cm, align=center]{le cerveau est informé (sensation) mais ne commande pas} (cortex.north);
\end{tikzpicture}
\end{center}
\legende{Schéma comparatif des deux circuits moteurs. A : dans l'arc réflexe, le centre de commande est la moelle épinière ; le circuit est court, la réponse rapide et involontaire. B : dans le mouvement volontaire, la commande naît dans le cortex cérébral (aire motrice), descend vers la moelle épinière puis gagne le muscle. En pointillés : lors d'un réflexe, le cerveau reçoit l'information mais n'a pas déclenché le mouvement.}
\end{popfigure}

\begin{exoresolu}[Réflexe ou mouvement volontaire ?]
\textbf{Énoncé.} Pour chacune des situations suivantes, indique s'il s'agit d'un mouvement volontaire ou d'un réflexe, en justifiant par le centre nerveux impliqué : (a) un élève lève la main pour répondre à une question ; (b) un enfant retire vivement le pied après avoir marché sur une épingle ; (c) un chauffeur de taxi à Douala freine en voyant un piéton ; (d) la jambe se détend quand le médecin frappe le tendon rotulien.
\tcblower
\textbf{Solution détaillée.}
\begin{itemize}
\item (a) \textbf{Mouvement volontaire} : le geste est décidé consciemment ; la commande naît dans le cortex cérébral (aire motrice).
\item (b) \textbf{Réflexe} (réflexe de retrait) : la réponse est rapide, involontaire et protectrice ; le centre de commande est la moelle épinière. La douleur n'est ressentie qu'après le retrait, quand l'information atteint le cerveau.
\item (c) \textbf{Mouvement volontaire} : freiner est un geste décidé (même s'il est devenu automatique par l'apprentissage) ; il implique le cortex cérébral. Attention : sa rapidité ne doit pas faire croire à un réflexe.
\item (d) \textbf{Réflexe} (réflexe rotulien) : mouvement involontaire, stéréotypé, dont le centre est la moelle épinière.
\end{itemize}
\textbf{Piège à éviter} : ne pas se fier à la rapidité du geste pour conclure ; le critère décisif est le \emph{centre de commande} (cerveau ou moelle) et le caractère décidé ou non du mouvement.
\end{exoresolu}

\begin{exercice}[Pour t'entraîner]
Un patient présente, après une lésion de la moelle épinière au niveau lombaire, une abolition du réflexe rotulien de la jambe droite, mais il peut encore, sur commande, tenter de contracter les muscles de la cuisse. Explique ces observations en précisant, pour chaque fonction, le centre nerveux impliqué et l'état des voies nerveuses concernées.
\end{exercice}

\begin{corrige}[Exercice : réflexe et commande volontaire après lésion médullaire]
Le réflexe rotulien est un réflexe spinal : son arc (récepteur, nerf sensitif, centre médullaire lombaire, nerf moteur, muscle) passe par la moelle épinière au niveau lombaire. La lésion détruit ou interrompt ce centre (ou ses racines nerveuses) : le réflexe est donc aboli. En revanche, la tentative de contraction volontaire des muscles de la cuisse montre que le cortex cérébral (aire motrice) est intact et que les voies descendantes vers les motoneurones des muscles de la cuisse (innervés plus haut) sont encore fonctionnelles. La dissociation observée illustre les deux niveaux de commande : médullaire pour le réflexe, cérébral pour le mouvement volontaire.
\end{corrige}

\begin{aretenir}[L'essentiel de la section]
\begin{itemize}
\item Un \cle{mouvement volontaire} est décidé et contrôlé consciemment ; sa commande naît dans le \textbf{cortex cérébral} (aires motrices). Un \cle{réflexe} est involontaire, rapide, et son centre est la \textbf{moelle épinière}.
\item L'\cle{encéphale} comprend le \textbf{cerveau} (deux hémisphères plissés en circonvolutions), le \textbf{cervelet} (très plissé, coordination motrice) et le \textbf{tronc cérébral} (pont, bulbe rachidien), prolongé par la \textbf{moelle épinière}.
\item Les circonvolutions augmentent la surface du cortex (environ \SI{0,25}{m^2} chez l'Homme) : hypothèse d'un lien entre complexité corticale et fonctions élaborées comme la motricité volontaire.
\item La dissection de l'encéphale de Mammifère (mouton, chèvre, porc) et la réalisation d'un schéma annoté (traits nets, légendes alignées, titre) sont des savoir-faire exigibles.
\end{itemize}
\end{aretenir}

Maintenant que nous savons \emph{ce qu'est} un mouvement volontaire et \emph{où} se situe son organe de commande, une question demeure : comment a-t-on prouvé que certaines zones précises du cortex commandent le mouvement ? C'est l'objet de la section suivante : la localisation des aires motrices et sensitives, établie par des expériences de stimulation et de lésions corticales.

\coursec{Localisation des aires motrices et sensitives : preuves expérimentales}

Dans la section précédente, la dissection de l'encéphale de Mammifère nous a permis d'identifier les grandes structures : les deux hémisphères cérébraux recouverts du cortex plissé en circonvolutions, le cervelet, le tronc cérébral. Mais une question demeure, et elle est fondamentale : ce cortex cérébral, qui apparaît à l'œil nu comme une masse grise uniforme, est-il \cle{fonctionnellement homogène} ? Autrement dit, n'importe quelle région du cortex peut-elle commander n'importe quel mouvement ? Ou bien, au contraire, le cortex est-il divisé en \cle{aires spécialisées}, chacune dédiée à une fonction précise ? C'est à cette question que répondent les expériences de stimulation et de lésion que nous allons étudier maintenant.

\courssub{Le problème posé : un cortex uniforme ou un cortex cartographié ?}

\begin{experience}[Situation-problème]
\textbf{Observation clinique.} Au Centre Hospitalier d'Essos à Yaoundé, deux patients victimes d'un accident vasculaire cérébral (AVC) présentent des symptômes différents : le premier ne peut plus bouger le bras droit, le second ne sent plus le contact des objets avec la main gauche, bien que ses muscles fonctionnent normalement.

\textbf{Question.} Si le cortex était homogène, une lésion de même taille devrait provoquer les mêmes troubles quel que soit son siège. Or les symptômes dépendent de la \emph{localisation} de la lésion.

\textbf{Hypothèse.} Le cortex cérébral est divisé en \cle{aires fonctionnelles spécialisées} : certaines régions commandent le mouvement (aires motrices), d'autres reçoivent les informations sensorielles (aires sensitives).

\textbf{Vérification expérimentale.} Deux approches complémentaires sont possibles : la \emph{stimulation} (on excite un point du cortex et on observe la réponse) et la \emph{lésion} (on détruit une région et on observe la fonction perdue).
\end{experience}

\courssub{Les expériences de stimulation électrique du cortex}

\textbf{Les expériences de Fritsch et Hitzig (1870) chez le chien.}

Les physiologistes allemands Gustav Fritsch et Eduard Hitzig furent les premiers à tester expérimentalement l'hypothèse de la spécialisation corticale. Ils appliquèrent de faibles courants électriques sur la surface du cortex d'un chien anesthésié, dont la boîte crânienne avait été ouverte.

\begin{experience}[Stimulation du cortex précentral chez le chien (Fritsch et Hitzig, 1870)]
\textbf{Protocole.} On stimule électriquement, un à un, différents points de la circonvolution située \emph{en avant} du sillon central (sillon de Rolando) du cortex gauche du chien.

\textbf{Observations.}
\begin{itemize}
\item La stimulation d'un point précis provoque la contraction d'un \emph{groupe musculaire précis} (par exemple les muscles de la patte) ;
\item la réponse apparaît du \emph{côté opposé} du corps : stimuler le cortex gauche fait bouger la patte droite ;
\item à chaque point stimulé correspond toujours le même groupe musculaire : la réponse est \emph{reproductible} ;
\item la stimulation d'autres régions du cortex ne provoque aucun mouvement.
\end{itemize}

\textbf{Interprétation.} Il existe une région corticale spécialisée dans la commande du mouvement : c'est l'\cle{aire motrice}. À l'intérieur de cette aire, les groupes musculaires sont représentés de façon ordonnée : c'est la \cle{somatotopie}.
\end{experience}

\textbf{Les travaux de Penfield (XX\textsuperscript{e} siècle) chez l'Homme.}

Le neurochirurgien canadien Wilder Penfield étendit ces résultats à l'Homme. Au cours d'opérations chirurgicales du cerveau (traitement de l'épilepsie), il stimula électriquement le cortex de patients \emph{éveillés}, sous anesthésie locale uniquement, et recueillit leurs réponses. L'avantage décisif par rapport à l'animal : le patient peut \emph{décrire} ce qu'il ressent, ce qui permet d'explorer aussi bien la motricité que la sensibilité.

\begin{center}
\begin{tabularx}{\linewidth}{|l|X|X|}
\hline
\rowcolor{popblueL} \textbf{Point stimulé} & \textbf{Réponse observée} & \textbf{Conclusion} \\
\hline
Circonvolution précentrale (partie supérieure) & Mouvement du pied et de la jambe du côté opposé & Aire motrice : commande des membres inférieurs \\
\hline
Circonvolution précentrale (partie moyenne) & Mouvement de la main, des doigts du côté opposé & Aire motrice : commande de la main \\
\hline
Circonvolution précentrale (partie inférieure) & Mouvement des lèvres, de la langue, du visage & Aire motrice : commande de la face \\
\hline
Circonvolution postcentrale & Le patient déclare des picotements localisés (main, lèvres...) du côté opposé & Aire sensitive : réception des sensations somatiques \\
\hline
Autres régions du cortex & Aucun mouvement, aucune sensation élémentaire & Aires non motrices et non sensitives élémentaires \\
\hline
\end{tabularx}
\end{center}

\textbf{Interprétation du tableau.} Les réponses motrices ne sont obtenues que lors de la stimulation de la circonvolution précentrale, et les sensations que lors de la stimulation de la circonvolution postcentrale. Le cortex n'est donc pas homogène : il est divisé en \cle{aires spécialisées}. De plus, à l'intérieur de chaque aire, le corps est représenté de manière ordonnée (pied en haut, face en bas) : il existe une \emph{carte du corps} dans le cortex.

\begin{definition}[Aire motrice]
L'\cle{aire motrice} (ou \cle{cortex moteur}) est la région du cortex cérébral située dans la \emph{circonvolution précentrale} (en avant du sillon central), dont la stimulation électrique provoque la contraction de groupes musculaires précis du côté opposé du corps, et dont la lésion entraîne la paralysie de ces mêmes muscles. Elle est le siège du \emph{commandement} des mouvements volontaires.
\end{definition}

\begin{definition}[Aire sensitive]
L'\cle{aire sensitive} (ou \cle{cortex somesthésique}) est la région du cortex cérébral située dans la \emph{circonvolution postcentrale} (en arrière du sillon central), qui reçoit les informations sensorielles somatiques (tact, pression, température, douleur, position des membres) provenant du côté opposé du corps. Sa stimulation provoque des sensations localisées ; sa lésion entraîne une perte de sensibilité localisée.
\end{definition}

\begin{popfigure}
\begin{tikzpicture}[scale=1.0]
% Contour de l'hémisphère cérébral (vue latérale)
\draw[thick,popdark,fill=popcream] (0,0)
  .. controls (-0.4,1.4) and (0.4,2.9) .. (2.2,3.1)
  .. controls (4.6,3.4) and (6.6,2.6) .. (6.9,1.3)
  .. controls (7.1,0.5) and (6.5,0.1) .. (5.6,0.05)
  .. controls (4.4,0) and (3.6,-0.35) .. (2.6,-0.3)
  .. controls (1.4,-0.25) and (0.4,-0.5) .. (0,0) -- cycle;
% Sillon central
\draw[very thick,popink] (3.9,3.05) .. controls (3.5,2.3) and (3.35,1.5) .. (3.3,0.6);
\node[popink,font=\small,align=center] at (2.55,0.35) {Sillon central\\(de Rolando)};
% Aire motrice (précentrale)
\draw[very thick,poporange,fill=poporange!35] (3.15,3.0) .. controls (2.75,2.25) and (2.6,1.5) .. (2.55,0.65)
   -- (3.3,0.6) .. controls (3.35,1.5) and (3.5,2.3) .. (3.9,3.05) -- cycle;
% Aire sensitive (postcentrale)
\draw[very thick,popblue,fill=popblue!30] (3.9,3.05) .. controls (3.5,2.3) and (3.35,1.5) .. (3.3,0.6)
   -- (4.05,0.55) .. controls (4.1,1.45) and (4.25,2.25) .. (4.65,3.0) -- cycle;
% Circonvolution postcentrale : deuxième limite
\draw[thick,popblue] (4.65,3.0) .. controls (4.25,2.25) and (4.1,1.45) .. (4.05,0.55);
% Légendes fléchées
\draw[->,poporange,thick] (2.85,1.9) -- (1.5,2.6) node[left,font=\small\bfseries,align=right,text=poporange]{Circonvolution\\précentrale\\AIRE MOTRICE};
\draw[->,popblue,thick] (4.0,1.9) -- (5.6,2.7) node[right,font=\small\bfseries,align=left,text=popblue]{Circonvolution\\postcentrale\\AIRE SENSITIVE};
% Orientation
\node[font=\footnotesize\itshape,text=popmuted] at (0.9,3.3) {Avant (frontal)};
\node[font=\footnotesize\itshape,text=popmuted] at (6.2,3.3) {Arrière (occipital)};
\end{tikzpicture}
\legende{Vue latérale d'un hémisphère cérébral humain. La circonvolution précentrale (orange), située en avant du sillon central, constitue l'aire motrice ; la circonvolution postcentrale (bleu), située en arrière, constitue l'aire sensitive. Les flèches indiquent les deux aires fonctionnelles.}
\end{popfigure}

\courssub{Les expériences de lésion : la preuve par la fonction abolie}

La stimulation montre qu'une structure est \emph{capable} de déclencher une fonction. Pour prouver qu'elle en est \emph{le siège obligé}, on réalise des expériences de lésion (ablation chirurgicale chez l'animal, ou étude des conséquences de lésions accidentelles chez l'Homme).

\begin{experience}[Ablation de l'aire motrice chez le singe]
\textbf{Protocole.} On détruit chirurgicalement, chez un singe, la partie de la circonvolution précentrale gauche dont la stimulation provoquait le mouvement de la main droite.

\textbf{Observations.} À son réveil, l'animal présente une \emph{paralysie de la main droite} : il ne peut plus saisir les objets avec cette main. La main gauche fonctionne normalement. Les muscles de la main droite ne sont pas atteints : ils se contractent si on stimule directement le nerf moteur qui les innerve.

\textbf{Interprétation.} La paralysie ne vient ni des muscles ni des nerfs, mais de la destruction du centre de commande cortical. L'aire motrice est donc bien le \emph{siège de la commande} du mouvement volontaire. Le déficit touche le côté \emph{opposé} à la lésion : le contrôle est croisé.

\textbf{Expérience complémentaire.} La destruction de la circonvolution postcentrale entraîne, elle, une \emph{perte de la sensibilité tactile} localisée du côté opposé : l'animal ne réagit plus au toucher de la région correspondante, bien qu'il puisse la mouvoir normalement. Motricité et sensibilité sont donc dissociables : elles relèvent de deux aires distinctes.
\end{experience}

\begin{methode}[Le raisonnement expérimental : stimulation et lésion]
Pour attribuer une fonction à une structure nerveuse, on combine deux approches complémentaires :
\begin{enumerate}
\item \textbf{La stimulation} : on excite artificiellement la structure (stimulation électrique). Si la fonction apparaît (mouvement, sensation), la structure est \emph{capable} de la mettre en jeu ;
\item \textbf{La lésion} : on détruit la structure (ablation expérimentale ou observation de lésions pathologiques). Si la fonction disparaît (paralysie, anesthésie), la structure est \emph{nécessaire} à cette fonction ;
\item \textbf{La conclusion} : si la stimulation provoque la fonction ET si la lésion l'abolit, alors la structure est le \emph{siège} de cette fonction ;
\item \textbf{La précision somatotopique} : en comparant point par point les effets, on établit la carte fonctionnelle de la structure.
\end{enumerate}
\end{methode}

\courssub{La somatotopie : l'homunculus moteur et l'homunculus sensitif}

En compilant les résultats de centaines de stimulations, Penfield a dressé une carte précise du corps le long des circonvolutions précentrale et postcentrale. Deux constats remarquables s'imposent.

\textbf{Premier constat : l'ordre des représentations.} Le corps est représenté de façon ordonnée mais \emph{inversée} : le pied et la jambe en haut de la circonvolution (face interne de l'hémisphère), puis le tronc, le bras, la main, et enfin le visage, les lèvres et la langue en bas (face latérale).

\textbf{Second constat : la disproportion des représentations.} La surface corticale dédiée à une région du corps n'est \emph{pas proportionnelle à la taille} de cette région, mais à la \emph{finesse du contrôle} moteur (aire motrice) ou à la \emph{densité des récepteurs} sensoriels (aire sensitive). La main et le visage, petits en volume, occupent d'immenses territoires corticaux, car ils réalisent des mouvements très fins (écriture, gestes de l'artisan, parole) et possèdent une sensibilité tactile très développée. Le tronc, volumineux, n'occupe qu'une bande étroite. Le corps ainsi déformé est appelé \cle{homunculus} (« petit homme ») moteur ou sensitif.

\begin{exemplebox}[Un exemple chiffré à retenir]
Chez l'Homme, la \textbf{main} et le \textbf{visage} (lèvres, langue) occupent à eux seuls \emph{plus de la moitié} de la surface de l'aire motrice, alors qu'ils représentent une fraction minime de la masse corporelle. C'est le prix cortical de la dextérité manuelle et de la parole, deux acquisitions majeures de l'espèce humaine.
\end{exemplebox}

\begin{popfigure}
\begin{tikzpicture}[scale=1.0]
% Bande représentant la circonvolution précentrale dépliée
\draw[very thick,poporange,fill=poporange!20] (0,0) rectangle (10,1.4);
% Zones proportionnelles (schématiques)
\draw[thick,poporange] (1.0,0) -- (1.0,1.4);
\draw[thick,poporange] (1.8,0) -- (1.8,1.4);
\draw[thick,poporange] (2.6,0) -- (2.6,1.4);
\draw[thick,poporange] (5.8,0) -- (5.8,1.4);
\draw[thick,poporange] (7.6,0) -- (7.6,1.4);
\draw[thick,poporange] (8.8,0) -- (8.8,1.4);
% Étiquettes dans les zones
\node[font=\footnotesize,align=center] at (0.5,0.7) {Pied\\Jambe};
\node[font=\footnotesize,align=center] at (1.4,0.7) {Tronc};
\node[font=\footnotesize,align=center] at (2.2,0.7) {Bras};
\node[font=\footnotesize\bfseries,align=center,text=popdark] at (4.2,0.7) {MAIN\\(très large)};
\node[font=\footnotesize,align=center] at (6.7,0.7) {Visage};
\node[font=\footnotesize\bfseries,align=center,text=popdark] at (8.2,0.7) {Lèvres};
\node[font=\footnotesize,align=center] at (9.4,0.7) {Langue};
% Orientation
\node[font=\footnotesize\itshape,text=popmuted,align=center] at (0.5,1.75) {Haut de la\\circonvolution};
\node[font=\footnotesize\itshape,text=popmuted,align=center] at (9.4,1.75) {Bas de la\\circonvolution};
% Petits personnages déformés : main géante
\draw[thick,popdark,fill=poppink!40] (4.2,-0.9) ellipse (0.85 and 0.55);
\node[font=\footnotesize,text=popdark] at (4.2,-1.7) {Main hypertrophiée};
\draw[thick,popdark,fill=poppink!40] (8.4,-0.9) ellipse (0.7 and 0.45);
\node[font=\footnotesize,text=popdark] at (8.4,-1.7) {Lèvres hypertrophiées};
\draw[thick,popdark,fill=popmuted!30] (1.4,-0.85) ellipse (0.35 and 0.5);
\node[font=\footnotesize,text=popmuted] at (1.4,-1.7) {Tronc réduit};
% Flèches de correspondance
\draw[->,popmuted,dashed] (4.2,0) -- (4.2,-0.3);
\draw[->,popmuted,dashed] (8.2,0) -- (8.4,-0.4);
\draw[->,popmuted,dashed] (1.4,0) -- (1.4,-0.3);
\end{tikzpicture}
\legende{Schéma de l'homunculus moteur : la circonvolution précentrale « dépliée » montre la carte somatotopique du corps. La largeur de chaque zone est proportionnelle à la surface corticale dédiée : la main et les lèvres (en gras) occupent l'essentiel du territoire, d'où un « petit homme » déformé aux mains et aux lèvres géantes.}
\end{popfigure}

\courssub{La latéralisation : un contrôle croisé du corps}

Toutes les expériences précédentes convergent vers un résultat capital : la stimulation de l'aire motrice \emph{gauche} provoque un mouvement du côté \emph{droit} du corps, et la lésion de l'aire motrice gauche paralyse le côté droit. De même, l'aire sensitive gauche reçoit les informations du côté droit. Ce \cle{contrôle croisé} (ou controlatéral) s'explique par le \emph{croisement des voies nerveuses} : les fibres motrices issues du cortex se croisent au niveau de la moelle allongée (bulbe rachidien) — c'est la \emph{décussation des pyramides} — et les fibres sensitives se croisent également avant d'atteindre le cortex. Nous détaillerons ce trajet dans la section suivante.

\begin{attention}[Erreur fréquente : le côté du déficit]
Beaucoup d'élèves écrivent qu'une lésion de l'aire motrice gauche paralyse le côté gauche. C'est \textbf{faux} ! Chaque hémisphère contrôle le côté \textbf{opposé} du corps : l'aire motrice gauche commande le côté \textbf{droit}. Retiens l'expression « \emph{contrôle croisé} » ou « \emph{controlatéral} ». Application clinique : un patient paralysé du bras droit après un AVC présente une lésion de l'hémisphère \emph{gauche}.
\end{attention}

\courssub{Les autres aires corticales (savoir admis)}

Le cortex cérébral comporte d'autres aires spécialisées, que nous signalons sans les développer : l'\cle{aire visuelle} (lobe occipital, à l'arrière du cerveau), l'\cle{aire auditive} (lobe temporal, sur le côté), et les \cle{aires du langage} (aire de Broca, commande motrice de la parole ; aire de Wernicke, compréhension du langage), situées chez la plupart des individus dans l'hémisphère gauche. Le reste du cortex est constitué d'\emph{aires associatives}, impliquées dans les fonctions supérieures (mémoire, raisonnement, décision).

\begin{savaistu}[Penfield et les patients éveillés]
Le neurochirurgien Wilder Penfield opérait ses patients épileptiques \emph{éveillés}, sous simple anesthésie locale — le tissu cérébral lui-même ne possède pas de récepteurs de douleur ! Pendant l'opération, il stimulait doucement le cortex et demandait au patient ce qu'il ressentait : « Je sens un fourmillement dans la main gauche », « Mon pied droit a bougé ». C'est ainsi, dialogue après dialogue, qu'il cartographia l'homunculus moteur et sensitif, encore utilisé aujourd'hui par tous les neurologues du monde, y compris dans les hôpitaux camerounais pour localiser la zone lésée lors d'un AVC.
\end{savaistu}

\begin{exoresolu}[Interpréter des résultats expérimentaux]
\textbf{Énoncé.} Chez un chien anesthésié, on stimule électriquement trois points A, B et C du cortex de l'hémisphère droit. On obtient les résultats suivants : stimulation de A (circonvolution précentrale, partie supérieure) : mouvement de la patte arrière gauche ; stimulation de B (circonvolution précentrale, partie inférieure) : mouvement des muscles de la face du côté gauche ; stimulation de C (circonvolution postcentrale) : aucun mouvement. Interprète ces résultats et dégage deux propriétés de l'organisation corticale.
\tcblower
\textbf{Solution détaillée.}
\begin{itemize}
\item \textbf{Stimulation de A et B.} Des mouvements sont obtenus uniquement à partir de la circonvolution précentrale : celle-ci constitue donc l'\emph{aire motrice}, siège de la commande du mouvement. À chaque point stimulé correspond un groupe musculaire précis et reproductible (patte en haut, face en bas) : l'aire motrice est organisée en \emph{carte somatotopique}.
\item \textbf{Stimulation de C.} La circonvolution postcentrale ne provoque aucun mouvement : elle n'est pas motrice (c'est l'aire sensitive, dont la stimulation provoque des sensations, non décelables chez un animal anesthésié qui ne peut pas les verbaliser).
\item \textbf{Côté des réponses.} L'hémisphère stimulé est le droit, et les réponses apparaissent du côté \emph{gauche} : le contrôle moteur est \emph{croisé} (controlatéral).
\end{itemize}
\textbf{Conclusion.} Deux propriétés sont dégagées : (1) le cortex est divisé en \emph{aires fonctionnelles spécialisées} ; (2) l'aire motrice présente une organisation \emph{somatotopique} et un contrôle \emph{croisé} du corps.
\end{exoresolu}

\begin{aretenir}[L'essentiel de la section]
\begin{itemize}
\item Le cortex cérébral n'est pas homogène : il est divisé en \cle{aires spécialisées} (motrice, sensitive, visuelle, auditive, du langage).
\item L'\cle{aire motrice} (circonvolution précentrale) commande les mouvements volontaires ; l'\cle{aire sensitive} (circonvolution postcentrale) reçoit les sensations somatiques.
\item La preuve repose sur deux approches complémentaires : la \emph{stimulation} (la fonction apparaît) et la \emph{lésion} (la fonction disparaît).
\item Chaque aire contient une carte ordonnée du corps : la \cle{somatotopie}, figurée par l'\cle{homunculus}. La surface corticale dédiée à une région est proportionnelle à la finesse du contrôle, pas à sa taille (main et lèvres : plus de la moitié de l'aire motrice).
\item Le contrôle est \cle{croisé} : chaque hémisphère commande et perçoit le côté \emph{opposé} du corps.
\end{itemize}
\end{aretenir}

Maintenant que nous savons \emph{où} naît la commande et \emph{où} aboutit l'information sensorielle, une question s'impose : quel \emph{chemin} l'influx nerveux parcourt-il réellement entre le cortex et les muscles ? C'est l'objet de la section suivante, consacrée au trajet de l'influx nerveux lors d'un mouvement volontaire.

\coursec{Trajet de l'influx nerveux lors d'un mouvement volontaire}

La section précédente nous a permis de localiser les \cle{aires motrices} et \cle{aires sensitives} du cortex cérébral : les stimulations électriques du cortex précentral déclenchent des mouvements du côté opposé du corps, et les lésions de ces zones entraînent des paralysies. Mais une question demeure, et elle est au c\oe ur de cette leçon : une fois l'ordre moteur élaboré dans le cortex, \textbf{par quel chemin l'influx nerveux parvient-il jusqu'au muscle} qui doit se contracter ? Pourquoi une blessure de la moelle épinière prive-t-elle un patient de tout mouvement volontaire alors que son cerveau est intact ? C'est à cette question de \emph{trajet} que nous répondons maintenant, étape par étape, comme on suit un message le long d'un réseau de communication.

\courssub{Le problème : du cerveau au muscle, un long voyage}

\begin{exemplebox}[Une situation d'observation]
Au Centre Hospitalier d'Essos à Yaoundé, un jeune motocycliste victime d'un accident présente une fracture vertébrale au niveau cervical. Son cerveau fonctionne parfaitement : il parle, comprend, \emph{veut} bouger ses jambes. Pourtant, aucun mouvement volontaire des membres inférieurs n'est possible. En revanche, un pincement de la plante du pied déclenche encore un léger retrait réflexe.
\end{exemplebox}

Cette observation clinique est riche d'enseignements. Elle montre que :
\begin{itemize}
\item la \cle{volonté} de bouger ne suffit pas : il faut une \textbf{voie de communication intacte} entre le cerveau et les muscles ;
\item cette voie passe par la \textbf{moelle épinière}, puisque sa section interrompt le mouvement volontaire ;
\item la moelle conserve des circuits locaux autonomes (les réflexes), distincts de la commande volontaire.
\end{itemize}

L'investigation scientifique consiste alors à suivre le message nerveux depuis son élaboration jusqu'à l'effecteur, en identifiant chaque relais et chaque structure impliquée. Nous distinguerons cinq étapes.

\courssub{Étape 1 : élaboration du programme moteur dans le cortex}

Tout mouvement volontaire — attraper un fruit, écrire, frapper un ballon — commence par une \textbf{décision}. Cette décision naît dans les \cle{aires d'association} du cortex (notamment le cortex préfrontal et les aires prémotrices), qui élaborent un \textbf{programme moteur} : quel muscle contracter, dans quel ordre, avec quelle force, vers quelle cible.

Ce programme est ensuite transmis à l'\cle{aire motrice primaire} (cortex précentral, circonvolution frontale ascendante), dont les grandes cellules pyramidales constituent les \textbf{neurones moteurs supérieurs}. Rappelons un résultat de la section précédente : le corps y est représenté de manière inversée et disproportionnée (homonculus moteur) — les zones de la main et du visage, très habiles, occupent une surface corticale considérable.

\begin{definition}[Mouvement volontaire]
Un \cle{mouvement volontaire} est un mouvement décidé et commandé par l'encéphale (cortex cérébral), transmis aux muscles squelettiques par l'intermédiaire de la moelle épinière, et contrôlé en permanence par les informations sensorielles de retour. Il se distingue du mouvement réflexe, involontaire et stéréotypé, dont le centre est médullaire.
\end{definition}

\courssub{Étape 2 : le cervelet et les noyaux gris centraux, chefs d'orchestre silencieux}

Le cortex moteur ne travaille jamais seul. Avant et pendant l'exécution du mouvement, une copie du programme moteur est envoyée au \cle{cervelet} et aux \cle{noyaux gris centraux} (structures profondes de l'encéphale : striatum, pallidum, etc.).

\begin{propriete}[Rôle du cervelet (admis)]
Le \cle{cervelet} ne déclenche \textbf{pas} le mouvement volontaire. Il \textbf{coordonne} et \textbf{corrige} le mouvement en temps réel : il compare en permanence le programme moteur envoyé par le cortex aux informations sensorielles revenant des muscles, des articulations et de l'organe de l'équilibre (vestibule de l'oreille interne), puis envoie des signaux correcteurs qui assurent la précision, la fluidité et l'équilibre. Les noyaux gris centraux participent au réglage du tonus musculaire et à la sélection des mouvements.
\end{propriete}

\begin{experience}[Preuve par la lésion : l'ataxie cérébelleuse]
\textbf{Observation clinique et expérimentale.} Chez un animal (chien, singe) dont le cervelet a été lésé ou partiellement ablaté, les mouvements volontaires restent \emph{possibles} mais deviennent \textbf{désordonnés} : démarche chancelante, gestes amples et imprécis qui dépassent la cible, tremblement à l'approche de l'objectif. Cet ensemble de signes constitue l'\cle{ataxie}. Chez l'Homme, on observe le même tableau après certaines hémorragies cérébelleuses.

\textbf{Interprétation.} Le mouvement n'est pas aboli : la \emph{commande} (cortex moteur) est intacte. Mais il est mal exécuté : la \emph{coordination} (cervelet) fait défaut.

\textbf{Conclusion.} Commande et coordination sont deux fonctions distinctes, assurées par deux structures différentes : le \textbf{cortex moteur commande}, le \textbf{cervelet coordonne}.
\end{experience}

\begin{attention}[Erreur classique à éviter absolument]
N'écris \textbf{jamais} que l'influx moteur volontaire « part du cervelet » ou « est élaboré par le cervelet ». La commande volontaire naît dans le \cle{cortex moteur} (aire motrice primaire). Le cervelet intervient sur une \textbf{boucle latérale} de contrôle et de correction. De même, les noyaux gris centraux ne commandent pas le mouvement : leur dysfonctionnement (maladie de Parkinson, étudiée à la section 4) perturbe le tonus et le déclenchement, sans détruire la volonté de bouger.
\end{attention}

\courssub{Étape 3 : la descente par les voies pyramidales et la décussation}

Une fois le programme validé, l'influx moteur quitte le cortex. Les axones des cellules pyramidales de l'aire motrice primaire s'engagent dans une longue voie descendante qui traverse la substance blanche cérébrale, le tronc cérébral (pédoncules cérébraux, pont, bulbe rachidien) et rejoint la moelle épinière.

\begin{definition}[Voie pyramidale]
La \cle{voie pyramidale} (ou \cle{faisceau cortico-spinal}) est l'ensemble des fibres nerveuses (axones myélinisés) reliant les neurones moteurs du cortex cérébral aux motoneurones de la moelle épinière. Elle traverse le tronc cérébral et, au niveau du \textbf{bulbe rachidien}, la majorité de ses fibres \textbf{croise la ligne médiane} : c'est la \cle{décussation des pyramides}. Chaque hémisphère cérébral commande donc les mouvements de la moitié \textbf{opposée} (controlatérale) du corps.
\end{definition}

\begin{experience}[Preuve expérimentale : section des voies pyramidales]
\textbf{Protocole.} Chez un animal de laboratoire, on sectionne chirurgicalement le faisceau pyramidal d'un seul côté, au niveau du bulbe, \emph{au-dessus} de la décussation.

\textbf{Observations.} L'animal présente une \textbf{paralysie volontaire des membres du côté opposé} à la section : il ne peut plus exécuter de mouvement délibéré de ce côté, alors que les muscles ne sont ni lésés ni déconnectés de la moelle.

\textbf{Interprétation.} La section interrompt la seule voie par laquelle la commande corticale atteint les motoneurones. Le croisement des résultats (lésion à gauche $\Rightarrow$ paralysie à droite) démontre la \textbf{décussation} : la voie change de côté au niveau du bulbe.

\textbf{Conclusion.} La voie pyramidale est la voie obligée de la motricité volontaire ; elle relie le cortex moteur d'un hémisphère aux motoneurones controlatéraux.
\end{experience}

Cette décussation explique un fait clinique fondamental que tu croiseras souvent dans les sujets de Baccalauréat : un accident vasculaire cérébral (AVC) touchant l'aire motrice \emph{droite} provoque une paralysie du côté \emph{gauche} du corps (hémiplégie controlatérale).

\begin{savaistu}[Une conduction à grande vitesse]
Les axones moteurs de la voie pyramidale et des nerfs spinaux sont de gros axones \textbf{myélinisés}. Grâce à la gaine de myéline et à la conduction saltatoire (l'influx « saute » d'un n\oe ud de Ranvier au suivant), la vitesse de conduction y atteint environ \SI{100}{m/s}, soit \SI{360}{km/h}. Ainsi, l'ordre de bouger le pied parcourt le mètre et demi qui sépare le cortex de la moelle lombaire en environ \SI{15}{ms} : la décision et le geste nous paraissent simultanés. À titre de comparaison, les fibres amyéliniques de la douleur diffuse conduisent à moins de \SI{1}{m/s}.
\end{savaistu}

\courssub{Étape 4 : le relais médullaire et la jonction neuromusculaire}

À l'extrémité de la voie pyramidale, l'influx moteur fait \textbf{relais} dans la moelle épinière. Sur une coupe transversale, la moelle présente une \textbf{substance grise} centrale en forme de papillon (corps cellulaires des neurones) entourée de \textbf{substance blanche} (faisceaux d'axones ascendants et descendants). Les prolongements ventraux de la substance grise sont les \cle{cornes ventrales} : c'est là que siègent les corps cellulaires des \cle{motoneurones} (neurones moteurs inférieurs).

\begin{definition}[Motoneurone]
Le \cle{motoneurone} est le neurone dont le corps cellulaire se trouve dans la corne ventrale de la moelle épinière et dont l'axone, empruntant la racine ventrale puis le \textbf{nerf spinal}, se termine sur une fibre musculaire squelettique au niveau de la \cle{jonction neuromusculaire} (plaque motrice). C'est la « voie finale commune » de tout mouvement : aucune commande, volontaire ou réflexe, ne peut contracter un muscle sans passer par lui.
\end{definition}

À la jonction neuromusculaire, l'arrivée de l'influx provoque la libération d'un neuromédiateur, l'\textbf{acétylcholine}, qui dépolarise la membrane de la fibre musculaire et déclenche sa \textbf{contraction}. Le muscle, \cle{effecteur} du mouvement, exécute alors l'ordre cortical.

\begin{popfigure}
\begin{center}
\begin{tikzpicture}[scale=1.05, every node/.style={font=\small}]
% Substance blanche (moelle)
\fill[popcream, draw=popline, thick] (0,0) ellipse (3.4 and 2.6);
% Substance grise (papillon)
\fill[popblueL, draw=popblue, thick]
  (0,1.05) .. controls (0.55,1.15) and (0.95,1.35) .. (1.15,1.7)
  .. controls (1.3,1.95) and (0.7,1.85) .. (0.5,1.35)
  .. controls (0.35,1.0) and (0.75,0.6) .. (1.35,0.15)
  .. controls (1.95,-0.3) and (1.9,-1.15) .. (1.35,-1.35)
  .. controls (0.85,-1.55) and (0.4,-0.9) .. (0.25,-0.5)
  .. controls (0.12,-0.2) and (-0.12,-0.2) .. (-0.25,-0.5)
  .. controls (-0.4,-0.9) and (-0.85,-1.55) .. (-1.35,-1.35)
  .. controls (-1.9,-1.15) and (-1.95,-0.3) .. (-1.35,0.15)
  .. controls (-0.75,0.6) and (-0.35,1.0) .. (-0.5,1.35)
  .. controls (-0.7,1.85) and (-1.3,1.95) .. (-1.15,1.7)
  .. controls (-0.95,1.35) and (-0.55,1.15) .. (0,1.05) -- cycle;
% Canal épendymaire
\fill[white, draw=popline] (0,0.35) circle (0.09);
% Motoneurone dans la corne ventrale droite
\fill[poporange] (1.05,-0.85) circle (0.16);
\draw[poporange, thick] (1.05,-0.85) -- (0.85,-0.55);
\draw[poporange, thick] (1.05,-0.85) -- (1.3,-0.6);
\draw[poporange, thick] (1.05,-0.85) -- (0.8,-1.05);
% Axone sortant par la racine ventrale
\draw[poporange, very thick, -{Stealth[length=3mm]}] (1.2,-0.95) .. controls (2.0,-1.4) and (2.6,-1.7) .. (3.5,-1.9) -- (5.2,-2.2);
% Racine dorsale (sensitif)
\draw[popgreen, very thick, -{Stealth[length=3mm]}] (5.2,2.1) -- (3.4,1.8) .. controls (2.6,1.6) and (2.0,1.5) .. (1.5,1.45);
\fill[popgreen] (3.4,1.8) circle (0.18);
% Flèche voie descendante vers motoneurone
\draw[popdark, very thick, -{Stealth[length=3mm]}] (0.2,2.6) -- (0.6,1.4) .. controls (0.8,0.6) and (0.95,-0.2) .. (1.0,-0.68);
% Labels
\node[popdark, font=\small\bfseries] at (0,2.95) {Substance blanche (voies ascendantes et descendantes)};
\node[popblue, font=\small\bfseries] at (-2.6,0.2) {Substance grise};
\node[poporange, font=\small\bfseries, align=center] at (5.6,-2.5) {Axone du motoneurone\\ (racine ventrale $\to$ nerf spinal)};
\node[popgreen, font=\small\bfseries, align=center] at (5.7,2.4) {Fibre sensitive\\ (racine dorsale, ganglion spinal)};
\node[popdark, font=\small\bfseries, align=center] at (-1.6,2.2) {Fibre de la voie\\ pyramidale descendante};
\node[poporange, font=\small\bfseries] at (0.4,-1.9) {Corne ventrale};
\node[font=\footnotesize, popmuted] at (0.55,0.35) {Canal};
\draw[-{Stealth[length=2mm]}, poporange] (0.75,-1.75) -- (1.0,-1.0);
\node[font=\footnotesize, popdark] at (1.05,-0.85) [above right=0.02cm] {\textbf{Motoneurone}};
\end{tikzpicture}
\end{center}
\legende{Coupe transversale schématique de la moelle épinière. La substance grise centrale (bleu) forme des cornes ventrales (motrices) et dorsales (sensitives). Le corps cellulaire du motoneurone (orange) siège dans la corne ventrale ; son axone sort par la racine ventrale et rejoint le muscle via le nerf spinal. Les fibres sensitives (vert) entrent par la racine dorsale, dont le renflement est le ganglion spinal. La fibre pyramidale descendante (noir) fait synapse avec le motoneurone.}
\end{popfigure}

\courssub{Étape 5 : le retour sensitif, ou comment le cerveau « sait » ce que fait le muscle}

Un mouvement volontaire n'est jamais un ordre lancé à l'aveugle. Dès que le muscle se contracte, des \textbf{récepteurs sensoriels} — fuseaux neuromusculaires dans les muscles, corpuscules dans la peau, récepteurs articulaires — sont stimulés. Ils génèrent des influx \textbf{afférents} (sensitives) qui remontent par les nerfs spinaux, la racine dorsale et les voies ascendantes de la substance blanche vers :
\begin{itemize}
\item l'\cle{aire sensitive} du cortex (cortex postcentral), qui informe la conscience sur la position et l'état des membres ;
\item le \cle{cervelet}, qui compare le mouvement réellement exécuté au programme prévu et émet les corrections nécessaires.
\end{itemize}

C'est un véritable \textbf{rétrocontrôle} : le mouvement est ajusté en permanence, comme un chauffeur de taxi à Douala corrige sans cesse sa trajectoire en fonction de ce qu'il voit et ressent. Sans ce retour sensitif (par exemple lors d'une atteinte des voies sensitives), les gestes deviennent imprécis même si la voie motrice est intacte.

\courssub{Schéma fonctionnel récapitulatif du circuit}

\begin{methode}[Construire un schéma fonctionnel du mouvement volontaire]
Pour obtenir tous les points au Baccalauréat, ton schéma fonctionnel doit respecter trois règles :
\begin{enumerate}
\item placer les structures \textbf{dans l'ordre anatomique réel} du trajet : récepteurs $\to$ aire sensitive $\to$ aires d'association $\to$ aire motrice $\to$ voie pyramidale (avec décussation) $\to$ corne ventrale $\to$ nerf spinal $\to$ muscle ;
\item relier les structures par des \textbf{flèches orientées} dans le sens de circulation de l'influx, jamais par de simples traits ;
\item faire figurer la \textbf{boucle de rétrocontrôle} : retour sensitif vers l'aire sensitive et le cervelet, et flèche correctrice du cervelet vers la voie motrice.
\end{enumerate}
Chaque flèche représente un influx nerveux ; chaque boîte, une structure anatomique précise.
\end{methode}

\begin{popfigure}
\begin{center}
\begin{tikzpicture}[
  font=\small,
  struct/.style={draw=popblue, fill=popblueL, thick, rounded corners=2pt, align=center, minimum height=0.85cm, inner sep=4pt},
  moteur/.style={draw=poporange, fill=poporangeL, thick, rounded corners=2pt, align=center, minimum height=0.85cm, inner sep=4pt},
  sens/.style={draw=popgreen, fill=popgreenL, thick, rounded corners=2pt, align=center, minimum height=0.85cm, inner sep=4pt},
  cereb/.style={draw=poppurple, fill=poppurpleL, thick, rounded corners=2pt, align=center, minimum height=0.85cm, inner sep=4pt},
  flux/.style={-{Stealth[length=3mm]}, very thick, popdark},
  fluxs/.style={-{Stealth[length=3mm]}, very thick, popgreen},
  fluxc/.style={-{Stealth[length=2.5mm]}, thick, poppurple}
]
% Chaîne principale
\node[struct, align=center] (assoc) at (0,0){Aires d'association\\ (programme moteur)};
\node[moteur, align=center] (cortex) at (4.1,0){Cortex moteur\\ (aire motrice primaire)};
\node[moteur, align=center] (pyr) at (8.2,0){Voie pyramidale\\ tronc cérébral\\ \textbf{décussation} (bulbe)};
\node[moteur, align=center] (moelle) at (12.3,0){Corne ventrale\\ de la moelle\\ (motoneurone)};
\node[moteur, align=center] (muscle) at (12.3,-2.6){Muscle effecteur\\ (contraction)};
\node[sens, align=center] (recept) at (8.2,-2.6){Récepteurs sensoriels\\ (fuseaux, peau, articulations)};
\node[sens, align=center] (sensitive) at (0,-2.6){Aire sensitive\\ (cortex postcentral)};
\node[cereb, align=center] (cerelet) at (4.1,-5.0){Cervelet\\ (comparaison, correction, équilibre)};
\node[cereb, align=center] (ngc) at (8.6,-5.0){Noyaux gris centraux\\ (tonus, sélection du mouvement)};
% Flèches motrices
\draw[flux] (assoc) -- (cortex);
\draw[flux] (cortex) -- (pyr);
\draw[flux] (pyr) -- (moelle);
\draw[flux] (moelle) -- node[right, font=\footnotesize]{nerf spinal, jonction neuromusculaire} (muscle);
% Retour sensitif
\draw[fluxs] (muscle) -- (recept);
\draw[fluxs] (recept.west) -- node[above, font=\footnotesize]{voies sensitives ascendantes} (sensitive.east);
\draw[fluxs] (sensitive) -- (assoc);
% Boucle cérébelleuse
\draw[fluxc] (cortex.south) .. controls (3.4,-1.6) and (3.4,-3.8) .. node[left, font=\footnotesize, align=center]{copie du\\ programme} (cerelet.north west);
\draw[fluxc] (recept.south) .. controls (7.6,-3.6) and (6.4,-4.4) .. node[below, font=\footnotesize, align=center]{informations\\ sensorielles} (cerelet.east);
\draw[fluxc] (cerelet.north east) .. controls (5.6,-3.6) and (5.6,-1.4) .. node[right, font=\footnotesize, align=center]{corrections} (cortex.south east);
\draw[fluxc] (ngc.north) .. controls (9.4,-3.4) and (10.6,-1.6) .. (pyr.south east);
\draw[fluxc] (cortex.south) ++(0.9,0) .. controls (6.2,-2.2) and (7.4,-3.6) .. (ngc.north west);
\end{tikzpicture}
\end{center}
\legende{Schéma fonctionnel du trajet de l'influx nerveux lors d'un mouvement volontaire. En orange, la chaîne motrice descendante : aires d'association $\to$ cortex moteur $\to$ voie pyramidale (avec décussation au bulbe) $\to$ motoneurone de la corne ventrale $\to$ muscle. En vert, le retour sensitif : récepteurs $\to$ aire sensitive. En violet, la boucle de contrôle : le cervelet et les noyaux gris centraux reçoivent une copie du programme et les informations sensorielles, puis corrigent le mouvement en cours.}
\end{popfigure}

\courssub{Exercice résolu et pièges de rédaction}

\begin{exoresolu}[Relier lésion et trouble]
Un patient présente, à la suite d'un accident, une \textbf{paralysie des mouvements volontaires du bras et de la jambe gauches}, sans trouble de la sensibilité ni de l'équilibre. Les réflexes locaux persistent. Un autre patient, lui, exécute tous ses mouvements mais de façon \textbf{désordonnée} : il tremble en attrapant un verre et chancelle en marchant.
\begin{enumerate}
\item Localise la lésion la plus probable chez le premier patient.
\item Quelle structure est atteinte chez le second ? Justifie.
\end{enumerate}
\tcblower
\textbf{1. Premier patient.} La paralysie est \emph{volontaire} (les réflexes, médullaires, persistent) et \emph{motrice pure} (sensibilité et équilibre normaux). La lésion touche donc la \textbf{commande motrice} ou sa voie de conduction. La paralysie étant \emph{gauche}, et la voie pyramidale se croisant au niveau du bulbe (\cle{décussation des pyramides}), la lésion siège du \textbf{côté droit} : soit dans l'\textbf{aire motrice primaire droite}, soit dans la \textbf{voie pyramidale droite au-dessus de la décussation} (par exemple au niveau du tronc cérébral). Une lésion médullaire est exclue car elle affecterait aussi la sensibilité.

\textbf{2. Second patient.} Les mouvements sont \emph{possibles} mais \emph{mal coordonnés} : c'est l'\cle{ataxie}. La commande corticale est donc intacte ; c'est la structure de \textbf{coordination et de correction} qui est lésée : le \textbf{cervelet}. Ce contraste illustre la distinction fondamentale entre commande (cortex moteur) et coordination (cervelet).
\end{exoresolu}

\begin{attention}[Pièges fréquents dans la description du trajet]
\begin{itemize}
\item \textbf{Oublier la décussation} : sans elle, impossible d'expliquer qu'une lésion cérébrale droite paralyse le côté gauche. Mentionne-la explicitement, avec son siège (bulbe rachidien).
\item \textbf{Confondre corne ventrale et corne dorsale} : la corne \emph{ventrale} contient les motoneurones (motricité) ; la corne \emph{dorsale} reçoit les fibres sensitives. Mnémotechnie : « V comme Ventral, V comme Vouloir bouger ».
\item \textbf{Faire passer l'influx moteur par la racine dorsale} : c'est faux ; l'axone du motoneurone sort par la \emph{racine ventrale}. La racine dorsale (avec son ganglion) est sensitive.
\item \textbf{Décrire un trajet sans relais} : le neurone cortical ne touche pas directement le muscle ; il fait \emph{relais} sur le motoneurone médullaire, qui est la « voie finale commune ».
\end{itemize}
\end{attention}

\begin{aretenir}[L'essentiel du trajet de l'influx volontaire]
\begin{itemize}
\item Le mouvement volontaire est \textbf{programmé} par les aires d'association et \textbf{commandé} par l'aire motrice primaire du cortex.
\item L'influx descend par la \textbf{voie pyramidale} (faisceau cortico-spinal), qui \textbf{croise la ligne médiane} au bulbe (décussation des pyramides) : chaque hémisphère commande le côté opposé du corps.
\item Le relais obligatoire est le \textbf{motoneurone} de la corne ventrale ; son axone emprunte le nerf spinal et stimule le muscle via la jonction neuromusculaire (acétylcholine).
\item Le \textbf{cervelet} et les \textbf{noyaux gris centraux} coordonnent et corrigent le mouvement sans le commander ; les \textbf{afférences sensitives} assurent le rétrocontrôle.
\item Preuves : section pyramidale $\Rightarrow$ paralysie volontaire controlatérale ; lésion cérébelleuse $\Rightarrow$ ataxie sans paralysie.
\end{itemize}
\end{aretenir}

Nous savons désormais par quel chemin précis l'ordre moteur atteint le muscle et comment ce trajet est contrôlé. Il reste à comprendre ce qui se passe lorsque ces structures tombent en panne : c'est l'objet des sections suivantes, consacrées à la maladie de Parkinson puis aux maladies de Huntington et d'Alzheimer.

\coursec{Dysfonctionnements : maladie de Parkinson}

Dans la section précédente, nous avons suivi le trajet de l'influx nerveux lors d'un mouvement volontaire : du cortex moteur vers la moelle épinière, puis vers le muscle squelettique, avec une participation des noyaux gris centraux et du cervelet au réglage du mouvement. Ce circuit, si bien organisé, peut cependant être altéré. Que se passe-t-il lorsque l'une de ces structures se dégrade ? La maladie de Parkinson, affection fréquente chez les personnes âgées — y compris au Cameroun, où le vieillissement de la population la rend de plus en plus visible dans nos hôpitaux de district — offre un exemple remarquable de lien entre une lésion encéphalique précise et un dysfonctionnement moteur caractéristique.

\courssub{Observation d'un cas clinique : les signes de la maladie}

\begin{exemplebox}[Un cas clinique au Centre Hospitalier d'Essos]
Monsieur N., 68 ans, retraité, consulte pour des tremblements de la main droite apparus il y a deux ans. Le médecin neurologue observe :
\begin{itemize}
\item un \cle{tremblement de repos} : la main tremble lorsqu'elle est posée sur la cuisse, mais le tremblement diminue lorsque le patient saisit un objet ;
\item une \cle{bradykinésie} (lenteur des mouvements) frisant l'\cle{akinésie} (difficulté à initier le mouvement) : le patient met plusieurs secondes à se lever de sa chaise ;
\item une \cle{rigidité musculaire} : les membres opposent une résistance à la mobilisation passive, dite « en tuyau de plomb » ;
\item une démarche à petits pas, le buste penché en avant, le visage figé et inexpressif (« masque parkinsonien »).
\end{itemize}
Le diagnostic clinique retenu est celui de la \cle{maladie de Parkinson}.
\end{exemplebox}

\textbf{Problème scientifique :} quelle structure du système nerveux est atteinte, et comment expliquer ce faisceau de symptômes ?

\begin{attention}[Un critère distinctif à ne jamais oublier]
Le tremblement parkinsonien survient \textbf{au repos} et \textbf{diminue pendant le mouvement volontaire}. C'est l'inverse du tremblement cérébelleux (tremblement d'action, qui apparaît pendant le mouvement). Au Baccalauréat, confondre ces deux types de tremblements est une erreur fréquente et lourde de conséquences : retiens l'association \emph{Parkinson = repos}, \emph{cervelet = action}.
\end{attention}

\courssub{La cause de la maladie : une dégénérescence des neurones dopaminergiques}

\begin{definition}[Maladie de Parkinson]
La \cle{maladie de Parkinson} est une affection dégénérative chronique du système nerveux central, due à la destruction progressive des neurones à \cle{dopamine} de la \cle{substance noire} (l'un des noyaux gris centraux, situé dans le mésencéphale). Elle se traduit par une triade de symptômes : \textbf{tremblement de repos}, \textbf{rigidité musculaire} et \textbf{akinésie} (ou bradykinésie).
\end{definition}

Comment cette cause a-t-elle été établie ? Par la démarche \cle{anatomopathologique}, c'est-à-dire l'examen des organes après le décès.

\begin{experience}[Preuves anatomo-pathologiques et biochimiques]
\begin{enumerate}
\item \textbf{Observation macroscopique à l'autopsie.} Chez un sujet sain, la substance noire du mésencéphale apparaît sombre, car ses neurones contiennent un pigment, la neuromélanine. Chez un patient parkinsonien décédé, cette zone apparaît \textbf{dépigmentée} : les neurones pigmentés ont disparu.
\item \textbf{Observation microscopique.} Les coupes histologiques confirment la perte massive des neurones dopaminergiques et la présence d'amas protéiques anormaux (corps de Lewy, constitués d'$\alpha$-synucléine) dans les neurones survivants.
\item \textbf{Dosages biochimiques.} Le dosage de la dopamine dans les noyaux gris centraux (notamment le striatum, cible des neurones de la substance noire) montre une \textbf{chute drastique} du taux de dopamine, la diminution dépassant \SI{80}{\%} (taux résiduel inférieur à \SI{20}{\%} de la normale) chez les patients symptomatiques.
\end{enumerate}
\textbf{Interprétation :} la destruction des neurones dopaminergiques de la substance noire entraîne un déficit en dopamine au niveau du striatum, responsable des troubles moteurs.
\end{experience}

\begin{exemplebox}[Un seuil chiffré remarquable]
Les symptômes moteurs de la maladie de Parkinson n'apparaissent que lorsque environ \textbf{70 à 80 \%} des neurones dopaminergiques de la substance noire sont déjà détruits. La maladie progresse donc silencieusement pendant des années avant de se déclarer : c'est la \emph{phase préclinique}. Cela explique l'importance du dépistage précoce des signes discrets (perte d'odorat, troubles du sommeil paradoxal, constipation chronique).
\end{exemplebox}

\begin{popfigure}
\begin{center}
\begin{tikzpicture}[scale=1.0, font=\small]
% Contour de l'encéphale (coupe sagittale schématique)
\draw[thick, popdark, fill=popcream] (0,0) .. controls (0.2,2.2) and (2.2,3.4) .. (4.6,3.0)
  .. controls (6.4,2.7) and (7.2,1.8) .. (7.0,0.9)
  .. controls (6.8,0.2) and (6.0,-0.1) .. (5.2,0.0)
  .. controls (4.6,-0.5) and (3.8,-0.7) .. (3.2,-0.4)
  .. controls (2.4,-0.9) and (1.2,-0.8) .. (0.6,-0.3)
  .. controls (0.1,-0.15) .. (0,0);
% Corps calleux
\draw[thick, popblue, fill=popblueL] (2.4,1.7) .. controls (3.4,2.3) and (4.8,2.2) .. (5.4,1.6)
  .. controls (4.8,1.9) and (3.4,2.0) .. (2.6,1.45) -- cycle;
% Noyaux gris centraux (striatum)
\draw[thick, poporange, fill=poporangeL] (3.3,0.9) ellipse (0.85 and 0.55);
% Thalamus
\draw[thick, poppurple, fill=poppurpleL] (4.6,0.7) ellipse (0.6 and 0.45);
% Tronc cérébral
\draw[thick, popdark, fill=popgreenL] (3.6,-0.4) .. controls (3.7,-1.3) and (4.1,-1.9) .. (4.0,-2.6)
  -- (4.7,-2.6) .. controls (4.8,-1.8) and (5.0,-1.0) .. (5.1,-0.3) -- cycle;
% Substance noire
\draw[thick, poppink, fill=poppinkL] (4.15,-0.75) ellipse (0.5 and 0.22);
% Cervelet
\draw[thick, popgold, fill=popgoldL] (1.6,-0.75) circle (0.85);
\draw[popgold!70!black] (1.0,-0.5) .. controls (1.6,-0.9) .. (2.2,-0.5);
\draw[popgold!70!black] (1.0,-1.0) .. controls (1.6,-1.3) .. (2.2,-1.0);
% Légendes fléchées
\draw[->, poporange, thick] (3.3,1.55) -- (3.3,2.6) node[above, poporange!80!black]{\textbf{1. Noyaux gris centraux (striatum)}};
\draw[->, poppink, thick] (4.65,-0.75) -- (6.3,-0.75) node[right, poppink!80!black, align=left]{\textbf{2. Substance noire}\\ (neurones à dopamine)};
\draw[->, popgold!80!black, thick] (1.0,-1.35) -- (0.6,-2.2) node[below, popgold!80!black]{\textbf{3. Cervelet}};
\draw[->, popgreen!60!black, thick] (4.35,-2.6) -- (4.35,-3.1) node[below, popgreen!50!black]{\textbf{4. Tronc cérébral}};
\draw[->, popblue, thick] (5.4,1.9) -- (6.3,2.6) node[right, popblue!80!black]{\textbf{5. Cortex (aire motrice)}};
% Projection dopaminergique
\draw[->, poppink!80!black, very thick, dashed] (4.15,-0.55) .. controls (3.9,0.1) .. (3.5,0.55);
\node[poppink!80!black, align=center] at (2.6,-0.15) {\footnotesize voie\\ \footnotesize dopaminergique};
\end{tikzpicture}
\end{center}
\legende{Coupe sagittale schématique de l'encéphale humain. \textbf{1} : noyaux gris centraux (dont le striatum), impliqués dans le déclenchement et le réglage des mouvements volontaires. \textbf{2} : substance noire du mésencéphale, siège des neurones dopaminergiques qui projettent vers le striatum (flèche pointillée) ; c'est leur dégénérescence qui cause la maladie de Parkinson. \textbf{3} : cervelet (coordination). \textbf{4} : tronc cérébral. \textbf{5} : cortex cérébral, où siège l'aire motrice.}
\end{popfigure}

\courssub{Lien structure-fonction : pourquoi le déficit en dopamine bloque-t-il le mouvement ?}

Rappelons un acquis de la section précédente : les \cle{noyaux gris centraux} (ou ganglions de la base) participent au \textbf{déclenchement} et au \textbf{réglage} des mouvements volontaires. Ils reçoivent des informations du cortex moteur et lui renvoient, via le thalamus, un signal qui « autorise » et module l'exécution du mouvement.

\begin{propriete}[Rôle de la dopamine — admise]
La \cle{dopamine} est un neurotransmetteur indispensable au réglage de la motricité par les noyaux gris centraux : libérée par les neurones de la substance noire au niveau du striatum, elle facilite l'initiation du mouvement volontaire et le réglage du tonus musculaire. (Propriété admise ; sa justification expérimentale complète dépasse le niveau exigible au Baccalauréat.)
\end{propriete}

Le raisonnement causal s'enchaîne alors logiquement :
\begin{enumerate}
\item les neurones dopaminergiques de la substance noire dégénèrent ;
\item le taux de dopamine dans les noyaux gris centraux (striatum) chute ;
\item le « signal d'autorisation » adressé au cortex moteur via le thalamus devient insuffisant ;
\item le patient peine à \textbf{initier} le mouvement (akinésie), ses mouvements sont \textbf{lents} (bradykinésie), le réglage du tonus est perturbé (\textbf{rigidité}), et l'équilibre des circuits moteurs est rompu, laissant apparaître un \textbf{tremblement de repos}.
\end{enumerate}

\begin{popfigure}
\begin{center}
\begin{tikzpicture}[scale=1.0]
\begin{axis}[
  width=0.85\linewidth, height=6.2cm,
  xlabel={Évolution de la maladie (années)},
  ylabel={Taux de dopamine striatale (\% de la normale)},
  xmin=0, xmax=20, ymin=0, ymax=110,
  xtick={0,5,10,15,20},
  ytick={0,20,40,60,80,100},
  grid=both, grid style={popline!40},
  legend style={at={(0.03,0.35)}, anchor=west, font=\small},
]
% Sujet sain
\addplot[popcurve, popblue, domain=0:20, samples=60]{100 - 0.8*x};
\addlegendentry{Sujet sain (vieillissement normal)}
% Sujet parkinsonien
\addplot[popcurve, poppink!80!black, domain=0:20, samples=100]{100/(1+exp(0.55*(x-11)))};
\addlegendentry{Sujet parkinsonien}
% Seuil symptomatique
\addplot[dashed, poporange, thick, domain=0:20, samples=2]{25};
\node[poporange!85!black, font=\footnotesize, anchor=south west] at (axis cs:0.3,27) {seuil d'apparition des symptômes ($\approx$ 70--80 \% de perte)};
% Zone symptomatique
\draw[->, poporange!85!black] (axis cs:13.5,12) -- (axis cs:16.5,12);
\node[poporange!85!black, font=\footnotesize] at (axis cs:15,7) {phase clinique};
\end{axis}
\end{tikzpicture}
\end{center}
\legende{Évolution du taux de dopamine dans le striatum en fonction de l'avancement de la maladie. Chez le sujet sain, la baisse liée au vieillissement est lente et reste sans conséquence clinique. Chez le sujet parkinsonien, la chute est rapide ; les symptômes moteurs n'apparaissent que lorsque la perte de neurones dopaminergiques atteint environ 70 à 80 \% (taux résiduel de dopamine inférieur à environ 20--30 \% de la normale).}
\end{popfigure}

\courssub{La preuve thérapeutique : l'effet de la L-DOPA}

Une preuve élégante du rôle de la dopamine vient de la \textbf{pharmacologie} : si les symptômes sont dus à un déficit en dopamine, alors restaurer ce déficit doit les atténuer.

\begin{experience}[Administration de L-DOPA à des patients parkinsoniens]
\textbf{Protocole :} on administre à des patients parkinsoniens de la \cle{L-DOPA} (lévodopa), précurseur naturel de la dopamine, par voie orale. On mesure l'évolution des scores moteurs (vitesse d'exécution de gestes standardisés, amplitude du tremblement, rigidité).

\textbf{Observations :} dans les heures qui suivent la prise, la rigidité et la bradykinésie diminuent nettement ; les patients retrouvent une capacité d'initiation du mouvement presque normale. L'arrêt du traitement fait réapparaître les symptômes.

\textbf{Interprétation :} la L-DOPA, transformée en dopamine dans l'encéphale, compense le déficit dopaminergique. L'amélioration réversible des symptômes confirme que ceux-ci résultent bien du manque de dopamine, et non d'une lésion du cortex moteur ou des voies pyramidales (qui, elles, causeraient une paralysie).
\end{experience}

\begin{savaistu}[Pourquoi la L-DOPA et pas la dopamine elle-même ?]
Utilisée depuis les années 1960, la \cle{L-DOPA} a révolutionné le destin des patients parkinsoniens. Mais pourquoi ne pas injecter directement de la dopamine ? Parce que la dopamine ne \textbf{traverse pas la barrière hémato-encéphalique}, ce filtre protecteur qui sépare le sang du tissu nerveux. La L-DOPA, elle, la franchit grâce à un transporteur spécifique, puis est transformée en dopamine sur place par une enzyme (la DOPA-décarboxylase). C'est un magnifique exemple de « cheval de Troie » pharmacologique ! Cette découverte a valu à Arvid Carlsson le prix Nobel de physiologie ou médecine en 2000.
\end{savaistu}

\courssub{Limitation du dysfonctionnement : les stratégies thérapeutiques}

Le titre de cette leçon parle de \emph{limitation} des dysfonctionnements : il faut être honnête, la médecine actuelle ne \textbf{guérit pas} la maladie de Parkinson, mais elle en limite efficacement les manifestations.

\begin{methode}[Les trois piliers de la prise en charge]
\begin{enumerate}
\item \textbf{Traitement médicamenteux.} La \cle{L-DOPA} (associée à un inhibiteur de la décarboxylase périphérique, comme la carbidopa ou la bensérazide, pour limiter les effets indésirables périphériques) reste le traitement de référence. On y ajoute des \cle{agonistes dopaminergiques}, molécules qui imitent l'action de la dopamine sur ses récepteurs.
\item \textbf{Kinésithérapie et rééducation.} Des exercices réguliers de marche, d'équilibre et d'élocution ralentissent la perte d'autonomie et préviennent les chutes, première cause de complications.
\item \textbf{Stimulation cérébrale profonde.} Dans les cas sévères ou lorsque les médicaments deviennent mal tolérés, on implante chirurgicalement des électrodes dans certains noyaux gris centraux (noyau sous-thalamique le plus souvent). Un stimulateur envoie des impulsions électriques qui rétablissent un fonctionnement plus régulier des circuits moteurs.
\end{enumerate}
\end{methode}

\begin{attention}[Traitement symptomatique, non curatif]
La L-DOPA \textbf{compense} le déficit en dopamine mais ne \textbf{stoppe pas} la dégénérescence des neurones : la maladie continue de progresser sous traitement. De plus, après plusieurs années, la L-DOPA perd en efficacité et provoque des mouvements anormaux involontaires (dyskinésies). Ne rédige jamais au Bac que « la L-DOPA guérit la maladie de Parkinson » : c'est faux. Écris : « la L-DOPA est un traitement \emph{symptomatique} qui améliore durablement la qualité de vie ».
\end{attention}

\courssub{Facteurs favorisants et prévention}

\begin{itemize}
\item \textbf{Le vieillissement} est le principal facteur de risque : la maladie touche environ 1 \% des personnes de plus de 65 ans, et sa fréquence augmente avec l'âge.
\item \textbf{Des facteurs génétiques} : certaines formes familiales rares sont liées à des mutations identifiées (gènes codant l'$\alpha$-synucléine, la parkine).
\item \textbf{Des facteurs environnementaux} : une exposition prolongée à certains pesticides augmente le risque — donnée à méditer dans les zones de cultures intensives de cacao ou de coton, où la protection des applicateurs de produits phytosanitaires est un enjeu de santé publique.
\end{itemize}

Il n'existe à ce jour \textbf{aucune prévention certaine} de la maladie de Parkinson. En revanche, le \cle{dépistage précoce} des signes avant-coureurs et la consultation neurologique dès l'apparition d'un tremblement de repos permettent de mettre en place le traitement au moment optimal et de préserver l'autonomie le plus longtemps possible. L'activité physique régulière et la limitation de l'exposition aux toxiques sont des mesures de bon sens recommandées à tous.

\begin{exoresolu}[Expliquer les symptômes d'un patient]
\textbf{Énoncé.} Un patient de 70 ans présente un tremblement des mains qui s'atténue lorsqu'il boit son verre d'eau, une grande lenteur pour se lever et une rigidité des membres. À l'autopsie d'un patient présentant les mêmes signes, on observe une dépigmentation de la substance noire et une chute de 85 \% du taux de dopamine striatale. (1) Nommer la maladie. (2) Expliquer le mécanisme reliant la lésion observée aux symptômes moteurs. (3) Proposer un traitement et préciser sa nature (curatif ou symptomatique).
\tcblower
\textbf{Solution détaillée.}

(1) Il s'agit de la \textbf{maladie de Parkinson} : la triade tremblement de repos (qui diminue pendant le mouvement volontaire, comme le précise l'énoncé), akinésie/bradykinésie et rigidité est caractéristique.

(2) La dépigmentation de la substance noire traduit la \textbf{dégénérescence des neurones dopaminergiques} qui s'y trouvent. Ces neurones libèrent normalement la dopamine au niveau du striatum (noyaux gris centraux), neurotransmetteur indispensable au déclenchement et au réglage des mouvements volontaires. La chute de 85 \% du taux de dopamine — supérieure au seuil de 70 à 80 \% de perte neuronale — déséquilibre le fonctionnement des noyaux gris centraux : le signal facilitant l'initiation du mouvement devient insuffisant (d'où l'akinésie et la lenteur), le réglage du tonus est perturbé (d'où la rigidité) et les circuits moteurs déséquilibrés génèrent un tremblement au repos.

(3) On propose un traitement par \textbf{L-DOPA}, précurseur de la dopamine capable de traverser la barrière hémato-encéphalique, éventuellement associé à des agonistes dopaminergiques et à une kinésithérapie. Il s'agit d'un traitement \textbf{symptomatique} : il compense le déficit en dopamine et améliore les symptômes, mais ne stoppe pas la dégénérescence neuronale ; il n'est donc pas curatif.
\end{exoresolu}

\begin{aretenir}[L'essentiel sur la maladie de Parkinson]
\begin{itemize}
\item Maladie dégénérative due à la destruction des \textbf{neurones dopaminergiques de la substance noire} (noyaux gris centraux, mésencéphale).
\item Triade symptomatique : \textbf{tremblement de repos}, \textbf{rigidité}, \textbf{akinésie}.
\item Preuves : dépigmentation de la substance noire à l'autopsie, chute du taux de dopamine striatale ($> 80 \%$), amélioration réversible sous L-DOPA.
\item Les symptômes apparaissent après destruction de \textbf{70 à 80 \%} des neurones (phase préclinique longue).
\item Traitement \textbf{symptomatique} : L-DOPA (+ inhibiteur décarboxylase périphérique), agonistes dopaminergiques, kinésithérapie, stimulation cérébrale profonde (noyau sous-thalamique).
\item Facteurs de risque : vieillissement, génétique (rares), environnement (pesticides) ; pas de prévention certaine, dépistage précoce essentiel.
\end{itemize}
\end{aretenir}

La maladie de Parkinson illustre parfaitement comment l'atteinte d'une structure encéphalique précise perturbe le contrôle de la motricité. Mais d'autres affections touchent le cerveau avec des conséquences différentes : c'est le cas des maladies de Huntington et d'Alzheimer, que nous étudierons dans la section suivante.

\coursec{Dysfonctionnements : maladies de Huntington et d'Alzheimer}

La maladie de Parkinson, étudiée précédemment, nous a montré comment la dégénérescence d'une structure bien précise — la substance noire, productrice de dopamine — entraîne des troubles moteurs caractéristiques : tremblements de repos, rigidité, akinésie. Mais le cerveau peut être atteint par d'autres maladies dégénératives, touchant d'autres structures et provoquant d'autres symptômes. Comment expliquer que certains patients présentent des mouvements involontaires incontrôlables, tandis que d'autres perdent progressivement la mémoire et l'autonomie intellectuelle ? Deux maladies vont nous permettre de répondre : la \cle{maladie de Huntington} et la \cle{maladie d'Alzheimer}.

\courssub{La maladie de Huntington : une chorée héréditaire}

\textbf{Observation clinique.}\par

Dans certaines familles, on observe que plusieurs membres, à l'âge adulte, développent les mêmes troubles : des gestes brusques, désordonnés, qu'ils ne peuvent contrôler, comme s'ils « dansaient » malgré eux. Peu à peu s'y ajoutent des changements de caractère, puis une perte des capacités intellectuelles. Cette maladie, décrite en 1872 par le médecin américain George Huntington, porte son nom.

\begin{definition}[Maladie de Huntington]
La \cle{maladie de Huntington} est une \cle{maladie héréditaire} à transmission \cle{autosomique dominante}, caractérisée par la \cle{dégénérescence} des neurones des \cle{noyaux gris centraux} (en particulier le \cle{striatum}, ensemble noyau caudé + putamen), puis du cortex cérébral. Elle se manifeste vers 35--50 ans par des \cle{mouvements involontaires} brusques et désordonnés appelés \cle{chorée}, auxquels s'ajoutent des troubles du comportement et un déclin intellectuel progressif.
\end{definition}

\textbf{Les symptômes.}\par

On distingue trois groupes de manifestations, qui évoluent ensemble au fil des années :
\begin{itemize}
\item \textbf{Troubles moteurs} : mouvements involontaires, brusques et désordonnés des membres, du visage et du tronc (la \cle{chorée}, du grec \emph{khoreia}, « danse ») ; plus tard, troubles de l'équilibre, de la marche et de la parole (dysarthrie) ;
\item \textbf{Troubles psychiques et du comportement} : irritabilité, dépression, anxiété, parfois troubles graves du jugement et de l'initiative ;
\item \textbf{Déclin intellectuel} : difficultés de concentration, pertes de mémoire, ralentissement psychomoteur, aboutissant à une démence.
\end{itemize}

Contrairement à la maladie de Parkinson, où les mouvements sont \emph{raréfiés} (akinésie), la maladie de Huntington se caractérise par un \emph{excès} de mouvements incontrôlés. Cette différence s'explique par la structure atteinte : le \cle{striatum}, qui participe au \emph{filtrage} et à l'\emph{inhibition} des mouvements parasites via la voie indirecte des noyaux de la base. Quand ses neurones GABAergiques dégénèrent, le thalamus n'est plus correctement inhibé : il surexcite le cortex moteur, et des ordres moteurs « sauvages » s'expriment.

\textbf{La cause : une mutation à expansion de tripletés.}\par

La maladie de Huntington est due à une \cle{mutation} d'un gène (\emph{HTT}) situé sur le bras court du chromosome 4 (4p16,3), codant une protéine appelée \cle{huntingtine}. La mutation consiste en une \textbf{expansion anormale d'un triplet CAG} (codant la glutamine) dans l'exon 1 : au-delà de 36--39 répétitions, l'allèle est muté (allèle \emph{H}). La protéine huntingtine mutée (polyglutamine allongée) est toxique : elle s'agrège dans le noyau et le cytoplasme des neurones du striatum, perturbant la transcription, le transport axonal et la fonction mitochondriale, ce qui provoque leur mort progressive par apoptose. La dégénérescence s'étend ensuite au cortex cérébral. L'allèle muté est \cle{dominant} : il suffit d'un seul exemplaire pour développer la maladie (pénétrance complète à l'âge adulte).

\textbf{La preuve génétique : l'étude des arbres généalogiques.}\par

Comment démontre-t-on le caractère héréditaire et dominant de la maladie ? Par l'\cle{analyse de généalogies} (pedigrees) dans les familles atteintes. Observons l'arbre ci-dessous.

\begin{popfigure}
\begin{center}
\begin{tikzpicture}[
  male/.style={rectangle, draw=popblue, thick, minimum size=0.55cm, inner sep=2pt},
  female/.style={circle, draw=poppink, thick, minimum size=0.55cm, inner sep=2pt},
  maleA/.style={rectangle, draw=popblue, fill=popblue, thick, minimum size=0.55cm, inner sep=2pt},
  femaleA/.style={circle, draw=poppink, fill=poppink, thick, minimum size=0.55cm, inner sep=2pt},
  every node/.style={font=\small},
  lbl/.style={font=\footnotesize\itshape, text=popmuted}
]
% Génération I
\node[maleA, label=above:{$I_1$}] (I1) at (0,0) {};
\node[female, label=above:{$I_2$}] (I2) at (1.6,0) {};
\draw (I1) -- (I2);
\coordinate (mI) at (0.8,0);
% Génération II
\draw (mI) -- (0.8,-0.6) -- (-1.4,-0.6);
\draw (-1.4,-0.6) -- (-1.4,-1.4);
\draw (0.8,-0.6) -- (3.0,-0.6);
\draw (3.0,-0.6) -- (3.0,-1.4);
\node[femaleA, label=left:{$II_1$}] (II1) at (-1.4,-1.4) {};
\node[male, label=left:{$II_2$}] (II2) at (-0.2,-1.4) {};
\draw (II1) -- (II2);
\node[maleA, label=right:{$II_3$}] (II3) at (3.0,-1.4) {};
% Génération III (enfants de II1 x II2)
\coordinate (mII) at (-0.8,-1.4);
\draw (mII) -- (-0.8,-2.0) -- (-2.6,-2.0);
\draw (-2.6,-2.0) -- (-2.6,-2.8);
\draw (-0.8,-2.0) -- (1.0,-2.0);
\draw (1.0,-2.0) -- (1.0,-2.8);
\node[maleA, label=below:{$III_1$}] at (-2.6,-2.8) {};
\node[female, label=below:{$III_2$}] at (-1.4,-2.8) {};
\node[male, label=below:{$III_3$}] at (-0.2,-2.8) {};
\node[femaleA, label=below:{$III_4$}] at (1.0,-2.8) {};
% Légende
\node[maleA] at (5.0,-0.2) {}; \node[right, font=\footnotesize] at (5.4,-0.2) {homme atteint};
\node[femaleA] at (5.0,-0.9) {}; \node[right, font=\footnotesize] at (5.4,-0.9) {femme atteinte};
\node[male] at (5.0,-1.6) {}; \node[right, font=\footnotesize] at (5.4,-1.6) {homme sain};
\node[female] at (5.0,-2.3) {}; \node[right, font=\footnotesize] at (5.4,-2.3) {femme saine};
\end{tikzpicture}
\end{center}
\legende{Arbre généalogique d'une famille atteinte de la maladie de Huntington. Carré = homme ; rond = femme ; symbole rempli = individu atteint ; trait horizontal = union ; trait vertical = descendance. On note que la maladie apparaît à \emph{chaque génération} et touche environ la moitié des enfants d'un parent malade : signature d'une transmission autosomique dominante.}
\end{popfigure}

L'interprétation de cet arbre repose sur deux indices décisifs :
\begin{enumerate}
\item la maladie frappe \textbf{toutes les générations} (pas de « saut » de génération, contrairement aux maladies récessives) ;
\item un parent atteint, presque toujours \cle{hétérozygote} $H/h$ (l'allèle muté $H$ étant rare, les homozygotes $H/H$ sont exceptionnels), transmet l'allèle muté à environ \textbf{la moitié de ses enfants}, garçons et filles indifféremment (le gène est porté par un \cle{autosome}, pas par un chromosome sexuel).
\end{enumerate}

\begin{savaistu}[Un risque de 50 \% et le test prédictif]
Un enfant dont l'un des parents est atteint de la maladie de Huntington a \textbf{50 \% de risque} de porter l'allèle muté, et donc de développer la maladie. Depuis l'identification du gène \emph{HTT} en 1993, un \cle{test génétique prédictif} existe : il permet à une personne à risque, avant l'apparition des symptômes, de savoir si elle porte l'expansion CAG. Ce test pose de lourdes questions éthiques (annonce d'une maladie incurable, absence de traitement curatif, impact sur la descendance, assurance) et n'est pratiqué qu'avec un accompagnement psychologique rigoureux, dans des centres spécialisés de génétique médicale (ex. : Centre de Génétique Humaine de la Faculté de Médecine de Yaoundé).
\end{savaistu}

\courssub{La maladie d'Alzheimer : une démence dégénérative}

\textbf{Observation clinique.}\par

Une grand-mère, autrefois vive et autonome, oublie d'abord les événements récents : elle repose trois fois la même question, égare ses affaires. Puis elle se perd dans son propre quartier de Yaoundé, ne reconnaît plus ses proches, peine à trouver ses mots (aphasie). Aux stades avancés, des troubles moteurs apparaissent (marche lente, rigidité) et elle devient totalement dépendante pour les actes de la vie quotidienne. C'est le tableau typique de la maladie d'Alzheimer, décrite en 1906 par le neurologue allemand Alois Alzheimer.

\begin{definition}[Maladie d'Alzheimer]
La \cle{maladie d'Alzheimer} est une \cle{démence dégénérative} liée à la \cle{mort progressive des neurones} du \cle{cortex cérébral} (aire d'association, cortex temporal, pariétal, frontal) et de l'\cle{hippocampe} (structure essentielle à l'encodage de la mémoire épisodique). Elle se traduit par des \cle{troubles de la mémoire} (d'abord des faits récents, mémoire épisodique), une \cle{désorientation} dans le temps et l'espace, des troubles du langage (aphasie), de la reconnaissance (agnosie), du geste (apraxie) et du jugement, puis, aux stades avancés, des troubles moteurs.
\end{definition}

\textbf{La cause : une dégénérescence corticale avec lésions histologiques caractéristiques.}\par

À l'autopsie (ou à l'imagerie cérébrale du vivant : TEP-amyloïde, IRM volumétrique), le cerveau d'un malade d'Alzheimer présente deux lésions microscopiques typiques, critères diagnostiques histologiques :
\begin{itemize}
\item des \cle{plaques amyloïdes} (ou plaques séniles) : amas \emph{extracellulaires} d'un peptide anormal, le peptide $\beta$-amyloïde (A$\beta$42), issu de la clivage anormal de la protéine précurseur APP (Amyloid Precursor Protein). Ces dépôts perturbent la synapse, déclenchent l'inflammation neurogliale et la toxicité neuronale ;
\item des \cle{dégénérescences neurofibrillaires} : enchevêtrements \emph{intracellulaires} de filaments hélicoïdaux appariés (protéine \emph{tau} hyperphosphorylée) à l'intérieur des neurones. La protéine tau normale stabilise les microtubules ; hyperphosphorylée, elle s'en détache, s'agrège et bloque le transport axonal (mitochondries, vésicules), conduisant à la mort neuronale.
\end{itemize}
La mort massive des neurones (perte synaptique précoce, puis perte cellulaire) entraîne une \cle{atrophie cérébrale} visible à l'imagerie (scanner, IRM) : le cortex rétrécit, les sillons (circonvolutions) s'élargissent et les ventricules (cavités internes liquides) se dilatent (hydrocéphalie ex vacuo). Par ailleurs, on observe un \textbf{déficit cholinergique} précoce (dégénérescence du noyau basal de Meynert), base du traitement symptomatique par inhibiteurs de l'acétylcholinestérase.

\begin{popfigure}
\begin{center}
\begin{tikzpicture}[scale=1]
% Cerveau sain
\begin{scope}
\fill[popgreenL] (0,0) ellipse (2.2 and 1.6);
\draw[popgreen, thick] (0,0) ellipse (2.2 and 1.6);
% sillons fins
\draw[popgreen, thick] (-1.6,0.5) .. controls (-1.0,0.9) and (-0.4,0.2) .. (0.2,0.7);
\draw[popgreen, thick] (-1.4,-0.5) .. controls (-0.8,-0.1) and (0.0,-0.8) .. (0.6,-0.3);
\draw[popgreen, thick] (0.4,1.0) .. controls (0.9,0.5) and (1.3,0.9) .. (1.7,0.4);
\draw[popgreen, thick] (0.8,-1.0) .. controls (1.2,-0.6) and (1.6,-0.9) .. (1.9,-0.5);
\fill[popgreen!50] (0,0) ellipse (0.5 and 0.3);
\node[below, font=\small\bfseries, text=popgreen] at (0,-1.8) {Cerveau sain};
\end{scope}
% Cerveau Alzheimer
\begin{scope}[xshift=6.5cm]
\fill[poporangeL] (0,0) ellipse (1.8 and 1.3);
\draw[poporange, thick] (0,0) ellipse (1.8 and 1.3);
% sillons élargis
\draw[poporange, line width=2.2pt] (-1.3,0.4) .. controls (-0.8,0.8) and (-0.3,0.1) .. (0.2,0.6);
\draw[poporange, line width=2.2pt] (-1.1,-0.4) .. controls (-0.6,0.0) and (0.1,-0.7) .. (0.5,-0.2);
\draw[poporange, line width=2.2pt] (0.3,0.9) .. controls (0.7,0.4) and (1.1,0.8) .. (1.4,0.3);
\fill[poporange!60] (0,0) ellipse (0.8 and 0.5);
\draw[poporange, thick] (0,0) ellipse (0.8 and 0.5);
\node[below, font=\small\bfseries, text=poporange] at (0,-1.8) {Cerveau Alzheimer};
\end{scope}
% flèche
\draw[-{Stealth}, very thick, popmuted] (2.7,0) -- (4.0,0) node[midway, above, font=\footnotesize\itshape] {dégénérescence};
\end{tikzpicture}
\end{center}
\legende{Comparaison schématique d'un cerveau sain et d'un cerveau atteint de la maladie d'Alzheimer (coupes horizontales). Chez le malade : atrophie du cortex (cerveau rétréci), \textbf{élargissement des sillons} (traits épais) et \textbf{dilatation des ventricules} (cavité centrale élargie), conséquences de la mort massive des neurones.}
\end{popfigure}

\begin{attention}[Alzheimer n'est pas un vieillissement normal]
Il est fréquent d'entendre : « il oublie, c'est normal à son âge ». \textbf{Faux !} La maladie d'Alzheimer \textbf{n'est pas une conséquence normale du vieillissement}. Un vieillissement normal s'accompagne d'oublis bénins et ponctuels (un prénom qu'on retrouve plus tard), sans perte d'autonomie ni altération des fonctions exécutives. La maladie d'Alzheimer est une \emph{pathologie} : destruction progressive et irréversible des neurones, conduisant à la dépendance totale. De nombreuses personnes très âgées (centenaires) conservent toutes leurs facultés intellectuelles.
\end{attention}

\courssub{Comparaison des trois maladies neurodégénératives étudiées}

Synthétisons dans un tableau les trois affections du contrôle de la motricité et des fonctions cérébrales rencontrées dans cette leçon.

\begin{center}
\begin{tabularx}{\linewidth}{|l|X|X|X|}
\hline
\rowcolor{popblueL} \textbf{Critère} & \textbf{Maladie de Parkinson} & \textbf{Maladie de Huntington} & \textbf{Maladie d'Alzheimer} \\
\hline
Structure atteinte en premier & Substance noire (pars compacta) $\rightarrow$ d\'eficit dopaminergique & Striatum (noyaux gris centraux : noyau caudé, putamen), puis cortex & Hippocampe et cortex cérébral (aire d'association) \\
\hline
Lésions histologiques & Corps de Lewy (agrégats d'$\alpha$-synucléine) dans les neurones survivants & Inclusions nucléaires de huntingtine mutée (polyQ) ; perte neuronale striatale massive & Plaques amyloïdes (A$\beta$ extracellulaire) + Dégénérescences neurofibrillaires (tau intracellulaire) \\
\hline
Symptôme dominant initial & Syndrome parkinsonien : akinésie, rigidité, tremblement de repos & Chorée (mouvements involontaires hyperkinétiques), troubles du comportement & Syndrome amnésique (mémoire récente), désorientation spatiotemporelle \\
\hline
Fonction d'abord touchée & Motricité volontaire (initiation, exécution) & Motricité volontaire (contrôle, inhibition) & Mémoire épisodique et fonctions cognitives supérieures \\
\hline
Transmission génétique & Majoritairement sporadique (facteurs environnementaux + susceptibilité génétique) & \textbf{Héréditaire autosomique dominante} (chr 4, gène \emph{HTT}, expansion CAG) & Majoritairement sporadique (après 65 ans) ; formes familiales rares (< 5 \%, autosomique dominante : gènes \emph{APP}, \emph{PSEN1}, \emph{PSEN2}) \\
\hline
Âge d'apparition typique & 55--65 ans & 35--50 ans (variable selon taille expansion CAG) & Généralement après 65 ans (précoce < 65 ans rare) \\
\hline
Traitement symptomatique principal & L-DOPA (précurseur dopamine) + agonistes dopaminergiques & Neuroleptiques (antidopaminergiques) pour la chorée ; antidépresseurs & Inhibiteurs de l'acétylcholinestérase (donépézil, rivastigmine) ; antagoniste NMDA (mémantine) \\
\hline
\end{tabularx}
\end{center}

\begin{attention}[Ne pas confondre les symptômes dominants -- Piège du Bac]
Au Baccalauréat, l'erreur classique consiste à attribuer les mêmes symptômes aux trois maladies. Retiens la logique structure-fonction :
\begin{itemize}
\item \textbf{Parkinson et Huntington touchent d'abord la MOTRICITÉ} (structures des noyaux de la base : substance noire, striatum).
\item \textbf{Alzheimer touche d'abord la MÉMOIRE et les FONCTIONS COGNITIVES} (cortex, hippocampe).
\item Dans les deux maladies motrices, les signes sont \textbf{opposés} : Parkinson = \emph{déficit} de mouvement (hypokinésie) ; Huntington = \emph{excès} de mouvement (hyperkinésie / chorée).
\end{itemize}
\end{attention}

\courssub{Limitation des dysfonctionnements : prise en charge et prévention}

À ce jour, \textbf{aucun de ces trois dysfonctionnements ne se guérit} : les neurones morts ne se régénèrent pas (absence de neurogenèse significative dans ces aires chez l'adulte). La prise en charge vise donc à \cle{ralentir} l'évolution, à \cle{soulager} les symptômes et à \cle{accompagner} le malade et sa famille.

\begin{itemize}
\item \textbf{Traitements médicamenteux symptomatiques} :
  \begin{itemize}
  \item Parkinson : L-DOPA (précurseur de la dopamine franchissant la BHE) + inhibiteurs de la DOPA-décarboxylase périphérique (carbidopa/benserazide) ; agonistes des récepteurs D2/D3 ; inhibiteurs de la MAO-B et COMT.
  \item Huntington : Neuroleptiques (halopéridol, tiapride) ou tétrabénazine (dépléteur de monoamines) pour atténuer la chorée ; antidépresseurs / anxiolytiques pour les troubles de l'humeur.
  \item Alzheimer : Inhibiteurs de l'acétylcholinestérase (donépézil, rivastigmine, galantamine) pour potentialiser la transmission cholinergique résiduelle ; mémantine (antagoniste des récepteurs NMDA, neuroprotecteur) aux stades modérés à sévères.
  \end{itemize}
\item \textbf{Stimulation cognitive et rééducation} : exercices de mémoire, lecture, jeux, activités sociales, ergothérapie, orthophonie, qui entretiennent les capacités résiduelles et l'autonomie (malades d'Alzheimer et de Huntington).
\item \textbf{Accompagnement familial et social} : aide à la vie quotidienne, sécurisation du domicile (prévention chutes, fugues), soutien psychologique des aidants — rôle essentiel des familles, notamment dans le contexte camerounais où les structures spécialisées (EHPAD, unités Alzheimer, centres de génétique) sont rares et coûteuses.
\item \textbf{Hygiène de vie préventive (facteurs de risque modifiables)} : activité physique régulière (effet neuroprotecteur prouvé), alimentation équilibrée type méditerranéenne (fruits, légumes, poisson, huile d'olive), lutte contre l'\cle{hypertension artérielle}, le \cle{diabète de type 2}, l'obésité, le tabagisme et la sédentarité, qui aggravent les dégénérescences vasculaires et neuronales (hypothèse vasculaire de l'Alzheimer). Niveau d'éducation et réserve cognitive : facteurs protecteurs.
\item \textbf{Recherche en cours (pistes thérapeutiques futures)} : thérapie génique (ARN interférent, ASO pour Huntington ; édition génomique), immunothérapie (anticorps anti-A$\beta$ : lecanemab, donanemab -- efficacité modérée aux stades précoces), cellules souches (remplacement neuronal), neuroprotection mitochondriale. Des essais cliniques sont en cours, mais l'accès reste limité.
\end{itemize}

\begin{exoresolu}[Interpréter un arbre généalogique et calculer un risque]
Énoncé. Dans une famille, le grand-père paternel est atteint de la maladie de Huntington ; la grand-mère paternelle est saine. Leur fils (le père du consultant) est atteint. Le père (atteint) et la mère (saine, sans antécédent familial) ont quatre enfants. On note $H$ l'allèle muté (dominant) et $h$ l'allèle normal.
\begin{enumerate}
\item Déduire les génotypes du grand-père paternel et du père.
\item Calculer le risque, pour chacun des quatre enfants de ce couple, de développer la maladie.
\item Expliquer pourquoi ce risque est indépendant du sexe de l'enfant.
\end{enumerate}
\tcblower
Solution détaillée.
\begin{enumerate}
\item L'allèle muté $H$ est \textbf{dominant} et rare dans la population : les individus atteints sont quasi-systématiquement hétérozygotes $H/h$ (les $H/H$ sont exceptionnels et plus sévères).
  \begin{itemize}
  \item Grand-père paternel (atteint) : génotype \textbf{$H/h$}.
  \item Grand-mère paternelle (saine) : génotype \textbf{$h/h$}.
  \item Leur fils, le père (atteint) : il a reçu l'allèle $h$ de sa mère (seule possibilité) et l'allèle $H$ de son père (pour être atteint). Son génotype est \textbf{$H/h$}.
  \end{itemize}
\item Croisement : Père $H/h$ $\times$ Mère $h/h$.
  \begin{itemize}
  \item Gamètes du père : $H$ (50 \%) et $h$ (50 \%) -- ségrégation méiotique.
  \item Gamètes de la mère : $h$ (100 \%) -- homozygote.
  \end{itemize}
  Fécondation possible :
  \begin{center}
  \begin{tabular}{|c|c|c|}\hline
  & $h$ (mère) & \\ \hline
  $H$ (père) & $H/h$ \textbf{(atteint)} & 1/2 \\ \hline
  $h$ (père) & $h/h$ (sain) & 1/2 \\ \hline
  \end{tabular}
  \end{center}
  Chaque enfant a donc une probabilité de \textbf{1/2 (soit 50 \%)} d'être de génotype $H/h$ et de développer la maladie, et 50 \% d'être $h/h$ et sain.
\item Le gène \emph{HTT} est situé sur le \textbf{chromosome 4 (autosome)}. La ségrégation des allèles et la transmission aux gamètes ne dépendent pas du sexe du parent transmetteur ni du sexe de l'enfant (pas de liaison au sexe). Le risque est donc identique pour les garçons et les filles.
\end{enumerate}
\end{exoresolu}

\begin{exercice}[Analyse d'un document IRM et généalogie]
\textbf{Document 1 :} Coupe IRM cérébrale (pondération T1, axe axial) de deux sujets de 70 ans.
\begin{itemize}
\item Sujet A : Sillons corticaux fins, ventricules latéraux de taille normale, hippocampe bien visible.
\item Sujet B : Sillons corticaux \textbf{largement élargis}, ventricules latéraux \textbf{très dilatés}, hippocampe \textbf{atrophique} (difficile à individualiser).
\end{itemize}
\textbf{Document 2 :} Arbre généalogique d'une famille (symboles standards : carré/homme, cercle/femme, noir/atteint). La maladie apparaît à chaque génération. $I_1$ (homme) et $I_2$ (femme) ont 4 enfants : $II_1$ (homme atteint), $II_2$ (femme saine), $II_3$ (homme atteint), $II_4$ (femme saine). $II_1$ (atteint) épouse $II_5$ (femme saine) : 3 enfants dont 1 atteint ($III_1$), 2 sains.

\begin{enumerate}
\item À quelle pathologie correspond le Document 1 (Sujet B) ? Justifier par deux arguments anatomiques.
\item Quelle est la transmission de la maladie du Document 2 ? Justifier par deux arguments généalogiques.
\item Le sujet $III_1$ (atteint) souhaite avoir un enfant avec une conjointe saine. Quel est le risque pour cet enfant ? Préciser le génotype de $III_1$.
\end{enumerate}
\end{exercice}

\begin{corrige}[Exercice n\degre{} -- Analyse d'un document IRM et généalogie]
\begin{enumerate}
\item \textbf{Maladie d'Alzheimer (stade modéré/sévère).}
  \textbf{Justifications :}
  \begin{itemize}
  \item \textbf{Atrophie corticale diffuse} : élargissement marqué des sillons (espace sous-arachnoïdien élargi par perte de substance grise).
  \item \textbf{Dilatation ventriculaire (hydrocéphalie ex vacuo)} et \textbf{atrophie hippocampique} : perte de substance blanche et grise pariétale/temporale médiale. Ces signes sont les marqueurs d'imagerie de la neurodégénérescence alzheimerienne.
  \end{itemize}
  (Le Sujet A présente une imagerie normale pour l'âge).
\item \textbf{Transmission autosomique dominante.}
  \textbf{Arguments :}
  \begin{itemize}
  \item \textbf{Présence de la maladie à chaque génération} ($I \rightarrow II \rightarrow III$) : pas de saut de génération, excluant une récessive.
  \item \textbf{Atteinte des deux sexes indifféremment} (hommes $I_1, II_1, II_3, III_1$ et femmes potentiellement) et \textbf{proportion ~1/2 des enfants atteints} chez $II_1 \times II_5$ (1/3 ici, compatible avec 1/2 sur petit effectif) : signature d'une dominante autosomique.
  \end{itemize}
\item $III_1$ est atteint, donc de génotype \textbf{$H/h$} (hétérozygote, allèle $H$ dominant rare).
  Croisement $III_1 (H/h) \times$ Conjointe saine ($h/h$).
  Risque pour l'enfant = \textbf{50 \%} (1/2) d'hériter de l'allèle $H$ et de développer la maladie, quel que soit le sexe.
\end{enumerate}
\end{corrige}

\begin{aretenir}[L'essentiel pour le Bac]
\begin{itemize}
\item \textbf{Maladie de Huntington} : \textbf{Héréditaire autosomique dominante} (chr 4, gène \emph{HTT}, expansion CAG $>$ 36). Risque 50 \% pour enfant de parent atteint. Dégénérescence \textbf{striatale} (noyaux gris centraux) $\Rightarrow$ \textbf{chorée} (hyperkinésie) + troubles psychiques + démence. Début 35--50 ans. Pas de traitement curatif.
\item \textbf{Maladie d'Alzheimer} : \textbf{Sporadique le plus souvent} (facteurs âge, vasculaires, APOE $\varepsilon$4). Dégénérescence \textbf{corticale + hippocampique}. Lésions : \textbf{plaques A$\beta$ (extracellulaires)} + \textbf{dégénérescences neurofibrillaires tau (intracellulaires)}. Atrophie cérébrale (IRM). Clinique : \textbf{amnésie hippocampique} (mémoire récente) $\rightarrow$ démence globale. Début $>$ 65 ans. Traitement : anti-acétylcholinestérase + mémantine (symptomatique).
\item \textbf{Distinction clé} : Parkinson (hypokinésie, substance noire/dopamine) vs Huntington (hyperkinésie, striatum/GABA) vs Alzheimer (amnésie/démence, cortex/hippocampe/acétylcholine).
\item \textbf{Prévention} : Contrôle facteurs vasculaires (HTA, diabète), activité physique, stimulation cognitive, alimentation saine.
\end{itemize}
\end{aretenir}

\coursec{Synthèse : relier structure, fonction et dysfonctionnement ; prévention}

Après avoir étudié les dysfonctionnements majeurs que sont les maladies de Huntington et d'Alzheimer, tu disposes maintenant de toutes les pièces du puzzle : tu connais les aires encéphaliques, le trajet de l'influx nerveux lors d'un mouvement volontaire, et les conséquences cliniques des atteintes de chaque structure. Il reste à assembler ces pièces en une démarche cohérente : celle du \emph{raisonnement clinique}, qui permet, à partir d'un symptôme observé chez un malade, de remonter à la structure lésée. C'est exactement le type de raisonnement exigé au Baccalauréat dans les exercices d'interprétation de documents. Cette synthèse te donnera aussi les moyens d'agir : prévenir les atteintes du système nerveux et comprendre comment la rééducation exploite la plasticité cérébrale pour limiter les séquelles.

\courssub{La démarche diagnostique : du symptôme à la structure lésée}

\textbf{L'intuition de départ.} Un médecin du Centre Hospitalier d'Essos à Yaoundé reçoit un patient qui tremble des mains au repos, marche à petits pas et présente une rigidité musculaire. Sans aucun examen d'imagerie, il évoque immédiatement une atteinte de la substance noire. Comment est-ce possible ? Parce que chaque structure nerveuse a une \cle{fonction précise}, et que sa destruction produit un \cle{symptôme caractéristique}. La correspondance symptôme–structure n'est pas un tour de magie : elle repose sur des décennies d'expériences de stimulation, de lésion et, plus récemment, d'imagerie cérébrale.

\begin{methode}[Raisonnement clinique type Bac : remonter du symptôme à la structure]
Face à un document décrivant un malade, procède toujours en quatre étapes :
\begin{enumerate}
\item \textbf{Identifier précisément le symptôme} : paralysie (perte totale du mouvement), tremblement au repos, chorée (mouvements brusques et involontaires), ataxie (troubles de l'équilibre et de la coordination), amnésie (perte de mémoire), etc.
\item \textbf{Formuler une hypothèse de localisation} en utilisant le tableau de correspondance structure–fonction (voir ci-dessous).
\item \textbf{Justifier par un fait expérimental} : citer l'expérience de stimulation ou de lésion qui a établi le rôle de cette structure (stimulation du cortex moteur par Penfield, lésion de la substance noire chez l'animal, etc.).
\item \textbf{Conclure prudemment} : « le symptôme observé est compatible avec une atteinte de \ldots », en vérifiant que le document (autopsie, IRM, scanner, dosage) confirme l'hypothèse.
\end{enumerate}
\end{methode}

\begin{attention}[Une affirmation sans preuve expérimentale vaut zéro]
Au Baccalauréat, écrire « le patient est paralysé parce que son cortex moteur est détruit » sans citer le fait expérimental qui établit le rôle moteur de cette aire est insuffisant. Toute affirmation sur la \cle{localisation d'une fonction} doit être appuyée par un fait : stimulation électrique provoquant un mouvement, lésion expérimentale entraînant un déficit, ou image (IRM, scanner) montrant la zone atteinte. Pense à la triade : \textbf{symptôme — structure — preuve}.
\end{attention}

\courssub{Tableau de synthèse : structures, fonctions et dysfonctionnements}

Le tableau ci-dessous rassemble les correspondances essentielles établies dans les sections précédentes. Il constitue ton outil de diagnostic rapide.

\begin{center}
\begin{tabularx}{\linewidth}{|l|X|X|}
\hline
\rowcolor{popblueL} \textbf{Structure atteinte} & \textbf{Fonction normale} & \textbf{Symptôme / maladie} \\
\hline
Cortex moteur primaire (aire 4 de Brodmann, gyrus précentral) & Élaboration et commande des mouvements volontaires du côté opposé du corps & Paralysie controlatérale (hémiplégie), souvent après un AVC \\
\hline
Cortex somesthésique primaire (gyrus postcentral, aires 1, 2, 3) & Réception et analyse des sensibilités cutanées et profondes & Perte de sensibilité (anesthésie) du côté opposé \\
\hline
Cervelet & Coordination des mouvements, équilibre, tonus musculaire & Ataxie : gestes imprécis, démarche ébrieuse, troubles de l'équilibre \\
\hline
Substance noire (partie compacte, mésencéphale) & Production de dopamine, régulation du mouvement via les noyaux gris centraux & Maladie de Parkinson : tremblement de repos, rigidité, akinésie \\
\hline
Striatum (noyau caudé, putamen) & Filtrage et modulation des programmes moteurs & Maladie de Huntington : chorée (mouvements involontaires brusques), troubles cognitifs \\
\hline
Cortex d'association (diffus) et hippocampe & Mémoire, fonctions cognitives supérieures & Maladie d'Alzheimer : amnésie progressive, perte d'autonomie \\
\hline
\end{tabularx}
\end{center}

\begin{exoresolu}[Un cas clinique au Centre Hospitalier d'Essos]
\textbf{Énoncé.} Un homme de 62 ans consulte pour une paralysie brutale du bras et de la jambe droits, survenue au réveil. La sensibilité est conservée, l'équilibre est normal, il n'y a ni tremblement ni mouvement involontaire. Le scanner (TDM) montre une zone d'infarctus (tissu nécrosé par interruption de la circulation sanguine) dans l'hémisphère cérébral gauche.
\begin{enumerate}
\item Quelle structure est vraisemblablement atteinte ? Justifier.
\item Expliquer pourquoi la paralysie est du côté droit alors que la lésion est à gauche.
\item Pourquoi le cervelet et la substance noire peuvent-ils être écartés ?
\end{enumerate}
\tcblower
\textbf{Solution détaillée.}
\begin{enumerate}
\item La structure atteinte est le \cle{cortex moteur primaire} (gyrus précentral, aire 4 de Brodmann) de l'hémisphère gauche. Justification expérimentale : les stimulations électriques de Penfield sur cette aire chez l'homme conscient déclenchent des mouvements localisés du côté controlatéral, et les lésions de cette même aire (AVC, tumeur) provoquent une paralysie des régions correspondantes du corps. Ici, le déficit est purement moteur (sensibilité conservée), ce qui exclut le cortex somesthésique.
\item Les voies pyramidales (faisceaux cortico-spinaux) issues du cortex moteur se croisent au niveau du bulbe rachidien (\cle{décussation des pyramides}) : chaque hémisphère commande donc la motricité du côté \emph{opposé} (controlatéral) du corps. Une lésion gauche entraîne une hémiplégie droite.
\item L'atteinte du cervelet provoquerait une ataxie (troubles de l'équilibre, de la coordination, démarche ébrueuse), absente ici. L'atteinte de la substance noire provoquerait le syndrome parkinsonien (tremblement de repos, rigidité, akinésie le plus souvent bilatéraux, d'installation progressive), ce qui ne correspond ni au mode de début brutal, ni au caractère strictement unilatéral observé.
\end{enumerate}
\textbf{Conclusion} : le tableau est celui d'un AVC ischémique ayant détruit le cortex moteur gauche, responsable d'une hémiplégie droite.
\end{exoresolu}

\courssub{Limiter les dysfonctionnements : la prévention}

Puisque les neurones de l'encéphale se régénèrent très mal (absence de neurogenèse significative chez l'adulte dans ces aires), la meilleure stratégie reste la \cle{prévention}. Trois grands axes sont à connaître.

\textbf{Prévention des traumatismes crâniens.} Les accidents de la route sont une cause majeure de lésions cérébrales au Cameroun, en particulier chez les usagers des moto-taxis (\emph{bendskins}). Le port du casque homologué, le respect des limitations de vitesse et l'interdiction de conduire sous l'emprise de l'alcool réduisent considérablement le risque de traumatisme crâno-encéphalique.

\begin{savaistu}[Le casque, un gilet de sauvetage pour le cerveau]
Les études épidémiologiques montrent que le port du casque réduit d'environ \SI{70}{\percent} le risque de lésion cérébrale grave en cas d'accident de moto. À Yaoundé et à Douala, où les moto-taxis assurent une grande part des déplacements urbains, le port systématique du casque par le conducteur \emph{et} le passager est l'une des mesures de santé publique les plus efficaces pour limiter les paralysies et les comas post-traumatiques.
\end{savaistu}

\textbf{Lutte contre les accidents vasculaires cérébraux (AVC).} L'AVC, première cause de paralysie acquise de l'adulte, résulte de l'obstruction (AVC ischémique) ou de la rupture (AVC hémorragique) d'un vaisseau cérébral. Ses facteurs de risque sont largement évitables : l'\cle{hypertension artérielle} non traitée (facteur de risque n°1), le tabac, l'alcool, la sédentarité, le diabète et l'hypercholestérolémie. Le contrôle régulier de la tension artérielle dans les centres de santé, une alimentation peu salée, riche en fruits et légumes, et la pratique d'une activité physique régulière constituent une prévention efficace et peu coûteuse.

\textbf{Prévention des infections neurologiques.} Certaines infections atteignent le système nerveux central : méningites bactériennes (prévenues par la vaccination \emph{Haemophilus influenzae} b, pneumocoque, méningocoque et l'hygiène), neurosyphilis, complications neurologiques du paludisme grave (accès pernicieux) ou du VIH/Sida (encéphalopathie, toxoplasmose). La vaccination des enfants (PEV), le traitement rapide des infections, la prise en charge précoce du paludisme (moustiquaires imprégnées, traitement rapide) et la prévention du VIH limitent ces atteintes.

\courssub{Rééducation et plasticité cérébrale}

\textbf{Observation.} De nombreux patients hémiplégiques après un AVC récupèrent une partie de leur motricité en quelques mois, alors que les neurones détruits ne se régénèrent pas. Comment expliquer cette récupération ?

\textbf{Interprétation.} L'imagerie fonctionnelle (IRMf, TEP) montre que, chez ces patients, des aires voisines de la zone lésée (péri-lésionnelles), voire des aires homologues de l'hémisphère controlatéral, s'activent anormalement lors des mouvements du membre rééduqué. Le cerveau \emph{réorganise} ses circuits synaptiques.

\begin{propriete}[Plasticité cérébrale — complément utile au Bac]
La \cle{plasticité cérébrale} est la capacité du système nerveux à réorganiser partiellement ses connexions synaptiques après une lésion : des aires intactes peuvent prendre en charge, au moins partiellement, les fonctions d'une aire détruite. Cette plasticité est :
\begin{itemize}
\item maximale chez le \textbf{sujet jeune} (un enfant hémiplégique récupère souvent remarquablement bien) ;
\item entretenue et amplifiée par l'\textbf{entraînement intensif et répété} : c'est le fondement scientifique de la rééducation ;
\item toujours \textbf{partielle} : elle ne restaure jamais totalement la fonction initiale quand la lésion est étendue.
\end{itemize}
\end{propriete}

La \cle{rééducation} exploite cette plasticité : la \textbf{kinésithérapie} (exercices moteurs répétés, mobilisation passive puis active) stimule la réorganisation des circuits moteurs ; l'\textbf{orthophonie} rééduque le langage après une lésion des aires de Broca ou de Wernicke ; l'\textbf{ergothérapie} aide à réapprendre les gestes de la vie quotidienne. Plus la prise en charge est précoce (dès la phase aiguë stabilisée), plus la récupération fonctionnelle est importante.

\begin{attention}[Plasticité ne signifie pas régénération]
Ne confonds pas plasticité et régénération : les neurones détruits ne sont \emph{pas} remplacés par de nouveaux neurones. Ce sont les neurones survivants qui modifient l'efficacité de leurs connexions synaptiques existantes ou en créent de nouvelles (épines dendritiques). Écris donc « réorganisation des circuits synaptiques » ou « recrutement d'aires vicariantes » et non « repousse des neurones » ou « régénération neuronale », erreur fréquente et sanctionnée au Bac.
\end{attention}

\courssub{Dépistage précoce et prise en charge pluridisciplinaire}

Pour les maladies neurodégénératives (Parkinson, Huntington, Alzheimer), aucun traitement curatif n'existe encore, mais un \cle{dépistage précoce} permet de ralentir l'évolution et d'améliorer la qualité de vie : la \ce{L-DOPA} (précurseur de la dopamine) chez le parkinsonien compense partiellement le déficit dopaminergique ; la prise en charge de la maladie d'Alzheimer associe inhibiteurs de l'acétylcholinestérase, antagonistes des récepteurs NMDA, stimulation cognitive et soutien familial. La prise en charge optimale est \cle{pluridisciplinaire} : médecin neurologue, kinésithérapeute, orthophoniste, ergothérapeute, psychologue et entourage familial collaborent. Dans le contexte camerounais, la solidarité familiale joue un rôle central, qu'il convient d'appuyer par une meilleure information sur la nature non contagieuse de ces maladies et sur l'existence de structures d'appui (associations de malades, services de neurologie des hôpitaux centraux et régionaux).

\courssub{Bilan de la leçon : carte mentale}

\begin{popfigure}
\begin{center}
\begin{tikzpicture}[
  font=\small,
  racine/.style={rectangle, rounded corners, draw=popdark, fill=popdark, text=white, align=center, minimum width=3.5cm, minimum height=1.0cm},
  niv1/.style={rectangle, rounded corners, draw=popblue, fill=popblueL, align=center, minimum width=3.0cm, minimum height=0.85cm},
  niv2/.style={rectangle, rounded corners, draw=poporange, fill=poporangeL, align=center, minimum width=3.0cm, minimum height=0.85cm},
  niv3/.style={rectangle, rounded corners, draw=popgreen, fill=popgreenL, align=center, minimum width=3.0cm, minimum height=0.85cm},
  lien/.style={->, thick, popmuted}
]
\node[racine, align=center] (R) at (0,0){\textbf{Contrôle de la motricité volontaire}};
\node[niv1, align=center] (S) at (-6.0,0){\textbf{Structures encéphaliques}\\ Cortex moteur (aire 4)\\ Cortex somesthésique\\ Cervelet\\ Noyaux gris centraux\\(Substance noire, Striatum)\\ Voie pyramidale};
\node[niv2, align=center] (D) at (6.0,1.8){\textbf{Dysfonctionnements caractéristiques}\\ Hémiplégie (AVC cortex moteur)\\ Ataxie (Cervelet)\\ Syndrome parkinsonien\\(Substance noire)\\ Chorée de Huntington\\(Striatum)\\ Démence (Alzheimer)};
\node[niv3, align=center] (P) at (6.0,-1.8){\textbf{Limitation des séquelles}\\ \textbf{Prévention :} Casque, HTA,\\ Vaccination, Hygiène\\ \textbf{Rééducation :} Plasticité\\ Kinésithérapie, Orthophonie\\ \textbf{Prise en charge :} Précoce,\\ Pluridisciplinaire};
\draw[lien] (R) -- (S);
\draw[lien] (R) -- (D);
\draw[lien] (R) -- (P);
\draw[lien, dashed, poporange] (S) -- node[above, sloped, font=\scriptsize, text=popdark]{Lésion $\rightarrow$} (D);
\draw[lien, dashed, popgreen] (D) -- node[right, font=\scriptsize, text=popdark]{Soins, Prévention} (P);
\end{tikzpicture}
\end{center}
\legende{Carte mentale de synthèse : les structures encéphaliques (à gauche) assurent le contrôle de la motricité volontaire ; leur lésion entraîne des dysfonctionnements caractéristiques (en haut à droite), que l'on peut limiter par la prévention, la rééducation fondée sur la plasticité cérébrale et une prise en charge pluridisciplinaire précoce (en bas à droite).}
\end{popfigure}

\begin{aretenir}[L'essentiel de la leçon]
\begin{itemize}
\item Chaque structure encéphalique a une fonction précise ; sa lésion produit un symptôme caractéristique : cortex moteur $\rightarrow$ paralysie controlatérale ; cervelet $\rightarrow$ ataxie ; substance noire $\rightarrow$ syndrome parkinsonien ; striatum $\rightarrow$ chorée de Huntington ; cortex d'association et hippocampe $\rightarrow$ maladie d'Alzheimer.
\item Le raisonnement diagnostique (type Bac) suit la triade \textbf{symptôme — structure — preuve expérimentale} (stimulation, lésion, imagerie).
\item La prévention (casque homologué, lutte contre l'hypertension artérielle, vaccination, hygiène, lutte antipaludique) est la stratégie la plus efficace, car les neurones cérébraux se régénèrent très mal chez l'adulte.
\item La plasticité cérébrale (réorganisation synaptique, recrutement d'aires vicariantes), maximale chez le sujet jeune, fonde la rééducation et permet une récupération fonctionnelle partielle après lésion.
\item Le dépistage précoce et une prise en charge pluridisciplinaire (médecin, kinésithérapeute, orthophoniste, psychologue, famille) améliorent le pronostic et la qualité de vie des maladies neurologiques.
\end{itemize}
\end{aretenir}

\newpage

\coursec{Activité d'intégration}

\courssub{Objectifs de l'activité}

À l'issue de cette activité d'intégration, tu seras capable de :
\begin{itemize}
\item interpréter des résultats de stimulations corticales et des données d'imagerie cérébrale pour localiser une lésion ;
\item identifier, à partir d'un tableau clinique, la structure encéphalique atteinte (aire motrice, aire sensitive, corps strié, substance noire, cortex diffus) ;
\item construire le schéma fonctionnel du trajet de l'influx nerveux dans le cadre de la motricité volontaire et y situer le niveau de chaque lésion ;
\item exploiter un arbre généalogique pour caractériser un mode de transmission héréditaire ;
\item proposer des mesures de limitation des dysfonctionnements et de prévention adaptées au contexte sanitaire camerounais.
\end{itemize}

\courssub{Situation}

Au service de neurologie de l'\cle{Hôpital Central de Yaoundé}, le Dr Ngo Biyong reçoit, le même matin, trois patients présentant des troubles de la motricité. Les moyens du plateau technique étant limités, le diagnostic repose essentiellement sur l'examen clinique, l'imagerie et le raisonnement scientifique. Tu es l'élève de Terminale D stagiaire du service : le médecin te confie le dossier ci-dessous et te demande de l'aider à établir le diagnostic de chaque patient.

\textbf{Document 1 : tableaux cliniques des trois patients}

\begin{center}
\begin{tabularx}{\linewidth}{|l|X|X|X|}
\hline
\rowcolor{popblueL} & \textbf{Patient A} & \textbf{Patient B} & \textbf{Patient C} \\
\hline Âge / sexe & 24 ans, homme & 62 ans, homme & 38 ans, femme \\
\hline Circonstances & Accident de moto (sans casque) il y a 3 semaines ; aucune fracture au scanner & Apparition progressive depuis 2 ans & Apparition progressive depuis 3 ans \\
\hline Signes observés & Paralysie complète du bras \textbf{droit} ; sensibilité du bras conservée ; aucun tremblement & Tremblement des mains \textbf{au repos} (disparaît pendant le mouvement volontaire), rigidité des membres, lenteur et pauvreté des gestes (akinésie) & Mouvements involontaires, brusques et désordonnés des bras et du visage (chorée) ; troubles de l'humeur \\
\hline Antécédents & Aucun & Aucun & Le père et une tante paternelle ont présenté les mêmes signes vers 40 ans \\
\hline
\end{tabularx}
\end{center}

\textbf{Document 2 : stimulations corticales et imagerie}

Au cours d'une intervention chirurgicale de référence, une équipe a stimulé électriquement (courant de \SI{2}{mA}, durée \SI{0,5}{ms}) différents points du cortex d'un sujet éveillé. Résultats :

\begin{center}
\begin{tabularx}{\linewidth}{|l|X|X|}
\hline
\rowcolor{popblueL} \textbf{Point stimulé} & \textbf{Localisation} & \textbf{Réponse observée} \\
\hline 1 & Circonvolution frontale ascendante (avant le sillon central de Rolando), hémisphère gauche & Contraction des muscles de la main droite \\
\hline 2 & Même circonvolution, hémisphère droit & Contraction des muscles de la jambe gauche \\
\hline 3 & Circonvolution pariétale ascendante (après le sillon central) & Sensation de fourmillement dans le bras opposé, \textbf{sans mouvement} \\
\hline 4 & Structure cérébrale profonde (corps strié) & Aucune réponse motrice directe observable \\
\hline
\end{tabularx}
\end{center}

L'imagerie (IRM) montre : pour le patient A, une destruction localisée de la circonvolution frontale ascendante de l'hémisphère \textbf{gauche} ; pour le patient B, une dégénérescence de la \cle{substance noire} (perte d'environ \SI{70}{\%} des neurones dopaminergiques) ; pour la patiente C, une atrophie marquée du \cle{corps strié} (volume réduit d'environ \SI{30}{\%} par rapport à un témoin du même âge).

\textbf{Document 3 : arbre généalogique de la patiente C}

\begin{popfigure}
\begin{center}
\begin{tikzpicture}[
male/.style={rectangle,draw=popdark,fill=popblueL,minimum size=0.55cm},
female/.style={circle,draw=popdark,fill=poppinkL,minimum size=0.55cm},
affected/.style={fill=poporange},
link/.style={draw=popdark,thick}]
% Génération I
\node[male,affected] (I1) at (0,0) {};
\node[female] (I2) at (1.2,0) {};
\draw[link] (I1) -- (I2);
% Génération II
\node[male] (II1) at (-1.2,-1.6) {};
\node[female,affected] (II2) at (0,-1.6) {};
\node[male] (II3) at (1.2,-1.6) {};
\draw[link] (0.6,0) -- (0.6,-0.8);
\draw[link] (-1.2,-0.8) -- (1.2,-0.8);
\draw[link] (-1.2,-0.8) -- (-1.2,-1.6);
\draw[link] (0,-0.8) -- (0,-1.6);
\draw[link] (1.2,-0.8) -- (1.2,-1.6);
% Génération III (patiente C)
\node[female,affected] (III1) at (0,-3.2) {};
\draw[link] (0,-1.6) -- (0,-3.2);
\node[right=0.15cm of III1, text=popdark] {\small C};
\node[above=0.05cm of I1, text=popdark] {\small I};
\node[left=0.05cm of II2, text=popdark] {\small II};
\node[left=0.05cm of III1, text=popdark] {\small III};
\end{tikzpicture}
\end{center}
\legende{Arbre généalogique de la patiente C. Carré : homme ; cercle : femme ; symbole rempli (orange) : individu atteint ; trait horizontal : union ; trait vertical : descendance. La maladie apparaît à chaque génération et touche les deux sexes.}
\end{popfigure}

\courssub{Consignes}

Par groupes de quatre, tu répondras aux questions suivantes, en rédigeant soigneusement tes justifications :

\begin{enumerate}
\item À partir du document 2, déduis des stimulations 1, 2 et 3 le rôle respectif des circonvolutions frontale et pariétale ascendantes. Que montre la comparaison des stimulations 1 et 2 quant à la latéralisation du contrôle moteur ? Pourquoi la stimulation 4 ne provoque-t-elle aucune réponse motrice directe ?
\item Pour le patient A : identifie la structure lésée et explique pourquoi la paralysie atteint le bras \emph{droit} alors que la lésion est \emph{gauche}. Pourquoi la sensibilité du bras est-elle conservée ?
\item Trace le schéma fonctionnel du trajet de l'influx nerveux lors d'un mouvement volontaire du bras, depuis l'aire motrice jusqu'au muscle effecteur, et situe sur ce schéma le niveau de la lésion de chacun des trois patients.
\item Pour le patient B : nomme la maladie, relie chaque signe clinique au déficit en dopamine et explique pourquoi le tremblement disparaît pendant le mouvement volontaire.
\item Pour la patiente C : à l'aide du document 3, détermine le mode de transmission de la maladie (dominant ou récessif, autosomique ou lié au sexe) en justifiant par deux arguments tirés de l'arbre. Nomme la maladie.
\item Propose, pour chaque patient, une mesure de limitation du dysfonctionnement (traitement ou prise en charge) et, à l'échelle du Cameroun, deux mesures de prévention des causes de tels dysfonctionnements.
\end{enumerate}

\courssub{Ressources à mobiliser}

\begin{itemize}
\item Un \cle{mouvement volontaire} est un mouvement décidé et commandé par le cerveau (cortex cérébral), par opposition au mouvement réflexe, involontaire, commandé par la moelle épinière.
\item L'\cle{aire motrice} (circonvolution frontale ascendante, en avant du sillon central de Rolando) commande les mouvements volontaires du côté \emph{opposé} du corps ; l'\cle{aire sensitive} (circonvolution pariétale ascendante, en arrière du sillon central) reçoit les messages de sensibilité venant du côté opposé du corps.
\item Trajet du mouvement volontaire : aire motrice $\rightarrow$ voie pyramidale (faisceau cortico-spinal, qui s'entrecroise — décussation — au niveau du bulbe rachidien) $\rightarrow$ moelle épinière $\rightarrow$ motoneurone $\rightarrow$ muscle squelettique.
\item Le \cle{corps strié} et la \cle{substance noire} (neurones dopaminergiques) sont des structures encéphaliques profondes qui modulent et facilitent le mouvement volontaire ; leur dysfonctionnement entraîne des troubles moteurs \emph{sans paralysie}.
\item Maladie de Parkinson : destruction des neurones à dopamine de la substance noire $\rightarrow$ tremblement de repos, rigidité, akinésie. Traitement : L-DOPA (précurseur de la dopamine).
\item Maladie de Huntington : maladie héréditaire \emph{autosomique dominante}, atrophie du corps strié $\rightarrow$ chorée (mouvements involontaires désordonnés).
\item Maladie d'Alzheimer : dégénérescence diffuse du cortex cérébral $\rightarrow$ troubles de la mémoire et du comportement, sans trouble moteur initial.
\item Lecture d'un pedigree : une maladie dominante apparaît à chaque génération ; un caractère autosomique touche indifféremment les deux sexes et se transmet aussi bien par le père que par la mère.
\end{itemize}

\courssub{Démarche et corrigé commenté}

\begin{exoresolu}[Le mystère des patients de l'Hôpital Central de Yaoundé]
Énoncé : à partir des trois documents du dossier médical, identifier la structure atteinte chez chaque patient, construire le schéma fonctionnel du trajet de l'influx volontaire, déterminer le mode de transmission de la maladie de la patiente C et proposer des mesures de limitation et de prévention.
\tcblower
\textbf{Question 1 — Rôles des circonvolutions.} La stimulation des points 1 et 2 (circonvolution frontale ascendante) provoque des \emph{mouvements} : c'est donc l'\cle{aire motrice}, siège de la commande des mouvements volontaires. La stimulation du point 3 (circonvolution pariétale ascendante) provoque une \emph{sensation sans mouvement} : c'est l'\cle{aire sensitive}, siège de la réception des sensibilités. La comparaison des points 1 et 2 montre que la stimulation de l'hémisphère gauche fait bouger la main \emph{droite} et celle de l'hémisphère droit la jambe \emph{gauche} : le contrôle moteur est \textbf{contralatéral} (croisé), car la voie pyramidale s'entrecroise (décussation) au niveau du bulbe rachidien. Enfin, la stimulation du point 4 (corps strié) ne déclenche aucun mouvement direct : le corps strié n'est pas un centre de \emph{commande} mais un centre de \emph{modulation} du mouvement, qui agit en influençant l'aire motrice.

\textbf{Question 2 — Patient A.} L'IRM montre une destruction de l'aire motrice de l'hémisphère gauche. En raison du contrôle contralatéral (décussation de la voie pyramidale au bulbe), la commande du côté droit du corps est interrompue : d'où la paralysie du bras droit. L'aire sensitive, située en arrière du sillon central, est intacte, ce qui explique que la sensibilité du bras soit conservée. Il s'agit d'une \cle{paralysie d'origine centrale} (lésion du motoneurone supérieur, dans le cerveau), et non d'une atteinte du nerf périphérique — ce qui est cohérent avec l'absence de fracture au scanner.

\textbf{Question 3 — Schéma fonctionnel.}

\begin{popfigure}
\begin{center}
\begin{tikzpicture}[box/.style={rectangle,draw=popdark,fill=popblueL,rounded corners,align=center,minimum height=0.8cm,font=\small},fleche/.style={-{Latex[length=2mm]},thick,draw=popdark}]
\node[box, align=center] (aire) at (0,0){Aire motrice\\(cortex frontal)};
\node[box,fill=popgoldL, align=center] (modul) at (4.2,1.6){Corps strié + substance noire\\(modulation, dopamine)};
\node[box, align=center] (pyra) at (8.4,0){Voie pyramidale\\(décussation au bulbe)};
\node[box, align=center] (moelle) at (12.6,0){Moelle épinière\\(motoneurone)};
\node[box,fill=popgreenL, align=center] (muscle) at (8.4,-2.2){Muscle effecteur\\(mouvement du bras)};
\draw[fleche] (aire) -- (pyra);
\draw[fleche] (pyra) -- (moelle);
\draw[fleche] (moelle) |- (muscle);
\draw[fleche,dashed,draw=poporange] (modul) -- (pyra) node[midway,above,font=\scriptsize,text=poporange]{facilitation};
\node[font=\scriptsize\bfseries,text=red] at (0,-0.9) {Lésion A};
\node[font=\scriptsize\bfseries,text=red] at (4.2,2.6) {Lésions B et C};
\end{tikzpicture}
\end{center}
\legende{Schéma fonctionnel du trajet de l'influx nerveux lors d'un mouvement volontaire. Flèches pleines : voie de commande (aire motrice $\rightarrow$ voie pyramidale $\rightarrow$ moelle $\rightarrow$ muscle). Flèche pointillée orange : modulation facilitatrice exercée par le corps strié et la substance noire. Lésion A : destruction de l'aire motrice (paralysie). Lésion B : dégénérescence de la substance noire (Parkinson). Lésion C : atrophie du corps strié (Huntington).}
\end{popfigure}

\textbf{Question 4 — Patient B.} Il s'agit de la \cle{maladie de Parkinson}. La perte d'environ \SI{70}{\%} des neurones dopaminergiques de la substance noire supprime la facilitation que la dopamine exerce sur les circuits du mouvement volontaire : d'où la \emph{lenteur et la pauvreté gestuelle} (akinésie) et la \emph{rigidité} musculaire. Le \emph{tremblement de repos} traduit l'activité déséquilibrée de ces circuits moteurs lorsqu'ils ne sont pas sollicités ; il disparaît pendant le mouvement volontaire car la commande corticale, restée intacte, réengage alors correctement ces circuits. Prise en charge : la \cle{L-DOPA}, précurseur de la dopamine capable de franchir la barrière hémato-encéphalique (que la dopamine elle-même ne franchit pas), compense partiellement le déficit.

\textbf{Question 5 — Patiente C.} Il s'agit de la \cle{maladie de Huntington} (chorée de Huntington), diagnostic cohérent avec la chorée et l'atrophie du corps strié. Analyse du pedigree : (i) la maladie apparaît à \emph{chaque génération} (I, II et III) et tout individu atteint a un parent atteint $\rightarrow$ transmission \textbf{dominante} (un allèle morbide suffit) ; (ii) les \emph{deux sexes} sont touchés (homme en I, femmes en II et III) et la transmission se fait ici \emph{par le père} vers sa fille, sans différence de distribution entre les sexes $\rightarrow$ le gène est porté par un \textbf{autosome}. La maladie est donc \textbf{autosomique dominante} : chaque enfant d'un parent atteint hétérozygote a une probabilité de $1/2$ de recevoir l'allèle morbide et d'être atteint.

\textbf{Question 6 — Limitation et prévention.}
\begin{itemize}
\item Patient A : rééducation par kinésithérapie (la plasticité cérébrale permet parfois une récupération partielle des fonctions) ; prévention : port obligatoire et contrôlé du casque à moto, sensibilisation à la sécurité routière (les accidents de la circulation sont une cause majeure de traumatismes crâniens chez les jeunes au Cameroun).
\item Patient B : traitement par L-DOPA, kinésithérapie d'entretien pour limiter la rigidité ; prévention : limitation de l'exposition aux pesticides dans les exploitations agricoles (facteur de risque suspecté), port d'équipements de protection.
\item Patiente C : prise en charge symptomatique, soutien psychologique, conseil génétique pour les descendants (information sur le risque de $1/2$, dépistage prédictif possible) ; prévention : information et accompagnement des familles à risque dans les centres de santé.
\end{itemize}
\end{exoresolu}

\begin{attention}[Erreurs fréquentes à éviter au Bac]
\begin{itemize}
\item Ne confonds pas \cle{aire motrice} (en avant du sillon central de Rolando) et \cle{aire sensitive} (en arrière du sillon central) : une lésion de l'aire motrice supprime le mouvement mais \emph{conserve} la sensibilité.
\item N'oublie jamais la \cle{décussation} : une lésion corticale gauche se traduit par un déficit du côté \emph{droit} du corps (contrôle contralatéral).
\item Parkinson et Huntington ne sont pas des paralysies : la voie pyramidale est intacte ; ce sont des troubles de la \emph{modulation} du mouvement (structures profondes : substance noire, corps strié).
\item Dans un pedigree, « présent à chaque génération » est l'argument de dominance ; « les deux sexes touchés et transmission par le père comme par la mère » est l'argument du caractère autosomique. Rédige ces deux arguments séparément.
\item La dopamine ne traverse pas la barrière hémato-encéphalique : c'est pourquoi on administre son précurseur, la L-DOPA, et non la dopamine elle-même.
\end{itemize}
\end{attention}

\begin{savaistu}[Parkinson, Huntington, Alzheimer : trois dégénérescences, trois tableaux]
Ces trois maladies sont des affections dégénératives du système nerveux, mais elles diffèrent par la structure atteinte et donc par les signes : la maladie de Parkinson (substance noire) touche le \emph{lancement} du mouvement, la maladie de Huntington (corps strié) provoque des mouvements \emph{involontaires}, et la maladie d'Alzheimer (cortex diffus) altère d'abord la \emph{mémoire} et le comportement, sans trouble moteur initial. Au Cameroun, le vieillissement de la population rend ces maladies de plus en plus fréquentes dans les services de neurologie des hôpitaux de Yaoundé et de Douala, alors que l'accès aux traitements (L-DOPA, prise en charge) reste inégal : d'où l'importance du diagnostic clinique rigoureux et de la prévention.
\end{savaistu}

\courssub{Bilan de l'activité}

Cette activité t'a permis de relier rigoureusement \cle{structure}, \cle{fonction} et \cle{dysfonctionnement} : une même fonction — le contrôle de la motricité volontaire — peut être altérée à trois niveaux différents (aire motrice corticale, substance noire, corps strié), et chaque niveau de lésion produit un tableau clinique distinct et reconnaissable. Tu as mobilisé les trois savoir-faire exigibles de la leçon : interpréter des expériences de stimulation et de lésion corticales, construire un schéma fonctionnel du trajet de l'influx nerveux, et exploiter un arbre généalogique. Tu as enfin compris que le raisonnement clinique du neurologue est une démarche scientifique à part entière : observation des signes, formulation d'une hypothèse diagnostique, confrontation aux données expérimentales (imagerie, stimulations, pedigree), conclusion et proposition de prise en charge. Cette démarche, appliquée au contexte sanitaire camerounais (accidents de moto, accès aux traitements, conseil génétique), te prépare directement aux épreuves d'intégration du Baccalauréat D.

\newpage

\coursec{Méthodes et exercices résolus}

\coursec{Limitation des dysfonctionnements des structures responsables du contrôle de la motricité}

\courssub{Savoir-faire 1 : Disséquer l'encéphale d'un Mammifère et réaliser un schéma annoté}

\begin{methode}[Réussir la dissection et le schéma d'un encéphale]
\begin{enumerate}
\item \cle{Préparer} le matériel : encéphale frais de mouton ou de bœuf (obtenu à l'abattoir), cuvette de dissection, scalpel, ciseaux fins, pinces, gants, blouse. Respecter les règles d'hygiène (lavage des mains, nettoyage du plan de travail).
\item \cle{Observer d'abord la face dorsale} : repérer les deux hémisphères cérébraux séparés par la \emph{scissure interhémisphérique} (fissure longitudinale), les circonvolutions (replis du cortex), le cervelet à l'arrière et le tronc cérébral en dessous.
\item \cle{Observer la face ventrale} : repérer les bulbes olfactifs, le chiasma optique (croisement des nerfs optiques), l'hypophyse, le pont et le bulbe rachidien prolongé par la moelle épinière.
\item \cle{Pratiquer une coupe sagittale médiane} (dans le plan de symétrie) : on distingue la substance grise (cortex, en périphérie) de la substance blanche (axones myélinisés, au centre), le corps calleux (pont de fibres reliant les deux hémisphères), l'hypothalamus et le tronc cérébral.
\item \cle{Réaliser le schéma} : traits nets au crayon, pas de hachures, proportions respectées ; légendes horizontales reliées par des traits fins qui ne se croisent jamais ; titre souligné en dessous du schéma.
\end{enumerate}
\begin{attention}
Un schéma anatomique n'est pas un dessin artistique : pas d'ombres, pas de couleurs fantaisistes, une légende par structure identifiée. Au Bac, un schéma sans titre ni légende perd systématiquement des points.
\end{attention}
\end{methode}

\begin{popfigure}
\begin{center}
\begin{tikzpicture}[scale=0.7, every node/.style={font=\small}]
% Vue Dorsale
\begin{scope}[xshift=-5cm]
\draw[popblue, thick] (-2.5,0) .. controls (-2.5,2) and (2.5,2) .. (2.5,0) .. controls (2.5,-1) and (-2.5,-1) .. (-2.5,0); % Hémisphères
\draw[popblue, thick] (-2.5,0) .. controls (-2.5,-0.5) and (2.5,-0.5) .. (2.5,0); % Scissure
\draw[popdark, thin] (-2.2,0.5) .. controls (-2,1) .. (-1.5,1.2);
\draw[popdark, thin] (-1,0.8) .. controls (-0.5,1.3) .. (0,1.4);
\draw[popdark, thin] (1,0.8) .. controls (0.5,1.3) .. (0,1.4);
\draw[popdark, thin] (2.2,0.5) .. controls (2,1) .. (1.5,1.2);
\draw[popdark, thin] (-2.2,-0.2) .. controls (-2,-0.5) .. (-1.5,-0.7);
\draw[popdark, thin] (2.2,-0.2) .. controls (2,-0.5) .. (1.5,-0.7);
% Cervelet
\draw[popgreen, thick] (-1.8,-1.2) .. controls (-1.8,-2) and (1.8,-2) .. (1.8,-1.2) .. controls (1.8,-1) and (-1.8,-1) .. (-1.8,-1.2);
\foreach \x in {-1.5,-1,-0.5,0,0.5,1,1.5} \draw[popgreen, thin] (\x,-1.2) -- (\x,-1.9);
\node[popblue] at (0,1.8) {\textbf{Face dorsale}};
\node[popblue, font=\tiny] at (0,0.5) {Hémisphères cérébraux};
\node[popblue, font=\tiny] at (0,-0.2) {Scissure interhémisphérique};
\node[popgreen, font=\tiny] at (0,-1.6) {Cervelet};
\node[popdark, font=\tiny] at (0,-2.3) {Tronc cérébral (sous cervelet)};
\end{scope}

% Vue Ventrale
\begin{scope}[xshift=0cm]
\draw[popblue, thick] (-2.5,0) .. controls (-2.5,2) and (2.5,2) .. (2.5,0) .. controls (2.5,-1) and (-2.5,-1) .. (-2.5,0);
\draw[popblue, thick] (-2.5,0) .. controls (-2.5,-0.5) and (2.5,-0.5) .. (2.5,0);
% Bulbes olfactifs
\draw[poporange, thick] (-1.5,1.8) circle (0.3);
\draw[poporange, thick] (1.5,1.8) circle (0.3);
% Chiasma optique
\draw[poppurple, thick] (-0.8,0.5) -- (-0.3,0.2) -- (-0.8,-0.1);
\draw[poppurple, thick] (0.8,0.5) -- (0.3,0.2) -- (0.8,-0.1);
% Hypophyse
\draw[poppink, thick] (0,0.2) ellipse (0.3 and 0.15);
% Tronc cérébral
\draw[popdark, thick] (-0.5,-0.5) -- (-0.5,-1.5); % Pont
\draw[popdark, thick] (0.5,-0.5) -- (0.5,-1.5);
\draw[popdark, thick] (-0.5,-1.5) -- (-0.5,-2.5); % Bulbe
\draw[popdark, thick] (0.5,-1.5) -- (0.5,-2.5);
\draw[popdark, thick] (-0.3,-2.5) -- (-0.3,-3.5); % Moelle
\draw[popdark, thick] (0.3,-2.5) -- (0.3,-3.5);
\node[popblue] at (0,2.3) {\textbf{Face ventrale}};
\node[poporange, font=\tiny] at (-1.5,2.3) {Bulbes olfactifs};
\node[poppurple, font=\tiny] at (0,0.8) {Chiasma optique};
\node[poppink, font=\tiny] at (0,0.2) {Hypophyse};
\node[popdark, font=\tiny] at (0,-1) {Pont};
\node[popdark, font=\tiny] at (0,-2) {Bulbe rachidien};
\node[popdark, font=\tiny] at (0,-3) {Moelle épinière};
\end{scope}

% Coupe Sagittale Médiane
\begin{scope}[xshift=5cm]
\draw[popblue, thick] (-1.5,0) .. controls (-1.5,2.5) and (1.5,2.5) .. (1.5,0) .. controls (1.5,-1) and (-1.5,-1) .. (-1.5,0); % Hémisphère
\draw[popblue!50!white, thick] (-1.3,0.2) .. controls (-1.3,2.3) and (1.3,2.3) .. (1.3,0.2); % Cortex (Grise)
\draw[popblue, thick] (-1.5,0) -- (-1.5,-3.5); % Ligne médiane
% Corps calleux
\draw[popgold, thick] (-1.3,0.5) .. controls (-1.3,0.8) and (1.3,0.8) .. (1.3,0.5);
\draw[popgold, thick] (-1.3,0.5) .. controls (-1.3,0.3) and (1.3,0.3) .. (1.3,0.5);
% Substance blanche (hachures légères)
\foreach \y in {0.6,0.9,1.2,1.5,1.8,2.1} \draw[popmuted, thin] (-1.2,\y) -- (1.2,\y);
% Cervelet
\draw[popgreen, thick] (0.5,-1) .. controls (0.5,-1.8) and (1.5,-1.8) .. (1.5,-1) .. controls (1.5,-0.8) and (0.5,-0.8) .. (0.5,-1);
\foreach \x in {0.7,0.9,1.1,1.3} \draw[popgreen, thin] (\x,-1) -- (\x,-1.7);
% Tronc / Hypothalamus
\draw[popdark, thick] (0,-0.5) -- (0,-1.5); % Tronc
\draw[popred, thick] (0,0) circle (0.2); % Hypothalamus
\node[popblue] at (0,2.8) {\textbf{Coupe sagittale médiane}};
\node[popblue!50!white, font=\tiny] at (1.8,1.5) {Substance grise (Cortex)};
\node[popmuted, font=\tiny] at (1.8,1) {Substance blanche};
\node[popgold, font=\tiny] at (1.8,0.5) {Corps calleux};
\node[popred, font=\tiny] at (0.5,0) {Hypothalamus};
\node[popdark, font=\tiny] at (0.5,-1) {Tronc cérébral};
\node[popgreen, font=\tiny] at (1.5,-1.5) {Cervelet};
\end{scope}
\end{tikzpicture}
\end{center}
\legende{Organisation macroscopique de l'encéphale de mammifère (mouton/bœuf). \textbf{Gauche :} Face dorsale (hémisphères, scissure, cervelet). \textbf{Centre :} Face ventrale (bulbes olfactifs, chiasma optique, hypophyse, tronc cérébral). \textbf{Droite :} Coupe sagittale médiane (cortex/substance grise périphérique, substance blanche centrale, corps calleux, cervelet, tronc).}
\end{popfigure}

\begin{exoresolu}[Dissection d'un encéphale de mouton]
Au laboratoire du lycée, on te remet un encéphale de mouton. 1) Décris le protocole permettant d'identifier les grandes parties de cet encéphale. 2) Cite quatre structures que l'on peut observer sur la face ventrale. 3) Quelle précaution faut-il prendre pour réaliser une coupe sagittale médiane ?
\tcblower
\textbf{1) Protocole.} On place l'encéphale dans la cuvette de dissection, face dorsale vers le haut. On observe d'abord sans toucher : les deux hémisphères cérébraux, volumineux et plissés en circonvolutions, séparés par la scissure interhémisphérique ; en arrière et en dessous, le cervelet, plus petit, à surface finement striée ; en avant du cervelet, le tronc cérébral. On retourne ensuite l'encéphale pour observer la face ventrale. Enfin, à l'aide du scalpel, on pratique une coupe sagittale médiane qui révèle l'organisation interne (substance grise périphérique, substance blanche centrale, corps calleux).

\textbf{2) Structures visibles sur la face ventrale.} On peut citer : les \cle{bulbes olfactifs} (à l'avant), le \cle{chiasma optique} (croisement en X des nerfs optiques), l'\cle{hypophyse} (petite glande sous l'hypothalamus), le \cle{pont} (renflement du tronc cérébral) et le \cle{bulbe rachidien} prolongé par la moelle épinière.

\textbf{3) Précaution pour la coupe.} Il faut sectionner exactement dans le plan de symétrie de l'encéphale, en suivant la scissure interhémisphérique comme repère, d'un geste franc et continu, sans scier, afin d'obtenir deux moitiés symétriques où le corps calleux et le tronc cérébral apparaissent nettement en coupe.

\textbf{Conclusion.} L'observation méthodique (face dorsale, face ventrale, puis coupe sagittale médiane) permet d'identifier les trois grandes parties de l'encéphale : le cerveau (hémisphères cérébraux), le cervelet et le tronc cérébral.
\end{exoresolu}

\courssub{Savoir-faire 2 : Interpréter des expériences de stimulation ou de lésions corticales}

\begin{methode}[Interpréter une expérience de stimulation ou de lésion du cortex]
\begin{enumerate}
\item \cle{Identifier la variable expérimentale} : quelle zone du cortex est stimulée (électriquement) ou détruite (lésion) ? De quel côté ?
\item \cle{Observer l'effet} : mouvement déclenché (stimulation) ou déficit apparu (lésion) — paralysie, anesthésie, trouble du langage... Préciser de quel côté du corps il se manifeste.
\item \cle{Appliquer la règle du croisement} : chaque hémisphère contrôle la moitié \emph{opposée} (controlatérale) du corps. Une stimulation du cortex moteur gauche provoque un mouvement du côté droit.
\item \cle{Conclure par déduction} : la zone stimulée ou lésée est le centre nerveux responsable de la fonction perturbée. Si des points différents commandent des muscles différents, on conclut à une \emph{organisation somatotopique} (chaque région du corps a sa zone corticale propre).
\item \cle{Comparer aires motrices et aires sensitives} : stimulation de l'aire motrice $\Rightarrow$ mouvement ; stimulation de l'aire sensitive $\Rightarrow$ sensation (fourmillement, chaleur) ; lésion de l'aire sensitive $\Rightarrow$ perte de sensibilité sans paralysie.
\end{enumerate}
\end{methode}

\begin{popfigure}
\begin{center}
\begin{tikzpicture}[scale=0.9, every node/.style={font=\small}]
% Hémisphère gauche vu de côté
\draw[popblue, thick] (-3,0) .. controls (-3,3) and (3,3) .. (3,0) .. controls (3,-1.5) and (-3,-1.5) .. (-3,0);
% Sillon central
\draw[popdark, thick, dashed] (0,3) -- (0,-1.5);
\node[popdark, rotate=90] at (0.3,1) {Sillon central (Rolando)};
% Aire Motrice (Précentrale) - Antérieure au sillon
\draw[poporange, fill=poporangeL, opacity=0.5] (-2.5,0.5) .. controls (-2.5,2.5) and (-0.5,2.5) .. (-0.5,0.5) -- cycle;
\node[poporange] at (-1.5,1.8) {\textbf{Aire motrice}};
\node[poporange, font=\tiny] at (-1.5,1.3) {(Gyrus précentral)};
% Homonculus Moteur (simplifié)
\node[font=\tiny] at (-2.2,2.3) {Pied};
\node[font=\tiny] at (-1.8,2.0) {Jambe};
\node[font=\tiny] at (-1.4,1.6) {Tronc};
\node[font=\tiny] at (-1.0,1.2) {Bras};
\node[font=\tiny] at (-0.7,0.8) {Main};
\node[font=\tiny] at (-0.6,0.4) {Visage};
\node[font=\tiny] at (-0.5,0.1) {Langue};

% Aire Sensitive (Postcentrale) - Postérieure au sillon
\draw[popgreen, fill=popgreenL, opacity=0.5] (0.5,0.5) .. controls (0.5,2.5) and (2.5,2.5) .. (2.5,0.5) -- cycle;
\node[popgreen] at (1.5,1.8) {\textbf{Aire sensitive}};
\node[popgreen, font=\tiny] at (1.5,1.3) {(Gyrus postcentral)};
% Homonculus Sensitif
\node[font=\tiny] at (0.7,2.3) {Pied};
\node[font=\tiny] at (1.0,2.0) {Jambe};
\node[font=\tiny] at (1.3,1.6) {Tronc};
\node[font=\tiny] at (1.6,1.2) {Bras};
\node[font=\tiny] at (1.9,0.8) {Main};
\node[font=\tiny] at (2.1,0.4) {Visage};
\node[font=\tiny] at (2.2,0.1) {Langue};

% Flèches croisées
\draw[->, thick, popdark] (-1.5,-0.5) -- (-1.5,-1) node[below, font=\tiny] {Contrôle côté DROIT};
\draw[->, thick, popdark] (1.5,-0.5) -- (1.5,-1) node[below, font=\tiny] {Sensibilité côté DROIT};
\end{tikzpicture}
\end{center}
\legende{Localisation des aires motrice et sensitive de l'hémisphère gauche. Le sillon central (de Rolando) les sépare. L'organisation somatotopique (homonculus) est inversée : le bas du corps (pied) est en haut, la face/langue en bas. Chaque hémisphère contrôle la moitié controlatérale du corps.}
\end{popfigure}

\begin{exoresolu}[Les expériences de stimulation corticale chez le singe]
Chez un singe anesthésié, on stimule électriquement différents points du cortex cérébral. Résultats :

\begin{center}
\begin{tabularx}{\linewidth}{|c|X|X|}
\hline
\rowcolor{popblueL} Expérience & Zone stimulée & Effet observé \\
\hline
1 & Point A de la circonvolution frontale ascendante gauche & Mouvement de la patte antérieure droite \\
\hline
2 & Point B de la même circonvolution gauche & Mouvement des muscles de la face, côté droit \\
\hline
3 & Point C de la circonvolution pariétale ascendante gauche & Le singe (anesthésie légère) réagit comme si on touchait sa patte droite \\
\hline
4 & Lésion du point A & Paralysie durable de la patte antérieure droite \\
\hline
\end{tabularx}
\end{center}

1) Interprète les expériences 1 et 2. 2) Que révèle la comparaison des expériences 1 et 3 ? 3) Déduis de l'expérience 4 le rôle du point A.
\tcblower
\textbf{1) Interprétation des expériences 1 et 2.} La stimulation de deux points distincts (A et B) d'une même circonvolution déclenche des mouvements de deux régions différentes du corps (patte, face). On en déduit que la circonvolution frontale ascendante (gyrus précentral) est une \cle{aire motrice} : elle commande les mouvements volontaires. De plus, chaque point de cette aire commande un groupe musculaire précis : c'est la preuve d'une \cle{organisation somatotopique} du cortex moteur (représentation du corps point par point, appelée \emph{homonculus moteur}). Enfin, la stimulation du côté gauche provoque un mouvement du côté droit : le contrôle moteur est \cle{controlatéral} (croisé), car les voies pyramidales se croisent au niveau du bulbe rachidien (décussation des pyramides).

\textbf{2) Comparaison des expériences 1 et 3.} La stimulation de la circonvolution frontale ascendante (point A, gyrus précentral) provoque un \emph{mouvement}, tandis que celle de la circonvolution pariétale ascendante (point C, gyrus postcentral) provoque une \emph{sensation} tactile. Il existe donc deux aires distinctes, séparées par le sillon central : l'\cle{aire motrice} (commande des mouvements) et l'\cle{aire sensitive} (réception et analyse des sensations cutanées), également organisée somatotopiquement (homonculus sensitif).

\textbf{3) Rôle du point A.} La destruction du point A entraîne une paralysie durable de la patte antérieure droite : le point A est donc le \cle{centre cortical moteur} de cette patte. Sans lui, l'ordre moteur volontaire ne peut plus être élaboré ni transmis au muscle : la paralysie est d'origine centrale (le muscle et le nerf périphérique sont intacts).

\textbf{Conclusion.} Ces expériences établissent que le cortex cérébral comporte des aires motrices et sensitives spécialisées, organisées somatotopiquement et contrôlant la moitié opposée du corps.
\end{exoresolu}

\courssub{Savoir-faire 3 : Expliquer le trajet du message nerveux lors d'un mouvement volontaire}

\begin{methode}[Construire le schéma fonctionnel du mouvement volontaire]
\begin{enumerate}
\item \cle{Définir le mouvement volontaire} : mouvement décidé, programmé et contrôlé par le cerveau, impliquant la contraction de muscles squelettiques.
\item \cle{Ordonner les étapes} dans le bon sens (de la décision au muscle) :
\begin{itemize}
\item élaboration du programme moteur dans l'\emph{aire motrice} (cortex, gyrus précentral) par les neurones pyramidaux (motoneurones supérieurs) ;
\item correction et coordination par le \emph{cervelet} et les \emph{noyaux gris centraux} (dont le locus niger / substance noire) ;
\item descente de l'influx par la \emph{voie pyramidale} (faisceau cortico-spinal latéral) dans la substance blanche hémisphérique, pédoncules cérébraux, protubérance, pyramides du bulbe ;
\item \emph{décussation} (croisement) des fibres au niveau du bulbe rachidien (décussation des pyramides) ;
\item descente dans la \emph{voie cortico-spinale latérale} de la moelle épinière controlatérale ;
\item relais dans la \emph{corne ventrale de la moelle épinière} sur le motoneurone (motoneurone inférieur) ;
\item \emph{nerf moteur} (racine ventrale) jusqu'au muscle effecteur (jonction neuromusculaire) qui se contracte.
\end{itemize}
\item \cle{Ajouter la boucle de retour (rétroaction)} : les récepteurs musculaires (fuseaux neuromusculaires) et articulaires (organes tendineux de Golgi) renvoient des informations sensitives (proprioception) vers le cervelet (voie spinocérébelleuse) et le cortex (aire sensitive via lemnisque médian et thalamus) pour le contrôle et l'ajustement du mouvement en temps réel.
\item \cle{Vérifier la logique} : toute lésion d'un maillon de cette chaîne entraîne un déficit précis et localisable (raisonnement « structure $\rightarrow$ fonction $\rightarrow$ dysfonctionnement »).
\end{enumerate}
\end{methode}

\begin{exoresolu}[Tracer le trajet de l'influx lors d'un mouvement volontaire]
Un élève décide d'attraper son stylo tombé par terre. 1) Définis le mouvement volontaire. 2) Décris, sous forme de schéma fonctionnel fléché, le trajet suivi par le message nerveux depuis la décision jusqu'à la contraction des muscles de la main droite. 3) Précise le rôle du cervelet dans ce mouvement.
\tcblower
\textbf{1) Définition.} Un \cle{mouvement volontaire} est un mouvement décidé et contrôlé par la volonté : il naît d'une décision prise au niveau du cortex cérébral et se réalise par la contraction coordonnée de muscles squelettiques.

\textbf{2) Schéma fonctionnel du trajet de l'influx.}

\begin{popfigure}
\begin{center}
\begin{tikzpicture}[node distance=0.4cm and 0.8cm,
box/.style={draw=popblue, fill=popblueL, rounded corners=3pt, align=center, font=\small, inner sep=4pt, minimum width=3.5cm},
arr/.style={-{Stealth[length=3mm]}, thick, popdark},
arrfb/.style={-{Stealth[length=2.5mm]}, thick, poporange, dashed}]
\node[box, align=center] (a){1. Décision \& Programmation\\ Aire motrice (Gyrus précentral)\\ \emph{Hémisphère GAUCHE}\\ (Neurones pyramidaux / Motoneurones supérieurs)};
\node[box, below=of a, align=center] (b){2. Coordination \& Régulation\\ Cervelet \& Noyaux gris centraux\\ (Locus niger / Striatum)};
\node[box, below=of b, align=center] (c){3. Descente\\ Voie pyramidale (Substance blanche,\\ Pédoncules, Pont, Pyramides du bulbe)};
\node[box, below=of c, align=center] (d){4. Décussation des pyramides\\ (Bulbe rachidien)\\ \emph{Croisement controlatéral}};
\node[box, below=of d, align=center] (e){5. Voie cortico-spinale latérale\\ (Moelle épinière, hémisphère DROIT)};
\node[box, below=of e, align=center] (f){6. Relais spinal\\ Corne ventrale de la moelle\\ (Motoneurone inférieur / $\alpha$)};
\node[box, below=of f, align=center] (g){7. Nerf moteur (Racine ventrale)\\ $\rightarrow$ Jonction neuromusculaire\\ $\rightarrow$ Muscles main DROITE};
\draw[arr] (a) -- (b);
\draw[arr] (b) -- (c);
\draw[arr] (c) -- (d);
\draw[arr] (d) -- (e);
\draw[arr] (e) -- (f);
\draw[arr] (f) -- (g);
% Boucle de retour
\draw[arrfb] (g.west) to[out=180,in=180, looseness=1.5] node[left, font=\scriptsize, align=center, text=poporange, xshift=-0.5cm]{Proprioception\\ (Fuseaux, Golgi)\\ $\rightarrow$ Voies sensitives\\ $\rightarrow$ Cervelet / Cortex sensitive} (a.west);
\end{tikzpicture}
\end{center}
\legende{Trajet du message nerveux moteur volontaire (main droite). Flèches pleines : voie descendante (motrice). Flèches pointillées orange : boucle de rétroaction sensitive (proprioception) assurant le contrôle et la correction du mouvement.}
\end{popfigure}

Justification des étapes : l'ordre moteur est élaboré par les neurones pyramidaux de l'aire motrice (motoneurones supérieurs) ; leurs axones descendent dans la voie pyramidale, se croisent au niveau du bulbe rachidien (décussation des pyramides), ce qui explique que l'hémisphère gauche commande la main droite ; l'influx descend dans le faisceau cortico-spinal latéral, fait relais sur le motoneurone (motoneurone inférieur) de la corne ventrale de la moelle épinière, puis emprunte la racine ventrale et le nerf moteur jusqu'aux muscles effecteurs.

\textbf{3) Rôle du cervelet.} Le cervelet ne déclenche pas le mouvement mais le \cle{coordonne} : il compare en permanence le programme moteur envoyé par le cortex (via les noyaux gris centraux et le pont) aux informations sensitives revenant des muscles et des articulations (proprioception), et corrige la force, l'amplitude, la vitesse et la trajectoire. Grâce à lui, le mouvement est précis, fluide et équilibré (attraper le stylo sans le rater, sans tremblement d'intention).

\textbf{Conclusion.} Le mouvement volontaire résulte d'une chaîne : cortex moteur (décision) $\rightarrow$ noyaux gris centraux / cervelet (régulation) $\rightarrow$ voie pyramidale croisée $\rightarrow$ moelle (relais) $\rightarrow$ nerf moteur $\rightarrow$ muscle, sous le contrôle permanent de la rétroaction sensitive (cervelet, aire sensitive).

\end{exoresolu}

\courssub{Savoir-faire 4 : Relier une lésion ou une dégénérescence à un dysfonctionnement moteur}

\begin{methode}[Raisonner « structure $\rightarrow$ fonction $\rightarrow$ symptôme »]
\begin{enumerate}
\item \cle{Identifier la structure atteinte} dans les données (noyaux gris centraux, cortex, neurones dopaminergiques du locus niger, striatum, hippocampe...).
\item \cle{Rappeler la fonction normale} de cette structure (ex. : le locus niger produit la dopamine, neuromédiateur indispensable au déclenchement harmonieux des mouvements via la voie directe/indirecte des noyaux gris).
\item \cle{Déduire le symptôme} par logique causale : si la structure ne fonctionne plus, sa fonction est perdue ou perturbée $\Rightarrow$ symptôme correspondant.
\item \cle{Confronter aux données cliniques} : vérifier que les symptômes décrits (tremblement, rigidité, akinésie pour Parkinson ; chorée pour Huntington ; pertes de mémoire pour Alzheimer) correspondent bien à la structure lésée.
\item \cle{Proposer la piste thérapeutique} cohérente : compenser le déficit (ex. : L-DOPA, précurseur de la dopamine, dans la maladie de Parkinson) ou prise en charge symptomatique.
\end{enumerate}
\end{methode}

\begin{exoresolu}[Maladie de Parkinson : du déficit en dopamine aux symptômes]
Madame N., 68 ans, consulte à l'hôpital de district pour un tremblement des mains au repos, une rigidité des membres et une grande lenteur des mouvements (akinésie). L'autopsie de patients atteints de la même maladie a montré une dégénérescence des neurones du \emph{locus niger} (substance noire), structure des noyaux gris centraux qui sécrète normalement la \emph{dopamine}. Le traitement par la L-DOPA (précurseur de la dopamine) améliore nettement l'état des malades. 1) Identifie la structure lésée et sa fonction. 2) Explique l'origine des symptômes. 3) Justifie l'efficacité de la L-DOPA.
\tcblower
\textbf{1) Structure lésée et fonction.} La structure atteinte est le \cle{locus niger} (substance noire, \emph{substantia nigra pars compacta}), ensemble de neurones des noyaux gris centraux situés dans le mésencéphale (partie supérieure du tronc cérébral). Normalement, ces neurones libèrent la \cle{dopamine} dans le striatum (noyau caudé et putamen). La dopamine est un neuromédiateur qui facilite la \emph{voie directe} (initiation du mouvement) et inhibe la \emph{voie indirecte} (arrêt du mouvement) au sein des noyaux gris centraux, permettant ainsi le déclenchement harmonieux des mouvements volontaires.

\textbf{2) Origine des symptômes.} La dégénérescence des neurones dopaminergiques entraîne une chute de la concentration en dopamine dans le striatum. Le déséquilibre entre voies directe et indirecte conduit à une hyperactivité de la sortie inhibitrice des noyaux gris (globus pallidus interne / substance noire réticulée) vers le thalamus, bloquant l'activation du cortex moteur. D'où la triade classique : \emph{rigidité musculaire} (hypertonie plastique, « en roue dentée »), l'\emph{akinésie} (difficulté à initier le mouvement, micrographie, visage figé) et le \emph{tremblement au repos} (4-6 Hz, disparaissant à l'action). Le raisonnement causal est : neurones du locus niger détruits $\Rightarrow$ déficit en dopamine $\Rightarrow$ déséquilibre voies directes/indirectes $\Rightarrow$ inhibition excessive du thalamus $\Rightarrow$ hypoactivation du cortex moteur $\Rightarrow$ troubles moteurs.

\textbf{3) Efficacité de la L-DOPA.} La dopamine elle-même ne traverse pas la barrière hémato-encéphalique (BHE) ; son précurseur métabolique, la \cle{L-DOPA} (L-3,4-dihydroxyphénylalanine), la franchit par transport actif et est décarboxylée en dopamine par les neurones dopaminergiques résiduels (via la DOPA-décarboxylase). Le déficit est ainsi partiellement compensé, ce qui explique l'amélioration spectaculaire de l'akinésie et de la rigidité (moins efficace sur le tremblement et l'instabilité posturale tardive). C'est un traitement \emph{symptomatique} : il ne stoppe pas la dégénérescence neuronale progressive.

\textbf{Conclusion.} La maladie de Parkinson illustre parfaitement la démarche « structure $\rightarrow$ fonction $\rightarrow$ dysfonctionnement » : la destruction d'une population de neurones précise (locus niger) entraîne un déficit chimique identifiable (dopamine), responsable de symptômes moteurs caractéristiques (syndrome parkinsonien), qu'un traitement ciblé (L-DOPA) permet d'atténuer.
\end{exoresolu}

\begin{exoresolu}[Distinguer Huntington et Alzheimer]
Deux patients présentent des troubles différents. Patient 1 : mouvements involontaires, brusques et désordonnés des membres (chorée), apparaissant vers 40 ans, avec troubles du caractère ; la maladie est héréditaire et s'accompagne d'une dégénérescence des neurones des noyaux gris centraux (striatum). Patient 2 : pertes de mémoire progressives, désorientation, difficultés de langage, liées à une dégénérescence diffuse des neurones du cortex cérébral avec plaques amyloïdes. 1) Identifie la maladie de chaque patient en justifiant. 2) Explique pourquoi les symptômes du patient 1 sont moteurs alors que ceux du patient 2 sont intellectuels.
\tcblower
\textbf{1) Identification.} Le patient 1 est atteint de la \cle{maladie de Huntington} (chorée de Huntington) : les signes caractéristiques sont la chorée (mouvements involontaires saccadés, irréguliers, débordant sur les mouvements volontaires), le début typique vers l'âge mûr (30-50 ans), le caractère héréditaire à transmission autosomique dominante (gène \emph{HTT} sur chromosome 4, expansion de répétitions CAG) et l'atteinte sélective des neurones GABAergiques du \emph{striatum} (noyau caudé et putamen). Le patient 2 est atteint de la \cle{maladie d'Alzheimer} : pertes de mémoire progressives (atteinte hippocampique précoce), désorientation spatio-temporelle, aphasie, apraxie, agnosie, associées à une dégénérescence corticale diffuse (cortex d'association, hippocampe) avec dépôts extracellulaires de plaques amyloïdes (peptide $\beta$-amyloïde) et enchevêtrements neurofibrillaires intracellulaires (protéine Tau hyperphosphorylée).

\textbf{2) Différence des symptômes.} Les symptômes dépendent de la \emph{localisation} de la dégénérescence, conformément au principe « structure $\rightarrow$ fonction ». Chez le patient 1, les neurones détruits appartiennent aux \emph{noyaux gris centraux} (striatum), structures clés de la régulation et de la sélection des programmes moteurs : leur destruction (perte de l'inhibition GABAergique) libère des mouvements involontaires incontrôlés (chorée) et perturbe le comportement. Chez le patient 2, ce sont les neurones du \emph{cortex cérébral} (siège des fonctions cognitives supérieures : mémoire épisodique dans l'hippocampe, langage dans les aires de Broca/Wernicke, fonctions exécutives dans le cortex préfrontal) qui dégénèrent : les symptômes sont donc cognitifs (démence de type cortical), sans trouble moteur initial.

\textbf{Conclusion.} Parkinson, Huntington et Alzheimer sont trois maladies neurodégénératives distinctes : même mécanisme général (mort neuronale progressive, agrégats protéiques : $\alpha$-synucléine pour Parkinson, Huntingtine mutée pour Huntington, $\beta$-amyloïde/Tau pour Alzheimer), mais localisations anatomiques différentes, donc symptômes cliniques différents — troubles moteurs hypokinétiques pour Parkinson, hyperkinétiques pour Huntington, troubles cognitifs pour Alzheimer.
\end{exoresolu}

\courssub{Synthèse : Relier structure, fonction et dysfonctionnement ; Prévention}

\begin{aretenir}[Tableau récapitulatif des trois maladies neurodégénératives]
\begin{center}
\begin{tabularx}{\linewidth}{|l|X|X|X|}
\hline
\rowcolor{popblueL} \textbf{Critère} & \textbf{Maladie de Parkinson} & \textbf{Maladie de Huntington} & \textbf{Maladie d'Alzheimer} \\
\hline
\cle{Structure atteinte} & Locus niger (Substantia nigra pars compacta) : neurones dopaminergiques & Striatum (Noyau caudé, Putamen) : neurones GABAergiques & Cortex cérébral diffus + Hippocampe : neurones cholinergiques / glutamatergiques \\
\hline
\cle{Anomalie biochimique} & Déficit en \cle{Dopamine} (voie nigro-striatale) & Déficit en \cle{GABA} / Acétylcholine (voie striato-nigrale) & Déficit en \cle{Acétylcholine} (voie septo-hippocampique) + Dépôts $\beta$-amyloïde, Tau \\
\hline
\cle{Type de trouble moteur} & \textbf{Hypokinésie} : Akinésie, Rigidité, Tremblement au repos (4-6 Hz) & \textbf{Hyperkinésie} : \cle{Chorée} (mouvements involontaires, choréo-athétose) & Aucun trouble moteur initial (stade tardif : marche instable) \\
\hline
\cle{Troubles associés} & Dépression, troubles autonomes, démence (stade avancé) & Troubles psychiatriques (dépression, irritabilité), démence subcorticale & \cle{Démence corticale} : Amnésie, Aphasie, Apraxie, Agnosie, troubles exécutifs \\
\hline
\cle{Transmission} & Majoritairement sporadique (facteurs génétiques + environnementaux) & \textbf{Autosomique dominante} (gène \emph{HTT}, chr 4, expansion CAG) & Majoritairement sporadique (facteur de risque : allèle \emph{APOE} $\varepsilon$4) ; formes familiales rares (APP, PSEN1/2) \\
\hline
\cle{Traitement principal} & \cle{L-DOPA} (précurseur dopamine) + Agonistes dopaminergiques + Inhibiteurs MAO-B/COMT & Symptomatique : Neuroleptiques (antidopaminergiques), Tétrabénazine (dépleteur vésiculaire) & Inhibiteurs de l'acétylcholinestérase (Donepezil, Rivastigmine) + Antagoniste NMDA (Mémantine) \\
\hline
\end{tabularx}
\end{center}
\end{aretenir}

\begin{popfigure}
\begin{center}
\begin{tikzpicture}[node distance=1.2cm and 1.5cm,
sbox/.style={draw=popblue, fill=popblueL, rounded corners=3pt, align=center, font=\small, inner sep=4pt, minimum width=3cm},
dbox/.style={draw=poporange, fill=poporangeL, rounded corners=3pt, align=center, font=\small, inner sep=4pt, minimum width=3.5cm},
sbox2/.style={draw=popgreen, fill=popgreenL, rounded corners=3pt, align=center, font=\small, inner sep=4pt, minimum width=3cm},
arr/.style={-{Stealth[length=2.5mm]}, thick, popdark}]
% Ligne Structure
\node[sbox, align=center] (s1){STRUCTURE\\ Locus niger (SNc)};
\node[sbox, right=of s1, align=center] (s2){STRUCTURE\\ Striatum (Caudé/Putamen)};
\node[sbox, right=of s2, align=center] (s3){STRUCTURE\\ Cortex / Hippocampe};
% Ligne Fonction
\node[dbox, below=of s1, align=center] (f1){FONCTION\\ Synthèse Dopamine\\ $\rightarrow$ Régulation motrice (Voie directe/indirecte)};
\node[dbox, below=of s2, align=center] (f2){FONCTION\\ Intégration GABA\\ $\rightarrow$ Sélection / Inhibition mouvements};
\node[dbox, below=of s3, align=center] (f3){FONCTION\\ Traitement cognitif\\ Mémoire, Langage, Raisonnement};
% Ligne Dysfonctionnement
\node[sbox2, below=of f1, align=center] (d1){DYSFONCTION\\ Parkinson\\ Akinésie, Rigidité, Tremblement};
\node[sbox2, below=of f2, align=center] (d2){DYSFONCTION\\ Huntington\\ Chorée, Troubles psy, Démence};
\node[sbox2, below=of f3, align=center] (d3){DYSFONCTION\\ Alzheimer\\ Amnésie, Aphasie, Désorientation};
% Flèches
\draw[arr] (s1) -- (f1);
\draw[arr] (f1) -- (d1);
\draw[arr] (s2) -- (f2);
\draw[arr] (f2) -- (d2);
\draw[arr] (s3) -- (f3);
\draw[arr] (f3) -- (d3);
% Étiquettes colonnes
\node[above=0.2cm of s1, font=\bfseries, popblue] {PARKINSON};
\node[above=0.2cm of s2, font=\bfseries, poporange] {HUNTINGTON};
\node[above=0.2cm of s3, font=\bfseries, popgreen] {ALZHEIMER};
\end{tikzpicture}
\end{center}
\legende{Modèle unificateur « Structure $\rightarrow$ Fonction $\rightarrow$ Dysfonctionnement » pour les trois maladies neurodégénératives. La spécificité des symptômes découle strictement de la topographie de la lésion neuronale.}
\end{popfigure}

\begin{savaistu}[Contexte camerounais et Prévention]
Au Cameroun, le vieillissement de la population entraîne une augmentation des maladies neurodégénératives. Le \cle{Service de Neurologie} de l'Hôpital Central de Yaoundé et de l'Hôpital Laquintinie de Douala sont les centres de référence pour le diagnostic et la prise en charge (L-DOPA disponible via pharmacies hospitalières, mais coût élevé pour les patients non assurés).
\begin{itemize}
\item \textbf{Maladie de Parkinson :} Prévalence estimée autour de 100-150/100 000 habitants. Pas de prévention primaire connue. L'activité physique régulière et l'évitement des pesticides (exposition agricole dans les zones cacaoyères/cotonnières) sont des pistes de réduction du risque environnemental.
\item \textbf{Maladie de Huntington :} Maladie génétique rare. Le \cle{conseil génétique} est crucial pour les familles porteuses de la mutation (test prédictif possible à l'âge adulte avec accompagnement psychologique). Pas de prévention de l'expression, mais diagnostic prénatal/préimplantatoire possible.
\item \textbf{Maladie d'Alzheimer :} Première cause de démence. Facteurs de risque modifiables (Commission Lancet 2020) : \cle{hypertension artérielle}, \cle{diabète}, \cle{obésité}, \cle{tabagisme}, \cle{inactivité physique}, \cle{isolement social}, \cle{déficit auditif}, \cle{niveau d'éducation faible}. La prévention passe par la lutte contre les maladies cardiovasculaires (très prévalentes au Cameroun), la stimulation cognitive (lecture, jeux, vie sociale) et la scolarisation des filles.
\end{itemize}
\end{savaistu}

\begin{attention}[Les 5 erreurs qui coûtent des points au Bac]
\begin{enumerate}
\item \cle{Confondre aire motrice et aire sensitive} : l'aire motrice (gyrus précentral, \emph{devant} le sillon central) \emph{commande} les mouvements ; l'aire sensitive (gyrus postcentral, \emph{derrière}) \emph{reçoit} les sensations. Les inverser dans une interprétation de stimulation est une faute éliminatoire.
\item \cle{Oublier le croisement des voies} : le cortex moteur gauche commande le côté \emph{droit} du corps (décussation des pyramides au bulbe rachidien). Écrire qu'une lésion du cortex gauche paralyse le côté gauche est faux.
\item \cle{Attribuer au cervelet le déclenchement du mouvement} : le cervelet \emph{coordonne} et \emph{corrige} (comparateur) ; c'est le cortex moteur qui \emph{décide} et \emph{déclenche} (initiateur). Ne pas écrire « le cervelet commande le mouvement volontaire ».
\item \cle{Confondre les trois maladies} : Parkinson = déficit en dopamine (locus niger), tremblement au repos + rigidité + akinésie (hypokinésie) ; Huntington = chorée héréditaire (striatum), mouvements involontaires (hyperkinésie) ; Alzheimer = démence corticale (cortex/hippocampe), plaques amyloïdes. Ne jamais écrire que la L-DOPA \emph{guérit} Parkinson : elle ne fait que compenser le déficit de façon symptomatique.
\item \cle{Négliger le schéma} : un schéma anatomique ou fonctionnel sans titre, sans légendes horizontales, ou avec des flèches de légende qui se croisent, perd des points même si le contenu est correct. Dans un schéma fonctionnel, le sens des flèches doit aller du centre nerveux vers l'effecteur (voie descendante) et le retour vers le centre (voie ascendante).
\end{enumerate}
\end{attention}

\begin{exercice}[Entraînement Bac - Cas cliniques et Schéma]
\textbf{Exercice 1 (Schéma) :} Réalise un schéma annoté de la face latérale de l'hémisphère gauche humain, montrant le sillon central, l'aire motrice, l'aire sensitive, et la représentation somatotopique (pied, main, visage) pour chacune. Indique par une flèche le côté du corps contrôlé.

\textbf{Exercice 2 (Analyse de documents) :} On pratique une micro-injection de traceur fluorescent (anterograde) dans l'aire motrice gauche d'un macaque. Après survie, on observe la fluorescence dans la corne ventrale de la moelle épinière.
\begin{enumerate}
\item De quel côté (gauche ou droit) la fluorescence est-elle observée ? Justifie.
\item Nomme le faisceau de substance blanche emprunté par les axones traceurs.
\item Quel type de synapse est réalisé dans la corne ventrale ? Sur quel type de neurone ?
\end{enumerate}

\textbf{Exercice 3 (Rédaction) :} Un patient de 55 ans présente depuis 2 ans des mouvements choréiques, une irritabilité et des troubles de la mémoire. Son père est décédé d'une maladie neurologique à 58 ans. L'IRM montre une atrophie des noyaux caudés.
\begin{enumerate}
\item Quel est le diagnostic probable ? Justifie par trois arguments.
\item Précise la transmission génétique et le chromosome concerné.
\item Explique le mécanisme physiopathologique reliant l'atteinte du striatum à la chorée.
\end{enumerate}
\end{exercice}

\begin{corrige}[Exercice 1]
Schéma attendu : Hémisphère gauche vu de profil. Sillon central (Rolando) vertical. Antérieur : Gyrus précentral (Aire motrice 4) avec homonculus : Pied (haut/médial) $\rightarrow$ Jambe $\rightarrow$ Tronc $\rightarrow$ Bras $\rightarrow$ Main (large) $\rightarrow$ Visage/Langue (bas/latéral). Postérieur : Gyrus postcentral (Aire sensitive 1,3,2) avec homonculus sensitif miroir. Flèche partant de l'aire motrice vers la mention « Côté DROIT du corps (controlatéral) ». Titre : « Aire motrice et sensitive de l'hémisphère gauche - Organisation somatotopique ». Légendes horizontales, traits de rappel non croisés.
\end{corrige}

\begin{corrige}[Exercice 2]
\begin{enumerate}
\item \textbf{Côté DROIT.} Les axones des neurones pyramidaux de l'aire motrice gauche descendent dans la voie pyramidale, se croisent (décussation) au niveau du bulbe rachidien (décussation des pyramides) et descendent controlatéralement dans le faisceau cortico-spinal latéral droit pour faire relais sur les motoneurones de la corne ventrale droite.
\item \textbf{Faisceau cortico-spinal latéral} (ou voie pyramidale croisée / faisceau pyramidal).
\item \textbf{Synapse excitatrice (cholinergique, nicotinique)} sur le \textbf{motoneurone $\alpha$} (motoneurone inférieur / final commun) dont le corps cellulaire est dans la corne ventrale de la moelle épinière.
\end{enumerate}
\end{corrige}

\begin{corrige}[Exercice 3]
\begin{enumerate}
\item \textbf{Maladie de Huntington (Chorée de Huntington).} Arguments : 1) Symptôme cardinal : chorée (mouvements involontaires, brusques, désordonnés). 2) Atteinte des noyaux gris centraux confirmée par l'IRM (atrophie des noyaux caudés = striatum). 3) Caractère héréditaire familial (père atteint) compatible avec une transmission autosomique dominante.
\item \textbf{Transmission autosomique dominante.} Gène \emph{HTT} (Huntingtine) localisé sur le \textbf{chromosome 4} (bras court 4p16.3). Mutation : expansion anormale de répétitions CAG ($>$ 36 répétitions, souvent 40-50).
\item \textbf{Mécanisme :} Le striatum (noyau caudé + putamen) est la porte d'entrée des noyaux gris centraux. Ses neurones de projection (GABAergiques) forment deux voies : la voie directe (facilitatrice du mouvement, récepteurs D1) et la voie indirecte (inhibitrice du mouvement, récepteurs D2). Dans Huntington, la dégénérescence touche \emph{préférentiellement les neurones de la voie indirecte} (projetant vers le globus pallidus externe). Cela entraîne une \textbf{désinhibition du thalamus} (voie indirecte moins active $\rightarrow$ GPe moins inhibé $\rightarrow$ STN/GPi moins excités $\rightarrow$ Thalamus moins inhibé) et une hyperactivité du cortex moteur $\rightarrow$ \cle{chorée} (mouvements involontaires non supprimés).
\end{enumerate}
\end{corrige}

\newpage

\coursec{Exercices}

\courssub{Je vérifie mes connaissances}

\begin{exercice}[QCM : les aires encéphaliques]
Pour chaque question, une seule réponse est exacte. Recopie la lettre correspondante.
\begin{enumerate}
\item L'aire motrice primaire (aire 4 de Brodmann) est située dans :
\begin{enumerate}
\item[a)] le lobe occipital ;
\item[b)] le lobe pariétal, en arrière du sillon central ;
\item[c)] le lobe frontal, en avant du sillon central (circonvolution frontale ascendante) ;
\item[d)] le lobe temporal.
\end{enumerate}
\item La décussation des pyramides a lieu au niveau :
\begin{enumerate}
\item[a)] du thalamus ;
\item[b)] du bulbe rachidien ;
\item[c)] du cervelet ;
\item[d)] du corps calleux.
\end{enumerate}
\item La maladie de Parkinson est due à la dégénérescence des neurones dopaminergiques :
\begin{enumerate}
\item[a)] de la substance noire (locus niger) ;
\item[b)] de l'hippocampe ;
\item[c)] du cortex visuel ;
\item[d)] du nerf vague.
\end{enumerate}
\item La chorée de Huntington se transmet selon un mode :
\begin{enumerate}
\item[a)] autosomique récessif ;
\item[b)] autosomique dominant ;
\item[c)] lié au chromosome X récessif ;
\item[d)] mitochondrial.
\end{enumerate}
\end{enumerate}
\end{exercice}

\begin{exercice}[Vrai ou faux — justifier]
Indique si chacune des affirmations suivantes est vraie ou fausse, puis justifie brièvement ta réponse.
\begin{enumerate}
\item Un mouvement volontaire est un mouvement déclenché sans participation de la volonté, en réponse à un stimulus.
\item La stimulation électrique d'un point précis du cortex moteur droit provoque la contraction de muscles du côté gauche du corps.
\item Le cervelet est le centre de commande qui initie le mouvement volontaire.
\item La L-DOPA, précurseur de la dopamine, franchit la barrière hémato-encéphalique, contrairement à la dopamine elle-même.
\item Dans la maladie d'Alzheimer, les premiers troubles sont essentiellement moteurs.
\end{enumerate}
\end{exercice}

\begin{exercice}[Définitions]
Définis en une ou deux phrases chacun des termes suivants :
\begin{enumerate}
\item mouvement volontaire ;
\item aire motrice primaire ;
\item homonculus moteur ;
\item akinésie ;
\item tremblement de repos.
\end{enumerate}
\end{exercice}

\begin{exercice}[Calcul direct : risque génétique]
Dans une famille, un homme atteint de la maladie de Huntington (allèle dominant $H$ ; allèle normal $h$) est hétérozygote. Il épouse une femme saine.
\begin{enumerate}
\item Écris le génotype de chaque parent.
\item Réalise l'échiquier de croisement.
\item Calcule, en pourcentage, la probabilité qu'un enfant de ce couple développe la maladie.
\end{enumerate}
\end{exercice}

\courssub{Je m'entraîne}

\begin{exercice}[Stimulations corticales chez le singe]
Au cours d'une expérience, on stimule électriquement différents points du cortex cérébral d'un singe anesthésié et l'on note les réponses observées.

\begin{center}
\begin{tabularx}{\linewidth}{|c|X|X|}
\hline
\rowcolor{popblueL} Point stimulé & Localisation corticale & Réponse observée \\
\hline
A & Circonvolution frontale ascendante droite, partie supérieure & Mouvement de la patte postérieure gauche \\
\hline
B & Circonvolution frontale ascendante droite, partie moyenne & Mouvement des doigts de la main gauche \\
\hline
C & Circonvolution frontale ascendante droite, partie inférieure & Mouvement des lèvres et de la langue (côté gauche) \\
\hline
D & Circonvolution pariétale ascendante (en arrière du sillon central) & Aucun mouvement ; le singe réagit comme s'il était touché \\
\hline
\end{tabularx}
\end{center}

\begin{enumerate}
\item Que peut-on déduire de la comparaison des résultats obtenus aux points A, B et C quant à l'organisation de l'aire motrice ?
\item Pourquoi les réponses aux points A, B, C s'observent-elles du côté gauche du corps alors que la stimulation est appliquée à droite ?
\item Quelle est la fonction de la zone stimulée au point D ? Justifie.
\end{enumerate}
\end{exercice}

\begin{exercice}[Lésions expérimentales : cortex moteur et cervelet]
On réalise deux séries d'expériences.
\begin{itemize}
\item \textbf{Expérience 1} : chez un chien, on sectionne les voies pyramidales au niveau du bulbe rachidien. L'animal ne peut plus exécuter de mouvements volontaires précis, mais conserve sa posture et sa démarche globale.
\item \textbf{Expérience 2} : chez un pigeon, on ablate le cervelet. L'animal est incapable de voler correctement, de se poser sans chuter et de picorer avec précision, bien que ses muscles et ses nerfs soient intacts.
\end{itemize}
\begin{enumerate}
\item À partir de l'expérience 1, déduis le rôle des voies pyramidales.
\item À partir de l'expérience 2, déduis le rôle du cervelet dans la motricité.
\item En comparant les deux expériences, distingue les rôles respectifs du cortex moteur et du cervelet dans le mouvement volontaire.
\end{enumerate}
\end{exercice}

\begin{exercice}[Schéma fonctionnel du mouvement volontaire]
Un élève attrape un stylo posé sur sa table.
\begin{enumerate}
\item Réalise le schéma fonctionnel du trajet suivi par le message nerveux, depuis la décision du mouvement jusqu'à la contraction des muscles de la main, en faisant figurer : aire motrice primaire, voie pyramidale, décussation des pyramides, moelle épinière, motoneurone, muscle effecteur.
\item Place sur ce schéma la boucle de contrôle cérébelleuse et précise son rôle.
\item Indique sur le schéma le siège de la lésion responsable d'une paralysie de la main droite, si celle-ci résulte d'une atteinte corticale.
\end{enumerate}
\end{exercice}

\begin{exercice}[Comparaison des trois maladies]
Recopie et complète le tableau comparatif suivant :

\begin{center}
\begin{tabularx}{\linewidth}{|l|X|X|X|}
\hline
\rowcolor{popblueL} Critère & Parkinson & Huntington & Alzheimer \\
\hline
Structure atteinte & & & \\
\hline
Symptômes principaux & & & \\
\hline
Origine (héréditaire ou non) & & & \\
\hline
Traitement symptomatique & & & \\
\hline
\end{tabularx}
\end{center}

Cite ensuite deux mesures permettant de limiter les lésions des structures nerveuses du contrôle de la motricité.
\end{exercice}

\courssub{Je me prépare au Bac}

\begin{exercice}[Type Bac : l'aire motrice du cortex]
\textbf{Document 1.} Chez un macaque éveillé, on enregistre l'activité électrique de neurones de la circonvolution frontale ascendante pendant que l'animal exécute un mouvement volontaire de préhension. On constate que certains neurones déchargent des potentiels d'action \textbf{avant} le début du mouvement de la main.

\textbf{Document 2.} On stimule électriquement des points du cortex d'un hémisphère cérébral humain au cours d'une intervention chirurgicale (patient éveillé, anesthésie locale). Les résultats sont consignés ci-dessous.

\begin{center}
\begin{tabularx}{\linewidth}{|c|X|X|}
\hline
\rowcolor{popblueL} Stimulation & Site cortical & Effet observé \\
\hline
1 & Frontale ascendante, tiers supérieur & Flexion de la jambe controlatérale \\
\hline
2 & Frontale ascendante, tiers moyen & Serrage des doigts controlatéraux \\
\hline
3 & Frontale ascendante, tiers inférieur & Mouvement de la face et de la langue controlatérales \\
\hline
4 & Pariétale ascendante & Sensation de toucher au niveau de la main controlatérale \\
\hline
5 & Lobe occipital & Perception d'éclairs lumineux \\
\hline
\end{tabularx}
\end{center}

\textbf{Document 3.} La surface corticale dédiée à la main et à la face est beaucoup plus grande que celle dédiée au tronc ou à la cuisse.

\begin{enumerate}
\item Définis le mouvement volontaire. (1 pt)
\item À partir du document 1, montre que la circonvolution frontale ascendante est le siège de la commande motrice et non d'une simple réaction au mouvement. (2 pts)
\item Exploite le document 2 pour dégager deux caractéristiques de l'organisation de l'aire motrice primaire (somatotopie et latéralisation). (3 pts)
\item Interprète le document 3 : à quoi correspond la disproportion des surfaces corticales ? (2 pts)
\item Réalise le schéma annoté d'un hémisphère cérébral en vue latérale, en y plaçant : sillon central, aire motrice primaire, aire sensitive primaire, lobe frontal, lobe pariétal, lobe occipital, lobe temporal. (2 pts)
\end{enumerate}
\end{exercice}

\begin{exercice}[Type Bac : étude de cas clinique à l'hôpital de Yaoundé]
Monsieur N., 62 ans, consulte à l'hôpital central de Yaoundé pour des tremblements de la main droite apparaissant au repos et disparaissant pendant le mouvement volontaire. Le médecin note également une rigidité des membres et une lenteur des mouvements (akinésie), ainsi qu'une marche à petits pas.

\textbf{Document 1.} Le dosage de la dopamine dans la substance noire (locus niger) donne les résultats suivants :

\begin{center}
\begin{tabular}{|l|c|c|}
\hline
\rowcolor{popblueL} & Témoin sain & Patient N. \\
\hline
Dopamine (ng/mg de tissu) & 3,2 & 0,6 \\
\hline
\end{tabular}
\end{center}

\textbf{Document 2.} La substance noire contient normalement des neurones dopaminergiques qui projettent sur le striatum, structure impliquée dans le déclenchement et l'amplitude des mouvements volontaires.

\textbf{Document 3.} Le médecin prescrit de la L-DOPA. Après quelques semaines de traitement, les tremblements et la rigidité du patient diminuent nettement. La dopamine administrée directement serait inefficace car elle ne franchit pas la barrière hémato-encéphalique, contrairement à la L-DOPA.

\begin{enumerate}
\item Identifie la maladie dont souffre Monsieur N. en justifiant à partir de ses symptômes. (2 pts)
\item À l'aide des documents 1 et 2, précise la structure atteinte et la nature du dysfonctionnement biochimique. (3 pts)
\item Explique le mode d'action de la L-DOPA et pourquoi on ne peut pas administrer la dopamine directement. (3 pts)
\item La maladie de Parkinson est-elle héréditaire dans la majorité des cas ? Cite une autre maladie neurodégénérative qui, elle, se transmet héréditairement, et précise son mode de transmission. (2 pts)
\end{enumerate}
\end{exercice}

\begin{exercice}[Type Bac : la chorée de Huntington dans une famille camerounaise]
Dans une famille suivie au centre hospitalier de Bafoussam, on établit l'arbre généalogique ci-dessous. La maladie se déclare généralement entre 35 et 50 ans : mouvements anormaux involontaires (chorée), troubles du comportement, déclin intellectuel progressif.

\begin{center}
\begin{tikzpicture}[
  femme/.style={circle,draw,minimum size=0.55cm,inner sep=1pt},
  homme/.style={rectangle,draw,minimum size=0.55cm,inner sep=1pt},
  atteint/.style={fill=popdark},
  every node/.style={font=\small}]
% Génération I
\node[femme,atteint] (I1) at (0,0) {};
\node[homme] (I2) at (1.4,0) {};
\draw (I1) -- (I2);
\node[below=2pt of I1] {I-1};
\node[below=2pt of I2] {I-2};
% Descendance
\draw (0.7,0) -- (0.7,-0.6) -- (0.7,-0.7);
\draw (0.7,-0.7) -- ++(-1.6,0) -- ++(3.2,0);
% Génération II
\node[femme] (II1) at (-0.9,-1.5) {};
\node[femme,atteint] (II2) at (0.7,-1.5) {};
\node[homme,atteint] (II3) at (2.3,-1.5) {};
\draw (-0.9,-0.7) -- (II1);
\draw (0.7,-0.7) -- (II2);
\draw (2.3,-0.7) -- (II3);
\node[below=2pt of II1] {II-1};
\node[below=2pt of II2] {II-2};
\node[below=2pt of II3] {II-3};
% Union II-2
\node[homme] (IIp) at (1.9,-1.5) {};
\draw (II2) -- (IIp);
\draw (1.3,-1.5) -- (1.3,-2.1);
\node[femme] (III1) at (1.3,-2.7) {};
\draw (1.3,-2.1) -- (III1);
\node[below=2pt of III1] {III-1 (20 ans)};
\end{tikzpicture}
\end{center}

Légende : cercle = femme ; carré = homme ; symbole rempli = individu atteint.

\begin{enumerate}
\item À partir de l'arbre généalogique, montre que l'allèle responsable de la maladie est dominant et porté par un autosome. (3 pts)
\item En notant $H$ l'allèle morbide et $h$ l'allèle normal, donne les génotypes de I-1, I-2 et II-2. Justifie. (3 pts)
\item L'individu III-1, âgé de 20 ans et encore asymptomatique, souhaite connaître son risque. Calcule la probabilité qu'elle ait reçu l'allèle $H$ de sa mère II-2. (2 pts)
\item Compare la chorée de Huntington et la maladie de Parkinson : structure atteinte, nature des troubles moteurs, origine génétique. (2 pts)
\end{enumerate}
\end{exercice}

\newpage

\coursec{Corrigés des exercices}

\coursec{Corrigés des exercices d'application — Leçon 10 : Limitation des dysfonctionnements des structures responsables du contrôle de la motricité}

\begin{corrige}[Exercice 1 — QCM : les aires encéphaliques]
\begin{enumerate}
\item \textbf{Réponse c}. L'aire motrice primaire (aire 4 de Brodmann) occupe la circonvolution frontale ascendante, dans le lobe frontal, immédiatement \textbf{en avant} du sillon central (scissure de Rolando). Le lobe occipital abrite l'aire visuelle, le lobe temporal l'aire auditive, et la région en arrière du sillon central est l'aire \emph{sensitive} primaire (somesthésique).
\item \textbf{Réponse b}. La décussation des pyramides (croisement des fibres de la voie pyramidale) a lieu au niveau du \textbf{bulbe rachidien} (moelle allongée), à la jonction entre l'encéphale et la moelle épinière. C'est ce croisement qui explique le contrôle controlatéral du corps.
\item \textbf{Réponse a}. La maladie de Parkinson résulte de la dégénérescence des neurones dopaminergiques de la \textbf{substance noire} (locus niger), qui projettent normalement sur le striatum (noyau caudé et putamen).
\item \textbf{Réponse b}. La chorée de Huntington se transmet selon un mode \textbf{autosomique dominant} : l'allèle morbide est porté par un chromosome non sexuel (chromosome 4, gène \emph{HTT}) et un seul exemplaire suffit à provoquer la maladie.
\end{enumerate}
\end{corrige}

\begin{corrige}[Exercice 2 — Vrai ou faux]
\begin{enumerate}
\item \textbf{Faux.} Un mouvement volontaire est précisément un mouvement \textbf{déclenché par la volonté}, sous le contrôle du cortex cérébral (aire motrice primaire), et non une réponse automatique à un stimulus. La définition donnée correspond à un \emph{réflexe} (mouvement involontaire, stéréotypé, sans conscience).
\item \textbf{Vrai.} En raison de la \textbf{décussation des pyramides} au niveau du bulbe rachidien, chaque hémisphère cérébral commande la motricité de la moitié \textbf{opposée} (controlatérale) du corps : le cortex moteur droit commande le côté gauche, et vice-versa.
\item \textbf{Faux.} Le centre qui \emph{initie} et commande le mouvement volontaire est l'\textbf{aire motrice primaire du cortex frontal} (circonvolution frontale ascendante). Le cervelet est un centre de \emph{coordination} : il compare le programme moteur élaboré par le cortex aux informations sensitives de retour (proprioception) et corrige le mouvement en cours (précision, équilibre, tonicité), sans jamais l'initier.
\item \textbf{Vrai.} La dopamine ne franchit pas la barrière hémato-encéphalique (BHE) ; son précurseur métabolique, la \textbf{L-DOPA} (L-3,4-dihydroxyphénylalanine), la franchit par transport actif, puis est transformée en dopamine dans l'encéphale par la DOPA-décarboxylase. C'est le principe du traitement symptomatique de référence de la maladie de Parkinson.
\item \textbf{Faux.} Dans la maladie d'Alzheimer, les premiers troubles sont \textbf{cognitifs} : pertes de mémoire récente (atteinte hippocampique), désorientation spatio-temporelle, troubles du langage (aphasie) et du jugement, liés à la dégénérescence des neurones de l'hippocampe et du cortex d'association. Les troubles moteurs (marche, déglutition) n'apparaissent qu'à des stades tardifs et sévères.
\end{enumerate}
\end{corrige}

\begin{corrige}[Exercice 3 — Définitions]
\begin{enumerate}
\item \textbf{Mouvement volontaire} : mouvement décidé consciemment et commandé par l'aire motrice primaire du cortex cérébral, transmis aux muscles squelettiques par la voie pyramidale (faisceau cortico-spinal) et coordonné par le cervelet ; il est sous le contrôle de la conscience et de la volonté.
\item \textbf{Aire motrice primaire} : région du cortex cérébral située dans la circonvolution frontale ascendante (aire 4 de Brodmann), en avant du sillon central (de Rolando), dont les neurones pyramidaux (cellules de Betz) émettent les messages nerveux moteurs commandant les mouvements volontaires de la moitié controlatérale du corps.
\item \textbf{Homonculus moteur} : représentation somatotopique, le long de l'aire motrice primaire, des différentes parties du corps sous forme d'un personnage dont les proportions reflètent la surface corticale dédiée à chaque région ; la main, les doigts et la face (langue, lèvres) y occupent une place disproportionnée, proportionnelle à la finesse (précision) de leur contrôle moteur et non à leur taille anatomique.
\item \textbf{Akinésie} : difficulté ou impossibilité d'initier et d'exécuter les mouvements volontaires, se traduisant par une lenteur (bradykinésie) et une pauvreté des gestes (hypokinésie), une mimique figée et une marche à petits pas ; symptôme majeur de la maladie de Parkinson.
\item \textbf{Tremblement de repos} : tremblement involontaire, rythmique et régulier (4--6 Hz) des extrémités (souvent les mains, type « roulage de pilule » pouce-index) apparaissant lorsque le membre est au repos complet, s'atténuant pendant le mouvement volontaire et disparaissant pendant le sommeil ; signe cardinal de la maladie de Parkinson.
\end{enumerate}
\end{corrige}

\begin{corrige}[Exercice 4 — Calcul direct : risque génétique (Huntington)]
\begin{enumerate}
\item La maladie est à transmission autosomique dominante. Un individu atteint est le plus souvent hétérozygote (l'homozygotie $H/H$ étant rare et sévère) : génotype $H/h$. La femme saine ne possède que l'allèle normal : $h/h$.
\[ \text{Père (atteint) : } H/h \qquad \text{Mère (saine) : } h/h \]
\item Échiquier de croisement (gamètes paternels en lignes, maternels en colonnes) :
\begin{center}
\begin{tabular}{|c|c|c|}
\hline
\rowcolor{popblueL} Gamètes $\backslash$ & $h$ & $h$ \\
\hline
$H$ & $H/h$ (atteint) & $H/h$ (atteint) \\
\hline
$h$ & $h/h$ (sain) & $h/h$ (sain) \\
\hline
\end{tabular}
\end{center}
\item Sur les quatre combinaisons équiprobables, deux correspondent au génotype $H/h$ (enfant qui développera la maladie) et deux au génotype $h/h$ (enfant sain) :
\[ P(\text{enfant atteint}) = \dfrac{2}{4} = \dfrac{1}{2} = 50\ \%. \]
\textbf{Probabilité : 50 \%} pour chaque grossesse, indépendamment du sexe de l'enfant (transmission autosomique).
\end{enumerate}
\end{corrige}

\begin{corrige}[Exercice 5 — Stimulations corticales chez le singe (expérience de Fritsch et Hitzig / Penfield)]
\begin{enumerate}
\item Les points de stimulation A, B et C appartiennent tous à la circonvolution frontale ascendante (en avant du sillon central) : c'est donc l'\textbf{aire motrice primaire}. De plus, à chaque point stimulé correspond le mouvement d'une partie du corps bien précise et constante : la patte postérieure pour le point supérieur (A), la main pour le point médian (B), la face pour le point inférieur (C). On en déduit que l'aire motrice est organisée de manière \textbf{somatotopique} : il existe une carte ordonnée du corps le long du cortex moteur (homonculus moteur), chaque région corticale commandant une région corporelle déterminée.
\item Les réponses motrices s'observent du côté \textbf{gauche} alors que la stimulation est appliquée sur l'hémisphère \textbf{droit} parce que les fibres de la voie pyramidale issues du cortex moteur \textbf{se croisent} (décussation) au niveau du bulbe rachidien avant de descendre dans la moelle épinière. Chaque hémisphère commande donc la motricité de la moitié \textbf{controlatérale} du corps.
\item La zone D, située en arrière du sillon central (circonvolution pariétale ascendante), est l'\textbf{aire sensitive primaire} (somesthésique, aire 3-1-2 de Brodmann) : sa stimulation ne provoque aucun mouvement mais une \emph{sensation} (ici, l'impression d'être touché au niveau de la main gauche). Elle reçoit et analyse les informations sensitives (tact, douleur, température, proprioception) provenant de la moitié controlatérale du corps.
\end{enumerate}
\end{corrige}

\begin{corrige}[Exercice 6 — Lésions expérimentales : cortex moteur et cervelet]
\begin{enumerate}
\item \textbf{Expérience 1 (section des voies pyramidales) :} L'ablation des mouvements volontaires précis (préension, mouvements fins) avec conservation de la posture et des mouvements automatiques (marche, griffage) montre que les voies pyramidales (faisceaux cortico-spinaux) sont les \textbf{voies de conduction spécifiques du message moteur volontaire} : elles transmettent l'influx nerveux de commande depuis l'aire motrice corticale vers les motoneurones de la corne antérieure de la moelle épinière.
\item \textbf{Expérience 2 (ablation du cervelet) :} Les muscles et les nerfs périphériques étant intacts, le pigeon conserve la capacité de contracter ses muscles, mais ses mouvements deviennent maladroits, imprécis, déséquilibrés (impossibilité de voler, chutes à l'atterrissage, picorage à côté du grain). Le cervelet est donc le centre de \textbf{coordination des mouvements volontaires} : il assure la précision, l'harmonie des gestes, l'équilibre et le maintien du tonus musculaire, en corrigeant en permanence le programme moteur grâce aux informations sensitives de retour (proprioception, vestibulaire, visuelle).
\item \textbf{Comparaison fonctionnelle} : Le \textbf{cortex moteur} (via la voie pyramidale) \emph{initie, programme et commande} le mouvement volontaire — sans lui, le mouvement volontaire précis disparaît (paralysie/parésie). Le \textbf{cervelet} ne commande pas le mouvement mais le \emph{contrôle, le corrige et l'harmonise} en temps réel — sans lui, le mouvement subsiste mais devient ataxique (maladroit, désordonné, hypermétrie). Le cortex est le « chef d'orchestre » qui donne la partition, le cervelet est le « correcteur » qui assure l'exécution juste.
\end{enumerate}
\end{corrige}

\begin{corrige}[Exercice 7 — Schéma fonctionnel du mouvement volontaire]
\begin{enumerate}
\item \textbf{Schéma fonctionnel} (les numéros indiquent l'ordre chronologique du trajet de l'influx nerveux) :
\begin{popfigure}
\begin{tikzpicture}[node distance=0.35cm and 0.9cm,
  boite/.style={draw=popblue, fill=popblueL, rounded corners=2pt, align=center, font=\small, inner sep=3pt},
  fleche/.style={-{Stealth[length=2.5mm]}, thick, popdark}]
\node[boite, align=center] (cortex){1. Aire motrice primaire\\ (cortex frontal gauche)\\ \emph{Décision, programmation,  commande}};
\node[boite, below=of cortex, align=center] (pyr){2. Voie pyramidale\\ (faisceau cortico-spinal descendant)};
\node[boite, below=of pyr, align=center] (dec){3. Décussation des pyramides\\ (bulbe rachidien)\\ \emph{Croisement des fibres}};
\node[boite, below=of dec, align=center] (moelle){4. Moelle épinière\\ (relai interneuronal)};
\node[boite, below=of moelle, align=center] (moto){5. Motoneurone alpha\\ (corne antérieure)};
\node[boite, below=of moto, align=center] (muscle){6. Muscles effecteurs\\ (main droite)};
\draw[fleche] (cortex) -- (pyr);
\draw[fleche] (pyr) -- (dec);
\draw[fleche] (dec) -- (moelle);
\draw[fleche] (moelle) -- (moto);
\draw[fleche] (moto) -- (muscle);
% Boucle cérébelleuse
\node[boite, right=of pyr, fill=popgreenL, draw=popgreen, xshift=1.2cm, align=center] (cer){Cervelet\\ \emph{Coordination, correction,  équilibre, tonus}};
\draw[fleche, popgreen, dashed] (cortex.east) -- ++(0.6,0) |- (cer.north);
\draw[fleche, popgreen, dashed] (cer.south) |- ++(-0.6,0) -- (moelle.east);
% Lésion
\node[boite, right=of dec, fill=poporangeL, draw=poporange, xshift=1.2cm, align=center] (les){Lésion aire motrice\\ \textbf{GAUCHE} (région main)\\ $\Rightarrow$ Paralysie main DROITE};
\draw[fleche, poporange, dashed] (les.west) -- (dec.east);
\end{tikzpicture}
\legende{Schéma fonctionnel du trajet du message nerveux moteur volontaire (exemple : mouvement de la main droite). 1 : Commande dans l'aire motrice primaire gauche ; 2 : Descente par la voie pyramidale ; 3 : Croisement des fibres au bulbe (décussation des pyramides) ; 4 : Relai dans la moelle épinière ; 5 : Motoneurone spinal ; 6 : Contraction des muscles effecteurs de la main droite. En vert (pointillés) : boucle de contrôle cérébelleuse (copie du programme + rétroaction proprioceptive $\rightarrow$ corrections). En orange (pointillés) : siège de la lésion responsable d'une hémiplégie controlatérale (main droite).}
\end{popfigure}
\item La \textbf{boucle cérébelleuse} (en vert) reçoit une copie de l'influx moteur descendant (collatérales des fibres pyramidales) et les informations sensitives de retour (proprioceptives, vestibulaires, visuelles) ; elle élabore des messages correcteurs qu'elle renvoie vers la moelle épinière (via les voies rubro-spinale, vestibulo-spinale) et vers le cortex (via le thalamus). Son rôle : \textbf{coordonner, corriger et harmoniser} le mouvement en cours (précision, synchronisation, adaptation de la force, équilibre, tonus postural).
\item La paralysie de la main \textbf{droite} résulte d'une lésion de l'aire motrice primaire de l'hémisphère \textbf{gauche} (région de l'homonculus dédiée à la main, face latérale de la circonvolution frontale ascendante), en raison de la décussation des pyramides : la lésion est du côté \emph{controlatéral} à la paralysie (point orange sur le schéma).
\end{enumerate}
\end{corrige}

\begin{corrige}[Exercice 8 — Comparaison des trois maladies neurodégénératives]
Tableau de synthèse complété :
\begin{center}
\begin{tabularx}{\linewidth}{|l|X|X|X|}
\hline
\rowcolor{popblueL} Critère & \textbf{Maladie de Parkinson} & \textbf{Chorée de Huntington} & \textbf{Maladie d'Alzheimer} \\
\hline
Structure encéphalique atteinte & Neurones dopaminergiques de la \textbf{substance noire} (locus niger), projetant sur le \textbf{striatum} (voie nigro-striatale) & Neurones GABAergiques du \textbf{striatum} (noyau caudé, putamen) et du \textbf{cortex cérébral} (atrophie cortico-striatale) & Neurones de l'\textbf{hippocampe} et du \textbf{cortex d'association} (plaques séniles $\beta$-amyloïde, dégénérescence neurofibrillaire protéine Tau) \\
\hline
Symptômes moteurs principaux & \textbf{Syndrome parkinsonien} : Tremblement de repos (4--6 Hz), rigidité plastique (roue dentée), \textbf{akinésie} (bradykinésie, hypokinésie), marche à petits pas, instabilité posturale tardive & \textbf{Chorée} : Mouvements involontaires, brusques, irréguliers, amples, migratoires (hyperkinésie), grimaces, troubles de la marche et de la déglutition & Troubles moteurs \textbf{absents ou tardifs} (stade sévère : marche hésitante, rigidité, troubles de la déglutition) \\
\hline
Symptômes non moteurs & Dépression, troubles du sommeil, constipation, troubles autonomes, démence (stade avancé) & Troubles psychiatriques (dépression, irritabilité, psychose), \textbf{déclin intellectuel progressif} (démence subcorticale) & \textbf{Troubles cognitifs majeurs} : Amnésie antérograde, aphasie, apraxie, agnosie, troubles des fonctions exécutives, désorientation \\
\hline
Origine étiologique & Majoritairement \textbf{sporadique} (facteurs environnementaux : pesticides, âge, stress oxydatif) ; formes familiales rares (gènes \emph{SNCA}, \emph{LRRK2}, \emph{PARKIN}...) & \textbf{Héréditaire} : \textbf{Autosomique dominante} (chromosome 4, gène \emph{HTT}, expansion triplet CAG $>$ 35) ; pénétrance complète & Majoritairement \textbf{sporadique} (âge, facteur de risque ApoE4) ; formes familiales précoces rares (gènes \emph{APP}, \emph{PSEN1}, \emph{PSEN2}) \\
\hline
Traitement symptomatique actuel & \textbf{L-DOPA} (précurseur DA franchissant la BHE) + inhibiteurs de la DOPA-décarboxylase périphérique (carbidopa/benserazide) ; agonistes dopaminergiques ; stimulation cérébrale profonde (STN) & Aucun traitement curatif ni modificateur de la maladie ; neuroleptiques (halopéridol, rispéridone) ou tétrabénazine pour atténuer la chorée ; prise en charge psychiatrique & Aucun traitement curatif ; inhibiteurs de l'acétylcholinestérase (donépézil, rivastigmine) et antagoniste NMDA (mémantine) : effet modéré et transitoire sur la cognition \\
\hline
\end{tabularx}
\end{center}
\textbf{Mesures de limitation des lésions / prévention (au moins deux attendues) :}
\begin{itemize}
\item \textbf{Lutter contre les facteurs toxiques environnementaux :} éviter l'exposition professionnelle ou domestique aux pesticides (paraquat, rotenone), solvants organiques, métaux lourds ; limiter la consommation d'alcool et de tabac.
\item \textbf{Prévenir les lésions vasculaires cérébrales :} dépister et traiter précocement l'hypertension artérielle (principal facteur de risque d'AVC et de démence vasculaire), le diabète, l'hypercholestérolémie ; lutter contre la sédentarité.
\item \textbf{Stimuler la réserve cognitive et la plasticité cérébrale :} pratiquer une activité physique régulière (marche, sport), une activité intellectuelle et sociale soutenue (lecture, jeux, apprentissages), maintenir une alimentation équilibrée (type méditerranéen, riche en antioxydants, oméga-3).
\item \textbf{Recourir au conseil génétique :} dans les familles atteintes de chorée de Huntington (ou formes familiales d'Alzheimer/Parkinson), pour informer sur le risque de transmission, proposer le test prédictif (avec accompagnement psychologique) et discuter des options reproductives (DPI).
\end{itemize}
\end{corrige}

\begin{corrige}[Exercice 9 — Type Bac : l'aire motrice du cortex (analyse de documents)]
\begin{enumerate}
\item \textbf{Définition (1 pt).} Un mouvement volontaire est un mouvement décidé consciemment, commandé par l'aire motrice primaire du cortex cérébral (circonvolution frontale ascendante), transmis aux muscles squelettiques par la voie pyramidale (faisceau cortico-spinal) et coordonné par le cervelet.
\item \textbf{Document 1 — Activité neuronale pré-motrice (2 pts).} Les neurones de la circonvolution frontale ascendante déchargent des potentiels d'action \textbf{avant} le début effectif du mouvement de la main (activité pré-motrice). Si cette zone ne faisait que « réagir » au mouvement (réception d'influx sensitifs de retour), son activité surviendrait \emph{pendant ou après} l'exécution motrice. Le fait qu'elle \emph{précède} le mouvement prouve qu'elle en émet la \textbf{commande motrice} : c'est bien le siège de l'initiation et de la programmation du mouvement volontaire.
\item \textbf{Document 2 — Stimulations corticales (3 pts).}
\begin{itemize}
\item \textbf{Organisation somatotopique (1,5 pt) :} Les stimulations 1, 2 et 3, appliquées à des points différents de la circonvolution frontale ascendante, provoquent chacune le mouvement d'une partie précise du corps (jambe pour la stimulation 1 en haut, doigts/main pour la stimulation 2 au milieu, face pour la stimulation 3 en bas). L'aire motrice primaire contient donc une \emph{carte ordonnée du corps} (homonculus moteur), chaque région corticale commandant une région corporelle déterminée.
\item \textbf{Latéralisation / Contrôle controlatéral (1,5 pt) :} Tous les effets moteurs (stim 1, 2, 3) et sensitifs (stim 4) s'observent du côté \textbf{opposé} (controlatéral) à l'hémisphère stimulé. Cela s'explique par la \textbf{décussation des pyramides} au niveau du bulbe rachidien : les fibres de la voie pyramidale croisent la ligne médiane avant de descendre dans la moelle épinière.
\item On note aussi la \textbf{spécialisation fonctionnelle} des aires corticales : la circonvolution frontale ascendante (antérieure au sillon central) est motrice ; la circonvolution pariétale ascendante (postérieure au sillon central, stimulation 4) est sensitive (somesthésique) ; le lobe occipital (stimulation 5) est visuel.
\end{itemize}
\item \textbf{Document 3 — Homonculus moteur (2 pts).} La surface corticale dédiée à une partie du corps n'est pas proportionnelle à sa taille anatomique mais à la \textbf{finesse (précision) de son contrôle moteur} : la main et la face (langue, lèvres), capables de mouvements très fins, complexes et fractionnés, occupent une surface corticale considérable, alors que le tronc et la cuisse, aux mouvements grossiers et globaux, n'occupent qu'une faible surface. C'est l'homonculus moteur « difforme » (grosses mains, grosses lèvres, petit tronc).
\item \textbf{Schéma annoté (2 pts)} :
\begin{popfigure}
\begin{tikzpicture}[scale=1]
% Contour hémisphère gauche (vue latérale)
\draw[thick, popdark] (0,0) .. controls (0.3,1.6) and (1.6,2.6) .. (3.2,2.5)
  .. controls (4.8,2.4) and (6.2,1.6) .. (6.4,0.6)
  .. controls (6.5,-0.1) and (6.0,-0.5) .. (5.2,-0.6)
  .. controls (4.4,-0.7) and (4.2,-1.3) .. (3.4,-1.3)
  .. controls (2.6,-1.3) and (2.2,-0.6) .. (1.4,-0.5)
  .. controls (0.6,-0.4) .. (0,0);
% Sillon central (Rolando)
\draw[very thick, poporange] (3.4,2.45) .. controls (3.1,1.8) and (3.0,1.2) .. (2.9,0.5);
\node[poporange, font=\small, align=center] at (2.1,2.6) {Sillon central\\(de Rolando)};
\draw[poporange, -{Stealth[length=2mm]}] (2.6,2.5) -- (3.15,2.1);
% Aire motrice primaire (circonvolution frontale ascendante)
\draw[very thick, popblue, dashed] (2.9,2.5) .. controls (2.6,1.8) and (2.5,1.2) .. (2.4,0.5);
\node[popblue, font=\small, align=center] at (1.2,1.9) {Aire motrice\\ primaire (4)};
\draw[popblue, -{Stealth[length=2mm]}] (1.9,1.75) -- (2.55,1.5);
% Aire sensitive primaire (circonvolution pariétale ascendante)
\draw[very thick, popgreen, dashed] (3.9,2.4) .. controls (3.6,1.8) and (3.5,1.2) .. (3.4,0.5);
\node[popgreen, font=\small, align=center] at (5.1,2.7) {Aire sensitive\\ primaire (3-1-2)};
\draw[popgreen, -{Stealth[length=2mm]}] (4.9,2.35) -- (3.75,1.7);
% Lobes
\node[font=\small\bfseries, popblue] at (1.1,0.9) {Lobe frontal};
\node[font=\small\bfseries, popgreen] at (4.6,1.2) {Lobe pariétal};
\node[font=\small\bfseries, poppurple] at (5.9,0.1) {Lobe occipital};
\node[font=\small\bfseries, poppink] at (3.0,-0.9) {Lobe temporal};
% Scissure de Sylvius (lobe temporal)
\draw[popline, thick] (2.4,-0.4) .. controls (3.2,-0.2) and (4.4,-0.3) .. (5.3,-0.55);
% Flèches antéro-postérieur
\draw[<->, popmuted] (0.5,2.7) -- node[above, font=\tiny] {Antérieur} (2.5,2.7);
\draw[<->, popmuted] (4.0,2.7) -- node[above, font=\tiny] {Postérieur} (6.0,2.7);
\end{tikzpicture}
\legende{Hémisphère cérébral gauche en vue latérale (face antérieure à gauche). Le sillon central (orange) sépare le lobe frontal du lobe pariétal. L'aire motrice primaire (bleu, aire 4) occupe la circonvolution frontale ascendante, \textbf{en avant} du sillon central. L'aire sensitive primaire (vert, aires 3-1-2) occupe la circonvolution pariétale ascendante, \textbf{en arrière} du sillon central. La scissure latérale (Sylvius) délimite le lobe temporal.}
\end{popfigure}
\end{enumerate}
\textbf{Barème indicatif (sur 10 pts) :} 1 pt (définition complète) ; 2 pts (argument « activité avant le mouvement » explicitement lié à la commande) ; 3 pts (somatotopie 1,5 pt + latéralisation 1,5 pt, avec citation des numéros de stimulation) ; 2 pts (proportionnalité à la finesse du contrôle, pas à la taille) ; 2 pts (schéma : 7 éléments corrects — hémisphère, sillon central, aire motrice antérieure, aire sensitive postérieure, 3 lobes nommés, scissure de Sylvius, flèches antéro-postérieures).

\textbf{Points de vigilance (pièges fréquents) :}
\begin{itemize}
\item Ne pas écrire que le cortex moteur « reçoit » l'ordre — il l'\emph{émet} (commande descendante).
\item Citer explicitement la \textbf{décussation des pyramides} (au bulbe) pour justifier le caractère controlatéral ; ne pas dire « le cerveau gauche commande le côté gauche ».
\item Sur le schéma : placer l'aire motrice \textbf{en avant} du sillon central et l'aire sensitive \textbf{en arrière} (inversion fréquente et éliminatoire). Ne pas confondre lobe frontal et lobe pariétal.
\item Préciser « aire 4 de Brodmann » pour l'aire motrice, « aires 3-1-2 » pour l'aire sensitive.
\end{itemize}
\end{corrige}

\begin{corrige}[Exercice 10 — Type Bac : étude de cas clinique à l'hôpital de Yaoundé (Parkinson)]
\begin{enumerate}
\item \textbf{Diagnostic (2 pts).} Monsieur N. souffre de la \textbf{maladie de Parkinson}. Justification par la triade symptomatique classique (syndrome parkinsonien) : \textbf{tremblement de repos} (disparaissant pendant le mouvement volontaire et le sommeil), \textbf{rigidité} musculaire des membres (hypertonie plastique, « roue dentée ») et \textbf{akinésie} (lenteur d'initiation et d'exécution des mouvements, marche à petits pas, mimique figée). L'âge (62 ans) est compatible avec l'âge typique de début (55--65 ans).
\item \textbf{Structure atteinte et dysfonctionnement (3 pts).} D'après le document 1, le taux de dopamine dans la substance noire du patient est de \SI{0,6}{ng/mg} contre \SI{3,2}{ng/mg} chez le témoin, soit une chute d'environ :
\[ \dfrac{3,2 - 0,6}{3,2} \times 100 = \dfrac{2,6}{3,2} \times 100 \approx 81\ \% \text{ de diminution.} \]
D'après le document 2, la substance noire (locus niger) contient des neurones dopaminergiques projetant sur le striatum (noyau caudé et putamen), structure des ganglions de la base impliquée dans le déclenchement, l'amplitude et la fluidité des mouvements volontaires. La maladie résulte donc de la \textbf{dégénérescence des neurones dopaminergiques de la substance noire}, entraînant un \textbf{déficit majeur en dopamine} au niveau du striatum : la modulation inhibitrice/excitrice des circuits des ganglions de la base est perturbée, d'où l'akinésie, la rigidité et les tremblements.
\item \textbf{Mode d'action de la L-DOPA (3 pts).} La L-DOPA (L-3,4-dihydroxyphénylalanine) est le \textbf{précurseur métabolique direct de la dopamine}. Administrée par voie orale, elle \textbf{franchit la barrière hémato-encéphalique (BHE)} par un transporteur spécifique pour les acides aminés neutres à grande neutralité, pénètre dans l'encéphale et y est transformée en dopamine par l'enzyme \textbf{DOPA-décarboxylase} (présente dans les neurones dopaminergiques survivants et d'autres cellules) : elle \emph{compense le déficit dopaminergique} striatal, ce qui explique la diminution des tremblements et de la rigidité. La dopamine elle-même ne peut pas être administrée directement car elle \textbf{ne franchit pas la barrière hémato-encéphalique} (molécule polaire, pas de transporteur) : injectée par voie générale, elle n'atteindrait jamais le striatum et serait dégradée en périphérie (effets cardiovasculaires indésirables).
\item \textbf{Origine génétique (2 pts).} Non : dans la grande majorité des cas (\textasciitilde 90--95 \%), la maladie de Parkinson est \textbf{sporadique} (non héréditaire), multifactorielle (âge, exposition aux pesticides, stress oxydatif, facteurs génétiques de susceptibilité comme \emph{LRRK2}, \emph{GBA}) ; seules de rares formes familiales monogéniques existent. En revanche, la \textbf{chorée de Huntington} est une maladie neurodégénérative \textbf{héréditaire}, transmise selon le mode \textbf{autosomique dominant} (allèle morbide \emph{HTT} avec expansion CAG sur le chromosome 4 : un seul exemplaire de l'allèle muté suffit à déclencher la maladie).
\end{enumerate}
\textbf{Barème indicatif (sur 10 pts) :} 2 pts (nom de la maladie + citation explicite des 3 symptômes cardinaux) ; 3 pts (exploitation chiffrée du doc 1 : calcul de la baisse \textasciitilde 80 \% + identification substance noire + lien déficit DA / troubles moteurs via striatum) ; 3 pts (L-DOPA = précurseur + franchit BHE + transformation en DA ; justification de l'inefficacité de la DA : ne franchit pas BHE) ; 2 pts (Parkinson = sporadique vs Huntington = autosomique dominant).

\textbf{Points de vigilance :}
\begin{itemize}
\item Exploiter les \emph{valeurs chiffrées} du tableau (comparaison 0,6 vs 3,2 ng/mg) — une réponse qualitative sans chiffres perd des points.
\item Ne pas écrire que la L-DOPA « guérit » la maladie : c'est un traitement \emph{symptomatique} qui ne stoppe pas la dégénérescence neuronale.
\item Bien orthographier « \textbf{barrière hémato-encéphalique} » (pas hémato-encéphalique sans trait d'union, pas hématoencéphalique).
\item Distinguer clairement \emph{sporadique} (pas de transmission familiale évidente) et \emph{héréditaire} (transmission mendélienne).
\end{itemize}
\end{corrige}

\begin{corrige}[Exercice 11 — Type Bac : la chorée de Huntington dans une famille camerounaise]
\begin{enumerate}
\item \textbf{Mode de transmission (3 pts).}
\begin{itemize}
\item \textbf{Allèle dominant (1,5 pt) :} L'individu I-1 est atteinte et son conjoint I-2 est sain. Ils ont une fille saine (II-1) et deux enfants atteints (II-2, II-3). Si l'allèle morbide était \emph{récessif}, l'individu I-1 atteinte serait homozygote pour l'allèle morbide ($h/h$) et ne pourrait transmettre que cet allèle ; croisée avec un homme sain (nécessairement $H/H$ ou $H/h$ selon la fréquence, mais phénotypiquement sain), \emph{tous} les enfants seraient au moins hétérozygotes et donc \emph{porteurs sains} (si $H$ dominant) ou atteints (si $h$ dominant — contradiction). Le fait qu'un couple « atteint $\times$ sain » produise des enfants \emph{sains} (II-1) et des enfants \emph{atteints} (II-2, II-3) prouve qu'un seul exemplaire de l'allèle morbide suffit à s'exprimer : l'allèle est \textbf{dominant}. Notons $H$ l'allèle morbide dominant, $h$ l'allèle normal.
\item \textbf{Porté par un autosome (1,5 pt) :} La maladie touche indifféremment les \textbf{deux sexes} (femmes : I-1, II-2 ; homme : II-3) et se transmet de la \textbf{mère} (I-1) à son \textbf{fils} (II-3) ainsi qu'à sa fille (II-2). Une transmission liée au chromosome X (dominante ou récessive) ne permet pas d'expliquer simplement la transmission d'une mère atteinte à un fils atteint \emph{et} à une fille saine dans les proportions observées sans hypothèses supplémentaires complexes. L'argument décisif au Bac : « La maladie atteint les deux sexes et se transmet par les deux sexes (ici mère $\rightarrow$ fils et mère $\rightarrow$ fille), l'allèle est donc porté par un \textbf{autosome} (chromosome non sexuel, ici le chromosome 4). »
\end{itemize}
\textbf{Conclusion :} Transmission \textbf{autosomique dominante}.
\item \textbf{Génotypes (3 pts).} On note $H$ l'allèle morbide (dominant), $h$ l'allèle normal.
\begin{itemize}
\item \textbf{I-2} est sain : phénotype sain $\Rightarrow$ génotype \textbf{$h/h$} (homozygote normal).
\item \textbf{I-1} est atteinte : elle possède au moins un allèle $H$. Or elle a une fille saine II-1 ($h/h$) qui a reçu un allèle $h$ de chaque parent ; I-1 possède donc aussi un allèle $h$ : elle est \textbf{hétérozygote $H/h$}.
\item \textbf{II-2} est atteinte : elle possède au moins un $H$. Son père I-2 ($h/h$) ne peut lui transmettre qu'un $h$ ; elle a donc obligatoirement reçu $H$ de sa mère I-1 : elle est \textbf{hétérozygote $H/h$}.
\end{itemize}
\[ \text{I-1 : } H/h \qquad \text{I-2 : } h/h \qquad \text{II-2 : } H/h \]
\item \textbf{Risque pour III-1 (2 pts).} La mère II-2 est $H/h$ (atteinte) et le père est sain ($h/h$). Échiquier de croisement :
\begin{center}
\begin{tabular}{|c|c|c|}
\hline
\rowcolor{popblueL} Mère II-2 $\backslash$ Père & $h$ & $h$ \\
\hline
$H$ & $H/h$ (atteint) & $H/h$ (atteint) \\
\hline
$h$ & $h/h$ (sain) & $h/h$ (sain) \\
\hline
\end{tabular}
\end{center}
\[ P(\text{III-1 a reçu l'allèle } H) = \dfrac{2}{4} = \dfrac{1}{2} = 50\ \%. \]
III-1 a donc \textbf{une chance sur deux (50 \%)} de porter l'allèle $H$ et, dans ce cas, de développer la maladie (pénétrance complète) généralement vers 35--50 ans (d'où l'importance du conseil génétique pour ce jeune adulte asymptomatique de 20 ans).
\item \textbf{Comparaison Huntington / Parkinson (2 pts).}
\begin{center}
\begin{tabularx}{\linewidth}{|l|X|X|}
\hline
\rowcolor{popblueL} Critère & \textbf{Chorée de Huntington} & \textbf{Maladie de Parkinson} \\
\hline
Structure atteinte & Neurones GABAergiques du \textbf{striatum} (noyau caudé, putamen) et du cortex & Neurones dopaminergiques de la \textbf{substance noire} (projection nigro-striatale) \\
\hline
Troubles moteurs & \textbf{Hyperkinésie} : Mouvements involontaires, brusques, désordonnés, amples (chorée), grimaces & \textbf{Hypokinésie} : Tremblement de repos, rigidité, \textbf{akinésie} (lenteur, pauvreté), marche à petits pas \\
\hline
Origine génétique & \textbf{Héréditaire} : Autosomique dominante (chr 4, gène \emph{HTT}, expansion CAG) & Majoritairement \textbf{sporadique} (multifactorielle) ; formes familiales rares \\
\hline
\end{tabularx}
\end{center}
\end{enumerate}
\textbf{Barème indicatif (sur 10 pts) :} 3 pts (dominance justifiée par croisement I-1$\times$I-2 + enfant sain 1,5 pt + autosomie justifiée par sexes atteints + transmission mère-fils 1,5 pt) ; 3 pts (trois génotypes corrects et justifiés, 1 pt chacun) ; 2 pts (échiquier correct + probabilité 50 \% + mention pénétrance tardive) ; 2 pts (tableau comparatif : 3 lignes correctes).

\textbf{Points de vigilance :}
\begin{itemize}
\item Pour prouver la \textbf{dominance} : citer explicitement le croisement I-1 $\times$ I-2 et l'existence d'enfants sains (II-1) — une affirmation sans référence à l'arbre n'est pas une démonstration.
\item Pour prouver l'\textbf{autosomie} : ne pas se contenter de dire « les deux sexes sont atteints » (constat) ; il faut expliciter : « transmission par la mère à son fils (I-1 $\rightarrow$ II-3) incompatible avec une transmission liée à Y ; transmission à une fille saine (II-1) et un fils atteint (II-3) compatible avec l'autosomie » ou plus simplement « atteinte des deux sexes + transmission par les deux sexes $\Rightarrow$ autosome ».
\item Rappeler que III-1, asymptomatique à 20 ans, n'est pas forcément indemne : la maladie se déclarant tardivement (35--50 ans), le risque génétique (50 \%) persiste toute la vie.
\item Ne pas confondre \textbf{chorée} (mouvements rapides, dansants, involontaires = hyperkinésie) et \textbf{tremblement parkinsonien} (lent, rythmique, de repos = hypokinésie associée).
\end{itemize}
\end{corrige}

\newpage

\coursec{Fiche bilan}

\begin{popfigure}
\begin{tikzpicture}[
  mindmap,
  grow cyclic,
  every node/.style={concept, font=\small},
  root concept/.append style={concept color=popblue, font=\bfseries\small, minimum size=3.2cm},
  level 1/.append style={level distance=4.2cm, sibling angle=60, font=\footnotesize},
  level 2/.append style={level distance=2.6cm, sibling angle=45, font=\scriptsize},
  concept connection/.append style={color=popmuted}
]
\node[root concept]{Motricit\'e volontaire et enc\'ephale}
  child[concept color=poporange]{ node {Mouvement volontaire}
    child { node {Acte d\'ecid\'e, commande corticale} }
  }
  child[concept color=popgreen]{ node {Aires enc\'ephaliques}
    child { node {Aire motrice : gyrus pr\'ecentral (4)} }
    child { node {Aire sensitive : gyrus postcentral (3,1,2)} }
    child { node {Homunculus : somatotopie} }
  }
  child[concept color=poppurple]{ node {Preuves exp\'erimentales}
    child { node {Stimulations corticales (Penfield)} }
    child { node {L\'esions $\Rightarrow$ d\'eficits localis\'es} }
  }
  child[concept color=popblue]{ node {Trajet de l'influx}
    child { node {Voie pyramidale : cortex $\to$ moelle} }
    child { node {Relais : cervelet, noyaux gris} }
    child { node {Motoneurone $\to$ muscle} }
  }
  child[concept color=popgold]{ node {Parkinson}
    child { node {D\'eficit en dopamine (substance noire)} }
    child { node {Tremblement, akin\'esie, rigidit\'e} }
  }
  child[concept color=poppink]{ node {Huntington et Alzheimer}
    child { node {Huntington : chor\'ee, g\'en\'etique} }
    child { node {Alzheimer : plaques, d\'emence} }
  };
\end{tikzpicture}
\legende{Carte mentale de la le\c{c}on : le mouvement volontaire na\^it dans le cortex moteur, emprunte la voie pyramidale et peut \^etre perturb\'e par des affections localis\'ees (Parkinson, Huntington, Alzheimer).}
\end{popfigure}

\begin{bilan}
\textbf{D\'efinitions cl\'es}
\begin{itemize}
  \item \cle{Mouvement volontaire} : mouvement d\'ecid\'e et command\'e par le \cle{cortex c\'er\'ebral} (aire motrice), ex\'ecut\'e par les muscles squelettiques sous contr\^ole de la conscience.
  \item \cle{Aire motrice} (aire 4 de Brodmann) : zone du \cle{gyrus pr\'ecentral} du lobe frontal dont la stimulation provoque des mouvements localis\'es du c\^ot\'e oppos\'e du corps.
  \item \cle{Aire sensitive} (aires 3, 1, 2) : zone du \cle{gyrus postcentral} du lobe pari\'etal recevant les informations somesth\'esiques (tact, douleur, proprioception).
  \item \cle{Homunculus} : repr\'esentation somatotopique du corps sur le cortex ; la surface corticale est proportionnelle \`a la pr\'ecision du contr\^ole (main, langue) et non \`a la taille de l'organe.
  \item \cle{Voie pyramidale} : faisceau de neurones moteurs reliant le cortex moteur \`a la moelle \'epini\`ere, avec \cle{d\'ecussation} (croisement) au niveau du bulbe rachidien.
\end{itemize}

\textbf{Trajet de l'influx nerveux moteur (sch\'ema fonctionnel \`a savoir reproduire)}
\[ \text{Cortex moteur} \Rightarrow \text{voie pyramidale (d\'ecussation)} \Rightarrow \text{motoneurone spinal} \Rightarrow \text{muscle squelettique} \]
avec contr\^ole et coordination par le \cle{cervelet} (\'equilibre, pr\'ecision) et les \cle{noyaux gris centraux} (programmation du mouvement).

\textbf{Maladies du contr\^ole de la motricit\'e}
\begin{center}
\begin{tabularx}{\linewidth}{|l|X|X|X|}
\hline
\rowcolor{popblueL} \textbf{Maladie} & \textbf{Structure atteinte} & \textbf{Sympt\^omes} & \textbf{Cause} \\
\hline
Parkinson & Substance noire (noyaux gris) : d\'eficit en \cle{dopamine} & Tremblement de repos, rigidit\'e, akin\'esie & D\'eg\'en\'erescence neuronale ; traitement : L-DOPA \\
\hline
Huntington & Corps stri\'e (noyaux gris) & Chor\'ee (mouvements involontaires), d\'emence & Maladie h\'er\'editaire autosomique dominante \\
\hline
Alzheimer & Cortex c\'er\'ebral (diffus) & Perte de m\'emoire, d\'emence, apraxie & Plaques amylo\"ides, d\'eg\'en\'erescence neuronale \\
\hline
\end{tabularx}
\end{center}

\textbf{M\'ethodes express}
\begin{itemize}
  \item \emph{Interpr\'eter une stimulation corticale} : stimulation d'un point de l'aire motrice $\Rightarrow$ contraction d'un muscle pr\'ecis du c\^ot\'e \textbf{oppos\'e} $\Rightarrow$ conclusion : localisation des fonctions motrices et innervation crois\'ee.
  \item \emph{Interpr\'eter une l\'esion} : l\'esion localis\'ee $\Rightarrow$ d\'eficit pr\'ecis (paralysie, tremblement\ldots) $\Rightarrow$ la structure l\'es\'ee est impliqu\'ee dans la fonction abolie.
  \item \emph{Dissection de l'enc\'ephale} : rep\'erer cerveau (h\'emisph\`eres, sillons), cervelet, tronc c\'er\'ebral ; sch\'ema annot\'e obligatoire avec l\'egendes pr\'ecises.
\end{itemize}
\end{bilan}

\begin{aretenir}[Checklist avant le Bac]
\begin{itemize}
  \item $\square$ Je sais d\'efinir un mouvement volontaire et le distinguer d'un r\'eflexe.
  \item $\square$ Je situe l'aire motrice (gyrus pr\'ecentral, lobe frontal) et l'aire sensitive (gyrus postcentral, lobe pari\'etal).
  \item $\square$ Je sais que l'innervation motrice est crois\'ee (d\'ecussation pyramidale).
  \item $\square$ Je sais interpr\'eter une exp\'erience de stimulation corticale (Penfield).
  \item $\square$ Je sais interpr\'eter les cons\'equences d'une l\'esion corticale localis\'ee.
  \item $\square$ Je sais expliquer la notion d'homunculus (somatotopie).
  \item $\square$ Je sais reproduire le sch\'ema fonctionnel du trajet de l'influx moteur.
  \item $\square$ Je connais le r\^ole du cervelet et des noyaux gris centraux dans la coordination.
  \item $\square$ Je connais les sympt\^omes et la cause biochimique de la maladie de Parkinson (dopamine).
  \item $\square$ Je distingue Huntington (h\'er\'editaire, chor\'ee) et Alzheimer (d\'emence, plaques amylo\"ides).
  \item $\square$ Je sais r\'ealiser un sch\'ema annot\'e d'un enc\'ephale de Mammif\`ere.
  \item $\square$ Je r\'edige mes interpr\'etations en suivant la d\'emarche : observation $\to$ hypoth\`ese $\to$ interpr\'etation $\to$ conclusion.
\end{itemize}
\end{aretenir}

\findecours

\end{document}
